% Mechanically extracted after \mainmatter from the frozen standalone master.

\markboth{Exposé XXI. Root Data}{Exposé XXI. Root Data}
% BEGIN inputs/XXI/tex/components/00_expose_XXI_opening_through_scholium1113_en.tex
\chapter*{Exposé XXI. Root Data}
\addcontentsline{toc}{chapter}{Exposé XXI. Root Data}

\noindent
This Exposé\footnote{Editors' note: Version of 13 October 2024.}
collects, in the absence of a suitable reference,\footnote{%
Editors' note: For results on root systems (Sections 1--5), one may
consult \emph{Groupes et algèbres de Lie}, Chapter VI.}
known results on root data (that is, ``abstract'' root systems), most of
which will be used below.

\medskip
\noindent\textbf{Notation.}
We denote by \(\Qp\) the set of positive (or zero) rational numbers;
\(\mathbb Z\cap\Qp=\mathbb N\). Let \(V\) be a \(\mathbb Q\)-vector
space. If \(A\) (respectively \(B\)) is a subset of \(\mathbb Q\)
(respectively of \(V\)), we denote by \(A\cdot B\) the image of
\(A\otimes B\) under the morphism
\(\mathbb Q\otimes V\to V\), in other words the set of linear
combinations of elements of \(B\) with coefficients in \(A\). We write
\(-B=\{-1\}B\). We denote by \(E-F\) the set of elements of \(E\)
that do not belong to \(F\).

\section{Generalities}

\SGAInlineRefTarget{sga3:xxi:subsection:1-1}
\subsection{Definitions and first properties}

\begin{definition}\label{sga3:XXI:1.1.1}
Let \(M\) and \(M^\ast\) be two free finite-type \(\mathbb Z\)-modules
in duality. Write
\[
  V=M\otimes\mathbb Q,\qquad V^\ast=M^\ast\otimes\mathbb Q;
\]
these are two \(\mathbb Q\)-vector spaces in duality. We identify \(M\)
(respectively \(M^\ast\)) with a subset of \(V\) (respectively
\(V^\ast\)). The canonical bilinear form on
\(M^\ast\times M\) (respectively \(V^\ast\times V\)) is denoted
\((\, ,\,)\).

Let \(R\) be a finite subset of \(M\). Suppose given a map
\(\alpha\mapsto\alpha^\ast\) from \(R\) to \(M^\ast\); the set of
\(\alpha^\ast\), for \(\alpha\in R\), is denoted \(R^\ast\).
To every \(\alpha\in R\) associate the endomorphism \(s_\alpha\)
(respectively \(s_\alpha^\ast\)) of \(M\) and \(V\) (respectively of
\(M^\ast\) and \(V^\ast\)) given by
\begin{align}
 s_\alpha(x)&=x-(\alpha^\ast,x)\alpha,
 &s_\alpha&=\operatorname{id}-\alpha^\ast\otimes\alpha,
 \tag{1}\label{sga3:XXI:eq:1}\\
 s_\alpha^\ast(u)&=u-(u,\alpha)\alpha^\ast,
 &s_\alpha^\ast&=\operatorname{id}-\alpha\otimes\alpha^\ast.
 \tag{1${}^\ast$}\label{sga3:XXI:eq:1star}
\end{align}
The pair \((R,R^\ast)\) (more precisely, the pair
\((R,R\to M^\ast)\)) is called a \emph{root datum} in
\((M,M^\ast)\), and \((M,M^\ast,R,R^\ast)\) is called a root datum,
if the following axioms hold:
\begin{align*}
 (\mathrm{DR\,I})\quad&
    (\alpha^\ast,\alpha)=2
    &&\text{for every }\alpha\in R,\\
 (\mathrm{DR\,II})\quad&
    s_\alpha(R)\subset R,\qquad
    s_\alpha^\ast(R^\ast)\subset R^\ast
    &&\text{for every }\alpha\in R.
\end{align*}
The set \(R\) is called the \emph{root system} of the root datum
\(\cR=(M,M^\ast,R,R^\ast)\). The elements of \(R\) (respectively
\(R^\ast\)) are called the \emph{roots} (respectively
\emph{coroots}) of the root datum.
\end{definition}

\begin{remark}\label{sga3:XXI:1.1.2}
Axiom \((\mathrm{DR\,I})\) is equivalent to any one of the following
properties:
\begin{align}
 s_\alpha s_\alpha&=\operatorname{id},
 &s_\alpha^\ast s_\alpha^\ast&=\operatorname{id},
 \tag{2,2${}^\ast$}\label{sga3:XXI:eq:2}\\
 s_\alpha(\alpha)&=-\alpha,
 &s_\alpha^\ast(\alpha^\ast)&=-\alpha^\ast.
 \tag{3,3${}^\ast$}\label{sga3:XXI:eq:3}
\end{align}
\end{remark}

\begin{remark}\label{sga3:XXI:1.1.3}
Axioms \((\mathrm{DR\,I})\) and \((\mathrm{DR\,II})\) imply
\[
 R=-R,\qquad R^\ast=-R^\ast,\qquad
 0\notin R,\qquad 0\notin R^\ast.
\]
\end{remark}

\begin{lemma}\label{sga3:XXI:1.1.4}
The map \(R\to R^\ast\) is a bijection. More generally, if
\(\alpha,\beta\in R\) and
\((\alpha^\ast,x)=(\beta^\ast,x)\) for every \(x\in R\), then
\(\alpha=\beta\).
\end{lemma}

\begin{proof}
Indeed, one then has
\[
 s_\beta(\alpha)=\alpha-2\beta,\qquad
 s_\alpha(\beta)=\beta-2\alpha.
\]
It follows at once that
\[
 s_\beta s_\alpha(\alpha)
   =2\beta-\alpha=\alpha+2(\beta-\alpha),\qquad
 s_\beta s_\alpha(\beta-\alpha)
   =s_\beta(\beta-\alpha)=\beta-\alpha.
\]
Hence, by \((\mathrm{DR\,II})\),
\[
 (s_\beta s_\alpha)^n(\alpha)
   =\alpha+2n(\beta-\alpha)\in R.
\]
Since \(R\) is finite, \(\beta-\alpha=0\).
\end{proof}

\begin{corollary}\label{sga3:XXI:1.1.5}
The inverse map \(R^\ast\to R\) defines a root datum
\[
 \cR^\ast=(M^\ast,M,R^\ast,R),
\]
called the \emph{dual} of \(\cR\).\footnote{Editors' note:
\((\alpha^\ast)^\ast=\alpha\).}
\end{corollary}

\begin{definition}\label{sga3:XXI:1.1.6}
We denote by \(\Gamma_0(R)\) the subgroup of \(M\) generated by \(R\).
We denote by \(\VR(R)\) the vector subspace of \(V\) generated by \(R\),
that is,
\[
 \VR(R)=\Gamma_0(R)\otimes\mathbb Q.
\]
Applying these definitions to \(\cR^\ast\), one likewise obtains
\(\Gamma_0(R^\ast)\) and \(\VR(R^\ast)\).

The \emph{reductive rank} of \(\cR\) is the number
\[
 \rgred(\cR)
 =\rank(M)=\dim(V)=\dim(V^\ast)=\rank(M^\ast)=\rgred(\cR^\ast).
\]
The \emph{semisimple rank} of \(\cR\) is the number
\[
 \rgss(\cR)
 =\rank(R)=\rank(\Gamma_0(R))=\dim(\VR(R)).
\]
Thus \(\rgss(\cR)\leq\rgred(\cR)\).

We shall see below that
\(\rgss(\cR)=\rgss(\cR^\ast)\), that is, that \(\VR(R)\) and
\(\VR(R^\ast)\) have the same dimension.
\end{definition}

\begin{definition}\label{sga3:XXI:1.1.7}\leavevmode\par
The root datum \(\cR\) is called \emph{semisimple} (respectively
\emph{trivial}) if
\(\rgss(\cR)=\rgred(\cR)\) (respectively \(\rgss(\cR)=0\)).
Thus \(\cR\) is trivial if and only if \(R\) is empty. The trivial root
datum of reductive rank zero is denoted
\[
 0=(\{0\},\{0\},\varnothing,\varnothing).
\]
\end{definition}

\begin{definition}\label{sga3:XXI:1.1.8}
We denote by \(W(\cR)\) the group of transformations of \(M\) generated
by the \(s_\alpha\), \(\alpha\in R\). It is called the \emph{Weyl
group} of \(\cR\). We write
\[
 W^\ast(\cR)=W(\cR^\ast).
\]
Then \(W(\cR)\) acts on \(R\), \(\Gamma_0(R)\), \(\VR(R)\), \(M\),
and \(V\). If \(w\in W(\cR)\) and \(x\in M\) (respectively
\(x\in V\)), then
\[
 wx-x\in\Gamma_0(R)
 \quad\text{(respectively }wx-x\in\VR(R)\text{)};
\]
this follows immediately from formula \eqref{sga3:XXI:eq:1}.
Likewise for \(W^\ast(\cR)\).
\end{definition}

\begin{lemma}\label{sga3:XXI:1.1.9}
For every \(\alpha\in R\), \(x\in V\), and \(u\in V^\ast\), one has
\begin{equation}
 (s_\alpha^\ast(u),s_\alpha(x))=(u,x).
 \tag{4}\label{sga3:XXI:eq:4}
\end{equation}
\end{lemma}

\begin{proof}
Taking the scalar product of \eqref{sga3:XXI:eq:1} and
\eqref{sga3:XXI:eq:1star}, the left-hand side is
\[
 (u,x)+(u,\alpha)(\alpha^\ast,x)
       \bigl((\alpha^\ast,\alpha)-2\bigr)=(u,x).
\]
\end{proof}

\begin{remark}\label{sga3:XXI:1.1.10}
If one assumes \(0\notin R\) and \(0\notin R^\ast\), then
Lemma~\ref{sga3:XXI:1.1.9} is equivalent to
\((\mathrm{DR\,I})\).
\end{remark}

\begin{corollary}\label{sga3:XXI:1.1.11}
Let \(h\mapsto h^\vee\) denote the isomorphism from
\(\GL(M)\) onto \(\GL(M^\ast)\) that associates to \(h\) its
contragredient. Formula \eqref{sga3:XXI:eq:4} may then also be written
\begin{equation}
 (s_\alpha)^\vee=s_\alpha^\ast.
 \tag{5}\label{sga3:XXI:eq:5}
\end{equation}
\end{corollary}

\begin{corollary}\label{sga3:XXI:1.1.12}
The preceding isomorphism induces an isomorphism from \(W(\cR)\) onto
\(W^\ast(\cR)\).
\end{corollary}

\begin{scholium}\label{sga3:XXI:1.1.13}
By the preceding result, we shall identify \(W\) and \(W^\ast\), and
regard \(W\) as a group of transformations of
\[
 R,\ R^\ast,\ M,\ M^\ast,\ \Gamma_0(R),\ \Gamma_0(R^\ast),\
 V,\ V^\ast,\ \VR(R),\ \VR(R^\ast).
\]
We shall write \(s_\alpha\) for \(s_\alpha^\ast\).
\end{scholium}

% END inputs/XXI/tex/components/00_expose_XXI_opening_through_scholium1113_en.tex

% BEGIN inputs/XXI/tex/components/01_expose_XXI_section12_through_application212_en.tex
\subsection{The map \texorpdfstring{\(p\)}{p}}

\begin{lemma}\label{sga3:XXI:1.2.1}
Let \(p:M\to M^\ast\) (respectively \(V\to V^\ast\)) be the linear map
defined by
\begin{equation}
 p(x)=\sum_{u\in R^\ast}(u,x)u.
 \tag{6}\label{sga3:XXI:eq:6}
\end{equation}
Write \(\ell(x)=(p(x),x)\). Then:
\begin{align}
 \ell(x)&\geq0,\qquad \ell(\alpha)>0
     &&\text{for }\alpha\in R,
 \tag{7}\label{sga3:XXI:eq:7}\\
 \ell(wx)&=\ell(x)
     &&\text{for }w\in W,
 \tag{8}\label{sga3:XXI:eq:8}\\
 (p(x),y)&=(p(y),x)
     &&\text{for }x,y\in V,
 \tag{9}\label{sga3:XXI:eq:9}\\
 \ell(\alpha)\alpha^\ast&=2p(\alpha)
     &&\text{for }\alpha\in R.
 \tag{10}\label{sga3:XXI:eq:10}
\end{align}
\end{lemma}

\begin{proof}
The first three relations are immediate.\footnote{Editors' note:
Formula \eqref{sga3:XXI:eq:8} follows from
\eqref{sga3:XXI:eq:4}, \eqref{sga3:XXI:eq:2}, and
\((\mathrm{DR\,II})\).}
We prove the last one. By \eqref{sga3:XXI:eq:1}, for \(u\in V^\ast\),
\begin{align*}
 (u,\alpha)^2\alpha^\ast
  &=(u,\alpha)u-(u,\alpha)s_\alpha(u)\\
  &=(u,\alpha)u+(u,s_\alpha(\alpha))s_\alpha(u)\\
  &=(u,\alpha)u+(s_\alpha(u),\alpha)s_\alpha(u)
       \qquad\bigl(s_\alpha^2=\operatorname{id}\bigr).
\end{align*}
Since \(s_\alpha\) permutes \(R^\ast\) by
\((\mathrm{DR\,II})\), summing over \(u\in R^\ast\) gives the result.
\end{proof}

\begin{scholium}\label{sga3:XXI:1.2.2}
Relation \eqref{sga3:XXI:eq:10} says that the map
\(\alpha\mapsto\ell(\alpha)\alpha^\ast\) is the restriction to \(R\)
of a linear map from \(M\) to \(M^\ast\). In particular,
\[
 (-\alpha)^\ast=-\alpha^\ast.
\]
\end{scholium}

\begin{corollary}\label{sga3:XXI:1.2.3}
The map \(p\) induces an isomorphism from \(\VR(R)\) onto
\(\VR(R^\ast)\).
\end{corollary}

\begin{proof}
The map \(p\) sends \(\VR(R)\) onto \(\VR(R^\ast)\). Hence
\[
 \dim\VR(R)\geq\dim\VR(R^\ast).
\]
Applying this inequality to the dual root datum gives
\[
 \dim\VR(R)=\dim\VR(R^\ast),
\]
so the surjective map \(p\) is also bijective.
\end{proof}

\begin{corollary}\label{sga3:XXI:1.2.4}
One has \(\dim\VR(R)=\dim\VR(R^\ast)\), and therefore
\[
 \rgss(\cR)=\rgss(\cR^\ast).
\]
\end{corollary}

\begin{corollary}\label{sga3:XXI:1.2.5}
The bilinear form \((u,x)\) is nondegenerate on
\(\VR(R^\ast)\times\VR(R)\), and thus puts these
\(\mathbb Q\)-vector spaces in duality.
\end{corollary}

\begin{proof}
If \((u,x)=0\) for every \(u\in R^\ast\), then \(p(x)=0\).
\end{proof}

\begin{corollary}\label{sga3:XXI:1.2.6}
The symmetric bilinear form \((p(x),y)\) is positive and nondegenerate
on \(\VR(R)\).
\end{corollary}

\begin{corollary}\label{sga3:XXI:1.2.7}
The group \(W\) acts faithfully on \(R\) (and hence on the other sets
listed in Scholium~\ref{sga3:XXI:1.1.13}).
\end{corollary}

\begin{proof}
Let \(u\in V^\ast\), and suppose \(w\in W\) satisfies
\(w(\alpha)=\alpha\) for every \(\alpha\in R\). We prove \(w(u)=u\).
For every \(\alpha\in R\),
\[
 (w(u)-u,\alpha)
 =(w(u),\alpha)-(u,\alpha)
 =(u,w^{-1}(\alpha))-(u,\alpha)=0.
\]
But \(w(u)-u\in\VR(R^\ast)\). Since it is orthogonal to every root, it
vanishes by Corollary~\ref{sga3:XXI:1.2.5}.
\end{proof}

\begin{corollary}\label{sga3:XXI:1.2.8}
The group \(W\) is finite.
\end{corollary}

\begin{proposition}\label{sga3:XXI:1.2.9}
The operations of \(W\) respect the correspondence between roots and
coroots. In other words, for \(\alpha\in R\) and \(w\in W\),
\[
 w(\alpha^\ast)=w(\alpha)^\ast.
\]
\end{proposition}

\begin{proof}
It suffices to verify this for \(w=s_\beta\), \(\beta\in R\), that is,
to verify
\[
 s_\beta(\alpha^\ast)=s_\beta(\alpha)^\ast.
\]
Now \(\ell(s_\beta(\alpha))s_\beta(\alpha)^\ast/2\) equals
\begin{align*}
 p(s_\beta(\alpha))
 &=\sum_{u\in R^\ast}(u,s_\beta(\alpha))u\\
 &=\sum_{u\in R^\ast}(s_\beta(u),\alpha)u\\
 &=\sum_{u\in R^\ast}(u,\alpha)s_\beta(u)
  =s_\beta(p(\alpha)).
\end{align*}
Since \(\ell(s_\beta(\alpha))=\ell(\alpha)\), we obtain
\[
 s_\beta(\alpha)^\ast
 =s_\beta\!\left(\frac{2p(\alpha)}{\ell(\alpha)}\right)
 =s_\beta(\alpha^\ast).
\]
\end{proof}

\begin{corollary}\label{sga3:XXI:1.2.10}
If \(\alpha\in R\) and \(w\in W\), then
\[
 ws_\alpha w^{-1}=s_{w(\alpha)}.
\]
\end{corollary}

\begin{proof}
For \(x\in V\),
\[
 ws_\alpha w^{-1}(x)
 =x-(\alpha^\ast,w^{-1}(x))w(\alpha)
 =x-(w(\alpha^\ast),x)w(\alpha).
\]
By Proposition~\ref{sga3:XXI:1.2.9}, the last expression is
\[
 \begin{aligned}
 x-(w(\alpha)^\ast,x)w(\alpha)
   &=s_{w(\alpha)}(x).
 \end{aligned}
\]
\end{proof}

\begin{corollary}\label{sga3:XXI:1.2.11}\leavevmode\par
Let \(R'\subset R\). The quadruple
\((M,M^\ast,R',R'^\ast)\) is a root datum if and only if
\(\alpha,\beta\in R'\) implies \(s_\alpha(\beta)\in R'\).
\end{corollary}

\section{Relations between two roots}

\subsection{Proportional roots}

\begin{proposition}\label{sga3:XXI:2.1.1}
Let \(\alpha\) and \(\beta\) be two roots. The following conditions are
equivalent:
\begin{enumerate}
\item[(i)] there exists \(k\in\mathbb Q\) such that
  \(\alpha=k\beta\);
\item[(ii)] \(s_\alpha=s_\beta\).
\end{enumerate}
Under these conditions,
\(\alpha^\ast=k^{-1}\beta^\ast\), and \(k\) is one of
\[
 1,\ -1,\ 2,\ -2,\ \tfrac12,\ -\tfrac12.
\]
\end{proposition}

\begin{proof}
Assume first \SGARefLink{sga3:XXI:2.1.1:item:i:02}{(i)}. Then
\[
 \alpha^\ast
 =\ell(\alpha)^{-1}2p(\alpha)
 =k^{-2}\ell(\beta)^{-1}2k\,p(\beta)
 =k^{-1}\beta^\ast.
\]
This immediately gives \(s_\alpha=s_\beta\). Conversely, if
\(s_\alpha=s_\beta\), then
\[
 \alpha=s_\alpha(-\alpha)
       =s_\beta(-\alpha)
       =(\beta^\ast,\alpha)\beta-\alpha,
\]
so
\[
 2\alpha=(\beta^\ast,\alpha)\beta,
\qquad(\beta^\ast,\alpha)\in\mathbb Z;
\]
thus \SGARefLink{sga3:XXI:2.1.1:item:ii:02}{(ii)} implies \SGARefLink{sga3:XXI:2.1.1:item:i:02}{(i)}. Finally, if \(\alpha=k\beta\), then
\(\alpha^\ast=k^{-1}\beta^\ast\), whence
\[
 (\alpha^\ast,\beta)=2k^{-1},
 \qquad(\beta^\ast,\alpha)=2k.
\]
Thus \(2k\) and \(2k^{-1}\) are integers, which gives the stated list.
\end{proof}

\begin{application}\label{sga3:XXI:2.1.2}
Root data \(\cR\) with \(\rgss(\cR)=1\) are of one of the following two
types:
\begin{enumerate}
\item[\SGAInlineRefTarget{sga3:XXI:2.1.1:item:i:02}(i)] Type \(A_1\): there exists \(\alpha\in M\) such that the
  roots are \(\alpha\) and \(-\alpha\). The coroots are then
  \(\alpha^\ast\) and \(-\alpha^\ast\).
\item[\SGAInlineRefTarget{sga3:XXI:2.1.1:item:ii:02}(ii)] Type \(A'_1\): there exists \(\alpha\in M\) such that the
  roots are
  \(\alpha,-\alpha,2\alpha,-2\alpha\). The coroots are then
  \(\alpha^\ast,-\alpha^\ast,\alpha^\ast/2,-\alpha^\ast/2\).
\end{enumerate}
\end{application}

% END inputs/XXI/tex/components/01_expose_XXI_section12_through_application212_en.tex

% BEGIN inputs/XXI/tex/components/02_expose_XXI_definition213_through_definition223_en.tex
\begin{definition}\label{sga3:XXI:2.1.3}
A root \(\alpha\in R\) is called \emph{indivisible} if
\(\alpha/2\notin R\). The root datum \(\cR\) is called
\emph{reduced} if every root is indivisible.
\end{definition}

The root datum \(\cR\) is reduced if and only if \(\cR^\ast\) is
reduced. If \(\alpha\) is indivisible and \(w\in W\), then
\(w(\alpha)\) is indivisible.

\begin{definition}\label{sga3:XXI:2.1.4}
Let \(\alpha\in R\). If \(\alpha\) is indivisible, set
\(\operatorname{ind}(\alpha)=\alpha\); otherwise set
\(\operatorname{ind}(\alpha)=\alpha/2\).
\end{definition}

\begin{corollary}\label{sga3:XXI:2.1.5}
If \(\alpha\in R\) is indivisible and \(k\alpha\in R\), where
\(k\in\mathbb Q\), then \(k\in\mathbb Z\).
\end{corollary}

\begin{proposition}\label{sga3:XXI:2.1.6}
Let \(\cR\) be a root datum. Then
\[
 \operatorname{ind}(\cR)
 =\bigl(M,M^\ast,\operatorname{ind}(R),
         \operatorname{ind}(R)^\ast\bigr)
\]
is a reduced root datum, and
\[
 W\bigl(\operatorname{ind}(\cR)\bigr)=W(\cR).
\]
\end{proposition}

\begin{proof}
The quadruple \(\operatorname{ind}(\cR)\) is a root datum by
Corollary~\ref{sga3:XXI:1.2.11}, since the Weyl group permutes the
indivisible roots. The second assertion follows from
Proposition~\ref{sga3:XXI:2.1.1}.
\end{proof}

\begin{remark}\label{sga3:XXI:2.1.7}
If \(\cR\) is not reduced, then
\[
 \operatorname{ind}(R)^\ast\neq\operatorname{ind}(R^\ast),
 \qquad
 \operatorname{ind}(\cR^\ast)\neq\operatorname{ind}(\cR)^\ast.
\]
\end{remark}

\subsection{Orthogonal roots}

\begin{lemma}\label{sga3:XXI:2.2.1}
Let \(\alpha\) and \(\beta\) be two roots. Then
\begin{equation}
 \ell(\alpha)(\alpha^\ast,\beta)
 =\ell(\beta)(\beta^\ast,\alpha).
 \tag{11}\label{sga3:XXI:eq:11}
\end{equation}
\end{lemma}

\begin{proof}
This follows immediately from Lemma~\ref{sga3:XXI:1.2.1},
formulas \eqref{sga3:XXI:eq:9} and \eqref{sga3:XXI:eq:10}.
\end{proof}

\begin{corollary}\label{sga3:XXI:2.2.2}
Let \(\alpha,\beta\in R\). The following conditions are equivalent:
\begin{enumerate}
\item[\SGAInlineRefTarget{sga3:XXI:2.2.2:item:i}(i)] \((\alpha^\ast,\beta)=0\);
\item[(i bis)] \((\beta^\ast,\alpha)=0\);
\item[(ii)] \((p(\alpha),\beta)=0\);
\item[(iii)] \(s_\alpha(\beta)=\beta\);
\item[(iii bis)] \(s_\beta(\alpha)=\alpha\);
\item[(iv)] \(s_\alpha(\beta^\ast)=\beta^\ast\);
\item[(iv bis)] \(s_\beta(\alpha^\ast)=\alpha^\ast\);
\item[\SGAInlineRefTarget{sga3:XXI:2.2.2:item:v}(v)] \(s_\alpha\neq s_\beta\), and \(s_\alpha\) and \(s_\beta\)
  commute.
\end{enumerate}
\end{corollary}

\begin{proof}
All the equivalences are immediate except those involving \SGARefLink{sga3:XXI:2.2.2:item:v}{(v)}. We
first show that \SGARefLink{sga3:XXI:2.2.2:item:i}{(i)}, together with its equivalent condition (i bis),
implies \SGARefLink{sga3:XXI:2.2.2:item:v}{(v)}. One has
\[
 s_\alpha s_\beta(x)
 =x-(\beta^\ast,x)\beta-(\alpha^\ast,x)\alpha
   +(\beta^\ast,x)(\alpha^\ast,\beta)\alpha.
\]
If \((\alpha^\ast,\beta)=0\), then
\((\beta^\ast,\alpha)=0\), and consequently
\[
 s_\alpha s_\beta(x)=s_\beta s_\alpha(x).
\]

Conversely, suppose \SGARefLink{sga3:XXI:2.2.2:item:v}{(v)}. Then
\[
 s_\alpha=s_\beta s_\alpha s_\beta
          =s_{s_\beta(\alpha)}
 \qquad\text{by Corollary~\ref{sga3:XXI:1.2.10}.}
\]
By Proposition~\ref{sga3:XXI:2.1.1}, there is a
\(k\in\mathbb Q^\ast\) such that
\[
 \alpha=k\,s_\beta(\alpha)
       =k\alpha-k(\beta^\ast,\alpha)\beta.
\]
Since \(s_\alpha\neq s_\beta\), the roots \(\alpha\) and \(\beta\)
are not proportional, again by Proposition~\ref{sga3:XXI:2.1.1}.
It follows that \((\beta^\ast,\alpha)=0\).
\end{proof}

\begin{definition}\label{sga3:XXI:2.2.3}
Two roots satisfying the equivalent conditions of
Corollary~\ref{sga3:XXI:2.2.2} are called \emph{orthogonal}.
\end{definition}

% END inputs/XXI/tex/components/02_expose_XXI_definition213_through_definition223_en.tex

% BEGIN inputs/XXI/tex/components/03_expose_XXI_remark224_through_lemma311_en.tex
\begin{remark}\label{sga3:XXI:2.2.4}
The roots \(\alpha\) and \(\beta\) are orthogonal if and only if the
coroots \(\alpha^\ast\) and \(\beta^\ast\) are orthogonal.
\end{remark}

\begin{lemma}\label{sga3:XXI:2.2.5}
If \(\alpha\) and \(\beta\) are two orthogonal roots, then
\(\alpha+\beta\in R\) if and only if \(\alpha-\beta\in R\).
\end{lemma}

\begin{proof}
Indeed,
\[
 s_\beta(\alpha+\beta)
 =s_\beta(\alpha)+s_\beta(\beta)=\alpha-\beta.
\]
\end{proof}

\begin{lemma}\label{sga3:XXI:2.2.6}
Let \(\alpha\) and \(\beta\) be two nonorthogonal roots. For a coroot
\(\gamma^\ast\), define \(\ell(\gamma^\ast)\) to be the value of
\(\ell\) on the root \(\gamma^\ast\) of \(\cR^\ast\). Then
\begin{equation}
 \ell(\alpha)\ell(\alpha^\ast)
 =\ell(\beta)\ell(\beta^\ast).
 \tag{12}\label{sga3:XXI:eq:12}
\end{equation}
\end{lemma}

\begin{proof}
Multiply formula \eqref{sga3:XXI:eq:11} by the corresponding formula
for \(\cR^\ast\), and use
\((\gamma^\ast)^\ast=\gamma\) for every \(\gamma\in R\). This gives
\[
 (\beta^\ast,\alpha)(\alpha^\ast,\beta)
 \ell(\alpha)\ell(\alpha^\ast)
 =
 (\beta^\ast,\alpha)(\alpha^\ast,\beta)
 \ell(\beta)\ell(\beta^\ast).
\]
The common integer factor is nonzero because the roots are not
orthogonal, so it may be cancelled.
\end{proof}

\subsection{General case}

\begin{proposition}\label{sga3:XXI:2.3.1}
If \(\alpha\) and \(\beta\) are any two roots, then
\begin{equation}
 0\leq(\alpha^\ast,\beta)(\beta^\ast,\alpha)\leq4.
 \tag{13}\label{sga3:XXI:eq:13}
\end{equation}
If \(\alpha\) and \(\beta\) are neither proportional nor orthogonal,
then
\[
 1\leq(\alpha^\ast,\beta)(\beta^\ast,\alpha)\leq3.
\]
\end{proposition}

\begin{proof}
One has
\[
 4(p(\alpha),\beta)^2
 =\ell(\alpha)\ell(\beta)
   (\alpha^\ast,\beta)(\beta^\ast,\alpha).
\]
On the other hand, by Corollary~\ref{sga3:XXI:1.2.6}, the symmetric
bilinear form \((p(x),y)\) is positive and nondegenerate on
\(\VR(R)\). Hence
\[
 (p(x),y)^2\leq\ell(x)\ell(y).
\]
\footnote{Editors' note: equality holds if and only if \(x\) and
\(y\) are proportional.}
The first inequality follows. In the nonproportional, nonorthogonal
case, the product is a nonzero integer and cannot equal \(4\), which
gives the second inequality.
\end{proof}

\begin{corollary}\label{sga3:XXI:2.3.2}
Let \(\alpha\) and \(\beta\) be two nonorthogonal roots. If
\(\ell(\alpha)=\ell(\beta)\), then there is a \(w\in W\) such that
\(w(\alpha)=\beta\).
\end{corollary}

\begin{proof}
If \(\alpha\) and \(\beta\) are proportional, equality of their
\(\ell\)-values gives \(\alpha=\beta\) or \(\alpha=-\beta\); take
\(w=1\) or \(w=s_\alpha\), respectively. If they are neither
proportional nor orthogonal, formula \eqref{sga3:XXI:eq:11} and
Proposition~\ref{sga3:XXI:2.3.1} give
\[
 (\beta^\ast,\alpha)=(\alpha^\ast,\beta)=\pm1.
\]
If both integers are \(1\), take \(w=s_\beta s_\alpha s_\beta\). If
both are \(-1\), take \(w=s_\alpha s_\beta\).
\end{proof}

\begin{corollary}\label{sga3:XXI:2.3.3}
Let \(\alpha\) and \(\beta\) be roots. If
\(\alpha\neq\beta\) and \((\beta^\ast,\alpha)>0\), then
\(\alpha-\beta\) is a root. Respectively, if
\(\alpha\neq-\beta\) and \((\beta^\ast,\alpha)<0\), then
\(\alpha+\beta\) is a root.
\end{corollary}

\begin{proof}
The second case follows from the first by replacing \(\beta\) by
\(-\beta\). Suppose first that \(\alpha\) and \(\beta\) are
proportional and that \((\beta^\ast,\alpha)>0\). Then
\[
 \beta=\alpha,\qquad 2\beta=\alpha,\qquad\text{or}\qquad
 \beta=2\alpha.
\]
The first case is excluded. In the other two cases,
\(\alpha-\beta=\beta\in R\) and
\(\alpha-\beta=-\alpha\in R\), respectively.

If \(\alpha\) and \(\beta\) are not proportional, then
\((\alpha^\ast,\beta)\) and \((\beta^\ast,\alpha)\) are strictly
positive integers whose product is at most \(3\). Thus one of them
equals \(1\). If \((\beta^\ast,\alpha)=1\), then
\[
 \alpha-\beta=s_\beta(\alpha)\in R;
\]
if \((\alpha^\ast,\beta)=1\), then
\[
 \alpha-\beta=-s_\alpha(\beta)\in R.
\]
\end{proof}

\begin{lemma}\label{sga3:XXI:2.3.4}
Let \(\alpha\) and \(\beta\) be two nonproportional roots. If
\(\beta-\alpha\) is not a root, then
\(\beta+k\alpha\in R\) for \(k=-(\alpha^\ast,\beta)\), but not for
\(k=-(\alpha^\ast,\beta)+1\).
\end{lemma}

\begin{proof}
Indeed,
\[
 \beta-(\alpha^\ast,\beta)\alpha=s_\alpha(\beta)\in R,
\]
whereas
\[
 \beta+\bigl(-(\alpha^\ast,\beta)+1\bigr)\alpha
 =s_\alpha(\beta-\alpha)\notin R.
\]
\end{proof}

\begin{proposition}\label{sga3:XXI:2.3.5}
Let \(\alpha\) and \(\beta\) be two nonproportional roots. The set of
integers \(k\) for which \(\beta+k\alpha\in R\) is an interval
\([p,q]\), with \(p\leq0\) and \(q\geq0\), and
\[
 p+q=-(\alpha^\ast,\beta).
\]
\end{proposition}

\begin{proof}
For the first assertion it suffices, for example, to show that if
\(\beta+k\alpha\in R\) and \(k>0\) is an integer, then
\(\beta+(k-1)\alpha\in R\). This is trivial if \(k=1\). If
\(k\geq2\), then
\[
 (\alpha^\ast,\beta+k\alpha)
 =(\alpha^\ast,\beta)+2k\geq-3+4>0,
\]
and the assertion follows from Corollary~\ref{sga3:XXI:2.3.3}.

Let \([p,q]\) be the interval in question. Applying
Lemma~\ref{sga3:XXI:2.3.4} to \(\beta+p\alpha\), one obtains
\[
 q-p
 =-(\alpha^\ast,\beta+p\alpha)
 =-(\alpha^\ast,\beta)-2p.
\]
\footnote{Editors' note: in this equality, \(+2p\) was corrected to
\(-2p\).}
The asserted formula follows.
\end{proof}

\begin{remark}\label{sga3:XXI:2.3.6}
The preceding formula contains the qualitative statements of
Lemma~\ref{sga3:XXI:2.2.5} and
Corollary~\ref{sga3:XXI:2.3.3}.
\end{remark}

\begin{complements}\label{sga3:XXI:2.3.7}
\footnote{Editors' note: these complements, useful for the proof of
Proposition~\ref{sga3:XXI:3.1.5}, were added by the editors.}
By formulas \eqref{sga3:XXI:eq:9} and
\eqref{sga3:XXI:eq:10}, and by
Corollary~\ref{sga3:XXI:1.2.6}, the bilinear form on \(\VR(R)\)
defined by
\[
 \langle x,y\rangle=(p(x),y)
\]
is symmetric and positive definite. For every \(\alpha\in R\) and
\(y\in\VR(R)\),
\[
 \langle\alpha,y\rangle
 =\frac{\ell(\alpha)}2(\alpha^\ast,y),
 \qquad
 \ell(\alpha)=\langle\alpha,\alpha\rangle.
\]
Consequently, Corollary~\ref{sga3:XXI:2.3.3} gives the following
corollary.

\medskip
\noindent\textbf{Corollary.}\quad
\SGARefTarget{sga3:XXI:2.3.7:corollary}
Let \(\alpha\neq\beta\) in \(R\). If \(\alpha-\beta\notin R\), then
\[
 \langle\alpha,\beta\rangle\leq0.
\]
\end{complements}

\section{Simple roots, positive roots}

\subsection{Systems of simple roots}

\begin{lemma}\label{sga3:XXI:3.1.1}
Let \(\alpha\) and \(\alpha_i\), \(i=1,\ldots,n\), be roots, and
suppose that \(\alpha\) is distinct from the \(\alpha_i\). If
\[
 \alpha=\sum_{i=1}^n q_i\alpha_i,
 \qquad q_i\in\Qp,
\]
then there is an index \(i\) such that
\[
 q_i\neq0,\qquad(\alpha^\ast,\alpha_i)>0,\qquad
 \alpha-\alpha_i\in R.
\]
\end{lemma}

\begin{proof}
Write
\[
 2=(\alpha^\ast,\alpha)
   =\sum_{i=1}^n q_i(\alpha^\ast,\alpha_i).
\]
This proves the first two assertions. The third then follows from
Corollary~\ref{sga3:XXI:2.3.3}.
\end{proof}

% END inputs/XXI/tex/components/03_expose_XXI_remark224_through_lemma311_en.tex

% BEGIN inputs/XXI/tex/components/04_expose_XXI_proposition312_through_remark319_en.tex
\begin{proposition}\label{sga3:XXI:3.1.2}
Let \(\alpha\) and \(\alpha_i\), \(i=1,\ldots,n\), be roots. If
\[
 \alpha=\sum_{i=1}^n m_i\alpha_i,
 \qquad m_i\in\mathbb N,
\]
then there is a sequence \(\beta_1,\ldots,\beta_m\) of roots chosen
from among the \(\alpha_i\), not necessarily pairwise distinct, such
that, on setting
\[
 \gamma_p=\sum_{i=1}^p\beta_i,\qquad p=1,\ldots,m,
\]
one has \(\gamma_p\in R\) and \(\gamma_m=\alpha\).
\end{proposition}

\begin{proof}
We argue by induction on the integer
\(m'=\sum_i m_i\). If \(\alpha\) equals one of the \(\alpha_i\), say
\(\alpha_{i_0}\), which necessarily happens when \(m'=1\), take
\[
 m=1,\qquad\beta_1=\alpha_{i_0}.
\]
Otherwise, Lemma~\ref{sga3:XXI:3.1.1} gives an index \(i\) for which
\(m_i\neq0\) and \(\alpha-\alpha_i\) is a root. Then
\(m_i-1\in\mathbb N\) and
\[
 \alpha-\alpha_i
 =(m_i-1)\alpha_i+\sum_{j\neq i}m_j\alpha_j.
\]
Apply the induction hypothesis to \(\alpha-\alpha_i\), and append
\(\alpha_i\) to the resulting sequence.
\end{proof}

\begin{corollary}\label{sga3:XXI:3.1.3}
Let \(R'\subset R\). The following conditions are equivalent:
\begin{enumerate}
\item[\SGAInlineRefTarget{sga3:XXI:3.1.3:item:i}(i)] if \(\alpha,\beta\in R'\) and \(\alpha+\beta\in R\), then
  \(\alpha+\beta\in R'\);
\item[\SGAInlineRefTarget{sga3:XXI:3.1.3:item:ii}(ii)] \((\mathbb N\cdot R')\cap R\subset R'\).
\end{enumerate}
\end{corollary}

\begin{proof}
Clearly \SGARefLink{sga3:XXI:3.1.3:item:ii}{(ii)} implies \SGARefLink{sga3:XXI:3.1.3:item:i}{(i)}. The converse follows immediately from
Proposition~\ref{sga3:XXI:3.1.2}.
\end{proof}

\begin{definition}\label{sga3:XXI:3.1.4}
A set of roots satisfying the equivalent conditions of
Corollary~\ref{sga3:XXI:3.1.3} is called \emph{closed}.
\end{definition}

\begin{proposition}\label{sga3:XXI:3.1.5}
Let \(\Delta\subset R\) be a set of roots. The following conditions
are equivalent:
\begin{enumerate}
\item[\SGAInlineRefTarget{sga3:XXI:3.1.5:item:i}(i)] the elements of \(\Delta\) are indivisible and linearly
  independent, and
  \[
   R\subset(\Qp\cdot\Delta)\cup(-\Qp\cdot\Delta);
  \]
\item[\SGAInlineRefTarget{sga3:XXI:3.1.5:item:ii}(ii)] the elements of \(\Delta\) are linearly independent, and
  \[
   R\subset(\mathbb N\cdot\Delta)\cup(-\mathbb N\cdot\Delta);
  \]
\item[\SGAInlineRefTarget{sga3:XXI:2.3.7:corollary:item:iii}(iii)] every root can be written uniquely as a linear
  combination of the elements of \(\Delta\), with integer
  coefficients all of the same sign.
\end{enumerate}
\end{proposition}

\begin{proof}
\footnote{Editors' note: in the argument that follows, the order was
changed and the proof of the implication \SGARefLink{sga3:XXI:2.3.7:corollary:item:iii}{(iii)}\(\Rightarrow\)\SGARefLink{sga3:XXI:3.1.5:item:i}{(i)} was
spelled out.}
It is clear that \SGARefLink{sga3:XXI:3.1.5:item:ii}{(ii)} implies \SGARefLink{sga3:XXI:2.3.7:corollary:item:iii}{(iii)}. Write
\(\alpha_1,\ldots,\alpha_n\) for the distinct elements of \(\Delta\).

Assume \SGARefLink{sga3:XXI:3.1.5:item:i}{(i)}, and let \(\alpha\in R\). There is a unique expression
\[
 \varepsilon\alpha=\sum_i q_i\alpha_i,
 \qquad q_i\in\Qp,\qquad\varepsilon=\pm1.
\]
We show that the \(q_i\) are integers. This is immediate when all but
one of them vanish, by Corollary~\ref{sga3:XXI:2.1.5}. Otherwise,
\(\varepsilon\alpha\) is distinct from the \(\alpha_i\), and
Lemma~\ref{sga3:XXI:3.1.1} gives an \(i_0\) such that
\[
 \alpha'=\varepsilon\alpha-\alpha_{i_0}\in R,
 \qquad q_{i_0}\neq0.
\]
Thus
\[
 \alpha'
 =(q_{i_0}-1)\alpha_{i_0}
  +\sum_{i\neq i_0}q_i\alpha_i.
\]
At least one \(q_i\), with \(i\neq i_0\), is nonzero, so \SGARefLink{sga3:XXI:3.1.5:item:i}{(i)} forces
\(q_{i_0}-1\geq0\). Repeating this operation with \(\alpha'\), after
finitely many steps we find that all \(q_i\) are integers. This proves
\SGARefLink{sga3:XXI:3.1.5:item:ii}{(ii)}.

Assume now \SGARefLink{sga3:XXI:2.3.7:corollary:item:iii}{(iii)}. We first prove that
\(\alpha_1,\ldots,\alpha_n\) are linearly independent over
\(\mathbb Q\). Otherwise there would be an equality
\[
 x=\sum_{i\in I}a_i\alpha_i
  =\sum_{j\in J}b_j\alpha_j,
\]
where \(I\) and \(J\) are disjoint subsets of
\(\{1,\ldots,n\}\), at least one of them, say \(I\), is nonempty, and
\(a_i,b_j\in\mathbb N^\ast\). By the corollary in
Complements~\ref{sga3:XXI:2.3.7}, one has
\(\langle\alpha_i,\alpha_j\rangle\leq0\) for \(i\in I\) and \(j\in J\).
Therefore \(\langle x,x\rangle\leq0\), and hence \(x=0\).

Choose \(i_0\in I\). The two expressions
\[
 \alpha_{i_0}
 =\alpha_{i_0}+\sum_{i\in I}a_i\alpha_i
 =\alpha_{i_0}+\sum_{j\in J}b_j\alpha_j
\]
then imply, by \SGARefLink{sga3:XXI:2.3.7:corollary:item:iii}{(iii)}, that \(a_i=0=b_j\) for all \(i,j\), a
contradiction. Thus the elements of \(\Delta\) are linearly
independent over \(\mathbb Q\).

It remains to show that they are indivisible. Suppose that
\(\beta\in\Delta\) is divisible. Then \(\beta/2\in R\), so
\[
 \frac{\beta}{2}
 =\varepsilon\sum_{\alpha\in\Delta}m_\alpha\alpha,
 \qquad m_\alpha\in\mathbb N,\qquad\varepsilon=\pm1.
\]
Consequently,
\[
 \beta=2\varepsilon\sum_{\alpha\in\Delta}m_\alpha\alpha.
\]
Uniqueness gives \(m_\alpha=0\) for \(\alpha\neq\beta\) and
\(2\varepsilon m_\beta=1\), a contradiction.
\end{proof}

\begin{definition}\label{sga3:XXI:3.1.6}
A set \(\Delta\) of roots satisfying the conditions of
Proposition~\ref{sga3:XXI:3.1.5} is called a
\emph{system of simple roots}, or a \emph{base} of \(\cR\).
\end{definition}

If \(w\in W\) and \(\Delta\) is a system of simple roots, then
\(w(\Delta)\) is a system of simple roots.

\begin{remark}\label{sga3:XXI:3.1.7}
This definition involves only \(R\), not \(\cR\); indeed, it involves
only \(\operatorname{ind}(R)\).
\end{remark}

\begin{remark}\label{sga3:XXI:3.1.8}
If \(\Delta\) is a system of simple roots, then \(\Delta\) is a basis
of the free abelian group \(\Gamma_0(R)\). Hence
\[
 \operatorname{Card}(\Delta)=\rgss(\cR).
\]
\end{remark}

\begin{remark}\label{sga3:XXI:3.1.9}
The conditions of Proposition~\ref{sga3:XXI:3.1.5} are therefore
equivalent to the following:
\begin{enumerate}
\item[(i\('\))] the elements of \(\Delta\) are indivisible, their
  number is at most \(\rgss(\cR)\), and
  \[
   R\subset(\Qp\cdot\Delta)\cup(-\Qp\cdot\Delta);
  \]
\item[(ii\('\))] \(\operatorname{Card}(\Delta)\leq\rgss(\cR)\), and
  \[
   R\subset(\mathbb N\cdot\Delta)\cup(-\mathbb N\cdot\Delta).
  \]
\end{enumerate}
\end{remark}

% END inputs/XXI/tex/components/04_expose_XXI_proposition312_through_remark319_en.tex

% BEGIN inputs/XXI/tex/components/05_expose_XXI_corollary3110_through_definition325_en.tex
\begin{corollary}\label{sga3:XXI:3.1.10}
If \(\Delta\) is a system of simple roots, then
\(\operatorname{ind}(\Delta^\ast)\) is a system of simple coroots
\emph{(i.e.}\ a system of simple roots of \(\cR^\ast\)\emph{)}.
\end{corollary}

\begin{proof}
If
\[
 \beta=\sum_{\alpha\in\Delta}a_\alpha\alpha,
 \qquad a_\alpha\in\mathbb N,
\]
then
\[
 p(\beta)=\sum_{\alpha\in\Delta}a_\alpha p(\alpha).
\]
Consequently, by \eqref{sga3:XXI:eq:10},
\[
 \beta^\ast
 =\sum_{\alpha\in\Delta}
   a_\alpha\frac{\ell(\alpha)}{\ell(\beta)}\alpha^\ast,
\]
which shows that \(\operatorname{ind}(\Delta^\ast)\) satisfies
condition (i\('\)) of Remark~\ref{sga3:XXI:3.1.9}.
\end{proof}

By Proposition~\ref{sga3:XXI:2.1.1}, if \(\alpha^\ast\) is not
indivisible, then
\[
 \alpha^\ast=2u^\ast,
 \qquad
 u^\ast\in\operatorname{ind}(\Delta^\ast)
 \quad\text{and}\quad
 u=2\alpha.
\]
It follows that:\footnote{Editors' note: The preceding sentence was
added, and in Corollary 3.1.11 the expression
\(4a_\alpha\ell(\alpha)\) was replaced by
\(2a_\alpha\ell(\alpha)\).}

\begin{corollary}\label{sga3:XXI:3.1.11}
Let \(\Delta\) be a system of simple roots, and let
\[
 \beta=\sum_{\alpha\in\Delta}a_\alpha\alpha,
 \qquad a_\alpha\in\mathbb Z,
\]
be the expression of \(\beta\) relative to \(\Delta\). Then
\(\ell(\beta)\) divides \(2a_\alpha\ell(\alpha)\); it even divides
\(a_\alpha\ell(\alpha)\) if \(\alpha^\ast\) is indivisible.
\end{corollary}

Before continuing to state the properties of systems of simple roots,
we show that such systems exist.

\subsection{Systems of positive roots}

\begin{definition}\label{sga3:XXI:3.2.1}
A subset \(R_+\subset R\) is called a \emph{system of positive roots}
of \(\cR\) \emph{(or of \(R\); see
Remark~\ref{sga3:XXI:3.2.2}\emph{)}} if it satisfies:
\begin{enumerate}
\item[(P 1)] \(R_+\cap(-R_+)=\varnothing\);
\item[(P 2)] \(R_+\cup(-R_+)=R\);
\item[(P 3)] \((\Qp\cdot R_+)\cap R\subset R_+\).
\end{enumerate}
In particular, such a set is closed. We shall see later, in
Corollary~\ref{sga3:XXI:3.3.8}, that a closed set satisfying
(P 1) and (P 2) also satisfies (P 3), and hence is a system of
positive roots. If \(w\in W\) and \(R_+\) is a system of positive
roots, then \(w(R_+)\) is a system of positive roots.
\end{definition}

\begin{remark}\label{sga3:XXI:3.2.2}
This definition involves only \(R\). We shall also say that \(R_+\)
is a system of positive roots of \(R\).
\end{remark}

\begin{remark}\label{sga3:XXI:3.2.3}
Conditions (P 1) and (P 2) immediately give
\[
 \operatorname{Card}(R_+)=\frac{\operatorname{Card}(R)}{2}.
\]
It follows that if \(R_+\) and \(R'_+\) are two systems of positive
roots and \(R_+\subset R'_+\), then \(R_+=R'_+\).
\end{remark}

\begin{remark}\label{sga3:XXI:3.2.4}
If \(R_+\) is a system of positive roots, then \(R^\ast_+\) is a
system of positive coroots \emph{(i.e.}\ a system of positive roots
of \(R^\ast\)\emph{)}.
\end{remark}

This follows immediately from Lemma~\ref{sga3:XXI:1.1.4} and
Scholium~\ref{sga3:XXI:1.2.2}.

\begin{definition}\label{sga3:XXI:3.2.5}
Let \(\Delta\) be a system of simple roots. Set
\[
 \mathscr{P}(\Delta)
 =(\Qp\cdot\Delta)\cap R
 =(\mathbb N_+\cdot\Delta)\cap R.
\]
\end{definition}

% END inputs/XXI/tex/components/05_expose_XXI_corollary3110_through_definition325_en.tex

% BEGIN inputs/XXI/tex/components/06_expose_XXI_proposition326_through_lemma3210_en.tex
\begin{proposition}\label{sga3:XXI:3.2.6}
If \(\Delta\) is a system of simple roots, then
\(\mathscr{P}(\Delta)\) is a system of positive roots. If \(\Delta\)
is a system of simple roots and \(R_+\) is a system of positive roots,
then
\[
 \Delta\subset R_+
 \quad\Longleftrightarrow\quad
 R_+=\mathscr{P}(\Delta).
\]
\end{proposition}

\begin{proof}
The first assertion is immediate. If \(\Delta\subset R_+\), then
\(\mathscr{P}(\Delta)\subset R_+\) by (P 3), and therefore
\(\mathscr{P}(\Delta)=R_+\) by
Remark~\ref{sga3:XXI:3.2.3}. The rest is trivial.
\end{proof}

\begin{remark}\label{sga3:XXI:3.2.7}
Systems of positive roots exist: endow \(\VR(R)\) with the structure
of a totally ordered vector space. The set of roots that are
\(\geq0\) for this order is a system of positive roots.
\end{remark}

\begin{theorem}\label{sga3:XXI:3.2.8}
Let \(R_+\) be a system of positive roots. There exists a unique
system of simple roots \(\mathscr{S}(R_+)\) such that
\[
 \mathscr{S}(R_+)\subset R_+,
 \qquad\text{i.e.}\qquad
 R_+=\mathscr{P}\bigl(\mathscr{S}(R_+)\bigr).
\]
\end{theorem}

Uniqueness follows immediately from:

\begin{lemma}\label{sga3:XXI:3.2.9}
Let \(\Delta\) be a system of simple roots. An element
\(\alpha\in\mathscr{P}(\Delta)\) belongs to \(\Delta\) if and only if
\(\alpha\) is not the sum of two elements of
\(\mathscr{P}(\Delta)\).
\end{lemma}

\begin{proof}
This follows immediately from the definitions and
Proposition~\ref{sga3:XXI:3.1.2}.
\end{proof}

We now prove the existence of \(\mathscr{S}(R_+)\). Consider the set
of subsets \(T\) of \(R_+\) such that
\[
 T=\operatorname{ind}(T)
 \qquad\text{and}\qquad
 \Qp\cdot T\supset R_+.
\]
This set is nonempty, since it contains
\(\operatorname{ind}(R_+)\). Let \(\Delta\) be a minimal element of
this set under inclusion. We show that \(\Delta\) is a system of
simple roots; by Proposition~\ref{sga3:XXI:3.1.5}\SGARefLink{sga3:XXI:3.1.5:item:i}{(i)}, it is enough to
show that \(\Delta\) is a linearly independent subset of
\(M\otimes\mathbb Q\).

\begin{lemma}\label{sga3:XXI:3.2.10}
If \(\alpha,\beta\in\Delta\) and
\[
 q\alpha+q'\beta\in R,
 \qquad q,q'\in\mathbb Q,
\]
then \(qq'\geq0\).
\end{lemma}

\begin{proof}
If \(qq'<0\), then, after interchanging \(\alpha\) and \(\beta\) if
necessary, we may write
\[
 a\alpha-b\beta\in R_+,
 \qquad a,b\in\mathbb Q^\ast_+.
\]
By the defining property of \(\Delta\), there is a relation
\[
 a\alpha-b\beta
   =\sum_{\gamma\in\Delta}c(\gamma)\gamma,
 \qquad c(\gamma)\in\Qp.
\]
If \(a\leq c(\alpha)\), this becomes
\[
 -\beta
 =\frac{c(\alpha)-a}{b}\alpha
  +\sum_{\gamma\neq\alpha}\frac{c(\gamma)}{b}\gamma.
\]
Thus \(-\beta\in(\Qp\cdot\Delta)\cap R\), which is contained in
\(R_+\) by (P 3). This gives
\(\beta\in R_+\cap(-R_+)\), contradicting (P 1).

If instead \(a>c(\alpha)\), then
\[
 \bigl(a-c(\alpha)\bigr)\alpha
 =b\beta+\sum_{\gamma\neq\alpha}c(\gamma)\gamma,
\]
which shows that
\[
 \Delta\subset\Qp\cdot\bigl(\Delta\setminus\{\alpha\}\bigr),
\]
contrary to the minimality of \(\Delta\).
\end{proof}

% END inputs/XXI/tex/components/06_expose_XXI_proposition326_through_lemma3210_en.tex

% BEGIN inputs/XXI/tex/components/07_expose_XXI_lemma3211_through_lemma3216_en.tex
Recall from Complements~\ref{sga3:XXI:2.3.7} that the bilinear form on
\(\VR(R)\) defined by
\[
 \langle x,y\rangle=(p(x),y)
\]
is a Euclidean inner product. Moreover, for \(\alpha\in R\),
\[
 \langle\alpha,y\rangle
 =\frac{\ell(\alpha)}{2}(\alpha^\ast,y),
\]
so \(\langle\alpha,y\rangle\) and \((\alpha^\ast,y)\) have the same
sign.

\begin{lemma}\label{sga3:XXI:3.2.11}
If \(\alpha,\beta\in\Delta\), then
\[
 (\beta^\ast,\alpha)\leq0
 \qquad\text{and hence}\qquad
 \langle\beta,\alpha\rangle\leq0.
\]
\end{lemma}

\begin{proof}
Indeed,
\[
 s_\beta(\alpha)
 =\alpha-(\beta^\ast,\alpha)\beta\in R,
\]
and therefore \((\beta^\ast,\alpha)\leq0\) by
Lemma~\ref{sga3:XXI:3.2.10}.
\end{proof}

We now show that \(\Delta\) is linearly independent. Otherwise,
Lemma~\ref{sga3:XXI:3.2.11} would imply, exactly as in the proof of
Proposition~\ref{sga3:XXI:3.1.5}, the existence of a nontrivial
relation
\[
 \sum_{\alpha\in\Delta}m(\alpha)\alpha=0,
 \qquad m(\alpha)\in\mathbb N.
\]
If \(m(\alpha_0)\geq1\), then
\[
 -\alpha_0
 =\bigl(m(\alpha_0)-1\bigr)\alpha_0
  +\sum_{\alpha\neq\alpha_0}m(\alpha)\alpha.
\]
By (P 3), this gives
\(\alpha_0\in R_+\cap(-R_+)\), contradicting (P 1).
Thus \(\Delta\) is a base of \(\cR\), completing the proof of
Theorem~\ref{sga3:XXI:3.2.8}.

\begin{corollary}\label{sga3:XXI:3.2.12}
Let \(R_+\) be a system of positive roots, let \(\Delta\) be a base of
\(\cR\), and let \(w\in W\). Then
\[
 \Delta\subset R_+
 \quad\Longleftrightarrow\quad
 R_+=\mathscr{P}(\Delta)
 \quad\Longleftrightarrow\quad
 \Delta=\mathscr{S}(R_+),
\]
\[
 \mathscr{P}\bigl(\operatorname{ind}(\Delta^\ast)\bigr)
   =\mathscr{P}(\Delta)^\ast,
 \qquad
 \mathscr{S}(R^\ast_+)
   =\operatorname{ind}\bigl(\mathscr{S}(R_+)^\ast\bigr),
\]
and
\[
 \mathscr{S}\bigl(w(R_+)\bigr)=w\bigl(\mathscr{S}(R_+)\bigr),
 \qquad
 \mathscr{P}\bigl(w(\Delta)\bigr)=w\bigl(\mathscr{P}(\Delta)\bigr).
\]
\end{corollary}

\begin{definition}\label{sga3:XXI:3.2.13}
Once a system of simple roots \(\Delta\) has been chosen, the elements
of \(\mathscr{P}(\Delta)\) are called \emph{positive}. Once a system
of positive roots \(R_+\) has been chosen, the elements of
\(\mathscr{S}(R_+)\) are called \emph{simple}.
\end{definition}

\begin{corollary}\label{sga3:XXI:3.2.14}
Let \(R_+\) be a system of positive roots, and let \(\alpha\in R_+\).
The following conditions are equivalent:
\begin{enumerate}
\item[\SGAInlineRefTarget{sga3:XXI:3.2.14:item:i}(i)] \(\alpha\) is simple, that is,
  \(\alpha\in\mathscr{S}(R_+)\);
\item[\SGAInlineRefTarget{sga3:XXI:3.2.14:item:ii}(ii)] \(\alpha\) is not the sum of two elements of \(R_+\);
\item[(iii)] \(R_+\setminus\{\alpha\}\) is closed.
\end{enumerate}
\end{corollary}

\begin{proof}
The equivalence of \SGARefLink{sga3:XXI:3.2.14:item:i}{(i)} and \SGARefLink{sga3:XXI:3.2.14:item:ii}{(ii)} follows immediately from
Lemma~\ref{sga3:XXI:3.2.9}; the equivalence of \SGARefLink{sga3:XXI:3.2.14:item:ii}{(ii)} and \SGARefLink{sga3:XXI:2.3.7:corollary:item:iii}{(iii)} is
immediate from Definition~\ref{sga3:XXI:3.1.4}.
\end{proof}

\begin{definition}\label{sga3:XXI:3.2.15}
Let \(\Delta\) be a system of simple roots. The sum of the
coefficients in the expression of a root \(\alpha\) relative to
\(\Delta\) is called the \emph{order of \(\alpha\) relative to
\(\Delta\)} and is denoted \(\operatorname{ord}_\Delta(\alpha)\).
\end{definition}

One has
\[
 \alpha\in\mathscr{P}(\Delta)
 \quad\Longleftrightarrow\quad
 \operatorname{ord}_\Delta(\alpha)>0,
 \qquad
 \alpha\in\Delta
 \quad\Longleftrightarrow\quad
 \operatorname{ord}_\Delta(\alpha)=1.
\]

\begin{lemma}\label{sga3:XXI:3.2.16}
Let \(\Delta\) be a system of simple roots and
\(\alpha\in\mathscr{P}(\Delta)\). There is a sequence
\(\alpha_i\in\Delta\), \(i=1,\ldots,m\), such that
\[
 \alpha_1+\alpha_2+\cdots+\alpha_p\in\mathscr{P}(\Delta)
 \quad\text{for }p=1,\ldots,m,
\]
and
\[
 \alpha_1+\alpha_2+\cdots+\alpha_m=\alpha.
\]
Moreover, for every sequence satisfying these conditions,
\[
 m=\operatorname{ord}_\Delta(\alpha).
\]
\end{lemma}

\begin{proof}
This is immediate from Proposition~\ref{sga3:XXI:3.1.2}.
\end{proof}

% END inputs/XXI/tex/components/07_expose_XXI_lemma3211_through_lemma3216_en.tex

% BEGIN inputs/XXI/tex/components/08_expose_XXI_lemma331_through_corollary339_en.tex
\subsection{Characterization and conjugacy of systems of positive roots}

\begin{lemma}\label{sga3:XXI:3.3.1}
If \(\alpha\in\mathscr{P}(\Delta)\), \(\beta\in\Delta\), and
\(\alpha\) and \(\beta\) are not proportional, then
\[
 s_\beta(\alpha)\in\mathscr{P}(\Delta).
\]
\end{lemma}

\begin{proof}
One has
\[
 s_\beta(\alpha)=\alpha-(\beta^\ast,\alpha)\beta.
\]
At least one simple root other than \(\beta\) occurs with a nonzero,
and hence strictly positive, coefficient in the expression of
\(\alpha\). It occurs with the same coefficient in the expression of
\(s_\beta(\alpha)\); therefore \(s_\beta(\alpha)\) is also positive.
\end{proof}

\begin{corollary}\label{sga3:XXI:3.3.2}
If \(\beta\in\Delta\), then the reflection \(s_\beta\) permutes the
elements of \(\mathscr{P}(\Delta)\) that are not proportional to
\(\beta\).
\end{corollary}

\begin{lemma}\label{sga3:XXI:3.3.3}
Let
\[
 \alpha\in\mathscr{P}(\Delta)\setminus\Delta
\]
be indivisible. There exists \(\beta\in\Delta\) such that
\[
 s_\beta(\alpha)\in\mathscr{P}(\Delta)
 \qquad\text{and}\qquad
 \operatorname{ord}_\Delta\bigl(s_\beta(\alpha)\bigr)
   <\operatorname{ord}_\Delta(\alpha).
\]
\end{lemma}

\begin{proof}
By Lemma~\ref{sga3:XXI:3.1.1}, there is a \(\beta\in\Delta\) such
that \((\beta^\ast,\alpha)>0\). Since \(\alpha\notin\Delta\) and
\(\alpha\) is indivisible, it cannot be proportional to \(\beta\).
Thus \(s_\beta(\alpha)\in\mathscr{P}(\Delta)\) by
Lemma~\ref{sga3:XXI:3.3.1}, and
\[
 \operatorname{ord}_\Delta\bigl(s_\beta(\alpha)\bigr)
 =\operatorname{ord}_\Delta(\alpha)-(\beta^\ast,\alpha)
 <\operatorname{ord}_\Delta(\alpha).
\]
\end{proof}

\begin{corollary}\label{sga3:XXI:3.3.4}
If \(\alpha\in\mathscr{P}(\Delta)\) is indivisible, then there is a
sequence \(\beta_p\in\Delta\), \(p=1,\ldots,q\), such that
\[
 \alpha_p
 =s_{\beta_p}\cdots s_{\beta_1}(\alpha)
 \in\mathscr{P}(\Delta)
 \quad\text{for }p=1,\ldots,q,
\]
and \(\alpha_q\in\Delta\).
\end{corollary}

\begin{proof}
Apply Lemma~\ref{sga3:XXI:3.3.3} repeatedly, inducting on
\(\operatorname{ord}_\Delta(\alpha)\).
\end{proof}

\begin{proposition}\label{sga3:XXI:3.3.5}
The Weyl group is generated by the \(s_\alpha\), for
\(\alpha\in\Delta\). Every indivisible root is conjugate under an
element of the Weyl group to a simple root.
\end{proposition}

\begin{proof}
The second assertion follows immediately from
Corollary~\ref{sga3:XXI:3.3.4}. The first then follows from
Corollary~\ref{sga3:XXI:1.2.10} and
Proposition~\ref{sga3:XXI:2.1.1}.
\end{proof}

\begin{proposition}\label{sga3:XXI:3.3.6}
Let \(R_+\) be a system of positive roots, and let \(P'\subset R\) be
closed and satisfy (P 2). Then there exists \(w\in W\) such that
\[
 w(R_+)\subset P'.
\]
\end{proposition}

We first state the corollaries.

\begin{corollary}\label{sga3:XXI:3.3.7}
The Weyl group acts transitively on the set of systems of positive
roots, and on the set of systems of simple roots.
\end{corollary}

\begin{corollary}\label{sga3:XXI:3.3.8}
A subset of \(R\) is a system of positive roots if and only if it is
closed and satisfies (P 1) and (P 2).
\end{corollary}

\begin{corollary}\label{sga3:XXI:3.3.9}
Suppose that \(\Gamma_0(R)\) is endowed with an ordered-group
structure for which every root is either \(>0\) or \(<0\). Then the
set of positive roots for this order is a system of positive roots.
\end{corollary}

% END inputs/XXI/tex/components/08_expose_XXI_lemma331_through_corollary339_en.tex

% BEGIN inputs/XXI/tex/components/09_expose_XXI_proof336_through_definition33101_en.tex
\begin{proof}[Proof of Proposition~\ref{sga3:XXI:3.3.6}]
Choose \(w\in W\) such that
\(\operatorname{Card}(w(R_+)\cap P')\) is maximal. Replacing \(R_+\)
by \(w(R_+)\), we may assume that
\begin{equation}
 \operatorname{Card}(R_+\cap P')
 \geq
 \operatorname{Card}\bigl(R_+\cap s_\alpha(P')\bigr)
 \quad\text{for every }\alpha\in\mathscr{S}(R_+)=\Delta.
 \tag{\(\ast\)}\label{sga3:XXI:3.3.6:star}
\end{equation}
We show that \(R_+\subset P'\). Otherwise, since \(P'\) is closed,
there is an \(\alpha\in\Delta\) such that \(\alpha\notin P'\). As
\(P'\) satisfies (P 2), one then has \(-\alpha\in P'\), and also
\(-2\alpha\in P'\) if \(2\alpha\) is a root.

For every \(\beta\in R_+\cap P'\), one has \(\beta\neq\alpha\). If
\(2\alpha\) is not a root, then
\(s_\alpha(\beta)\in R_+\) by Lemma~\ref{sga3:XXI:3.3.1}; hence
\[
 s_\alpha(R_+\cap P')
 \subset R_+\cap s_\alpha(P').
\]
But \(R_+\cap s_\alpha(P')\) also contains
\(\alpha=s_\alpha(-\alpha)\), contradicting
\eqref{sga3:XXI:3.3.6:star}. If \(2\alpha\) is a root, the same
argument applies.
\end{proof}

Thus, to study closed sets of roots satisfying (P 2), it is enough to
consider those that contain a system of positive roots.

\begin{proposition}\label{sga3:XXI:3.3.10}
Let \(R_+\) be a system of positive roots, and let \(P'\) be a closed
subset of \(R\) containing \(R_+\). Put
\[
 \Delta=\mathscr{S}(R_+),
 \qquad
 \Delta'=\Delta\cap(-P').
\]
Then \(P'\) is the union of \(R_+\) and the set of roots that are
linear combinations of the elements of \(\Delta'\) with negative
coefficients.
\end{proposition}

\begin{proof}
We argue by induction on the order of a root
\(\gamma\in P'\setminus R_+\). If
\(\operatorname{ord}_\Delta(\gamma)=-1\), then
\(-\gamma\in\Delta\), and hence \(\gamma\in-\Delta'\).

Suppose that \(\operatorname{ord}_\Delta(\gamma)<-1\). By
Proposition~\ref{sga3:XXI:3.1.2}, there exists \(\beta\in\Delta\)
such that \(-\gamma-\beta\in R\). Then
\[
 0<\operatorname{ord}_\Delta(-\gamma-\beta)
  =\operatorname{ord}_\Delta(-\gamma)-1
  <\operatorname{ord}_\Delta(-\gamma).
\]
The first inequality shows that \(-\gamma-\beta\in R_+\).
Consequently, since \(R_+\cap(-R_+)=\varnothing\), and since
\(\gamma+\beta\) is the sum of two roots of \(P'\), the root
\(\gamma+\beta\) belongs to \(P'\setminus R_+\) and satisfies
\[
 \operatorname{ord}_\Delta(\gamma+\beta)
 >\operatorname{ord}_\Delta(\gamma).
\]
By induction, \(\gamma+\beta\) is a linear combination of the
elements of \(\Delta'\) with negative coefficients. Since
\(\gamma=(\gamma+\beta)-\beta\), it remains only to show that
\(\beta\in-P'\). But \(\gamma\in P'\) and
\(-(\gamma+\beta)\in R_+\subset P'\), so
\[
 -\beta=\gamma-(\gamma+\beta)\in P'.
\]
\end{proof}

\begingroup
\renewcommand{\thetheorem}{3.3.10.1}
\begin{definition}\label{sga3:XXI:3.3.10.1}
\footnote{Editors' note: This definition was inserted here.}
A subset \(R'\) of \(R\) is called \emph{symmetric} if
\[
 -R'=R'.
\]
\end{definition}
\endgroup
\addtocounter{theorem}{-1}

% END inputs/XXI/tex/components/09_expose_XXI_proof336_through_definition33101_en.tex

% BEGIN inputs/XXI/tex/components/10_expose_XXI_scholium3311_through_remark345_en.tex
\begin{scholium}\label{sga3:XXI:3.3.11}
Let \(P'\) be a set of roots satisfying (P 2) and closed. There exist a
system of simple roots \(\Delta\) and a subset \(\Delta'\) of
\(\Delta\) such that
\[
 P'=R\cap\bigl(\mathbb{N}\cdot\Delta
          \,\cup\, -\mathbb{N}\cdot\Delta'\bigr).
\]
If
\[
 R_{\Delta'}=(\mathbb{Z}\cdot\Delta')\cap R,
\]
which is a closed and symmetric subset of \(R\), then \(P'\) is the
disjoint union of \(R_{\Delta'}\) and the closed subset of
\(\mathscr{P}(\Delta)\) formed by the
\[
 \alpha\in\mathscr{P}(\Delta)\setminus
           \mathbb{N}\cdot\Delta',
\]
that is, the positive roots whose decomposition with respect to
\(\Delta\) ``contains at least one root of
\(\Delta\setminus\Delta'\)''.
\end{scholium}

\subsection{Closed and symmetric sets of roots}
\label{sga3:XXI:3.4}

\begin{proposition}\label{sga3:XXI:3.4.1}
Let
\[
 \cR=(M,M^\ast,R,R^\ast)
\]
be a root datum, and let \(R'\) be a closed and symmetric subset of
\(R\). Then:
\begin{enumerate}
\item[(i)] \(\cR'=(M,M^\ast,R',R'^\ast)\) is a root datum;
\item[(ii)] for every system of positive roots \(R_+\) of \(R\),
  \(R'_+=R_+\cap R'\) is a system of positive roots of \(R'\);
\item[(iii)] the Weyl group \(W(\cR')\) of \(\cR'\) is the subgroup
  of \(W(\cR)\) generated by the \(s_\alpha\), for
  \(\alpha\in R'\).
\end{enumerate}
\end{proposition}

\begin{proof}
The first assertion is immediate from
Corollary~\ref{sga3:XXI:1.2.11}; the second follows from
Corollary~\ref{sga3:XXI:3.3.8}; and the third is evident.
\end{proof}

\begin{corollary}\label{sga3:XXI:3.4.2}
Let \(w\in W(\cR')\). The order of \(w\) is the smallest integer
\(n>0\) such that
\[
 w^n(\alpha')=\alpha'
 \qquad\text{for every }\alpha'\in R'.
\]
\end{corollary}

\begin{proof}
Apply Corollary~\ref{sga3:XXI:1.2.7} to the root datum \(\cR'\).
\end{proof}

\begin{corollary}\label{sga3:XXI:3.4.3}
Let \(\alpha\) and \(\beta\) be two nonproportional roots. Let \(n\)
be the smallest integer \(>0\) such that
\[
 (s_\alpha s_\beta)^n(\alpha)=\alpha,
 \qquad
 (s_\alpha s_\beta)^n(\beta)=\beta.
\]
Then the subgroup \(W_{\alpha,\beta}\) of \(W\) generated by
\(s_\alpha\) and \(s_\beta\) is defined by the relations
\[
 s_\alpha^2=1,\qquad
 s_\beta^2=1,\qquad
 (s_\alpha s_\beta)^n=1.
\]
\end{corollary}

\begin{proof}
In view of Corollary~\ref{sga3:XXI:3.4.2}, it remains only to verify
the following lemma.
\end{proof}

\begin{lemma}\label{sga3:XXI:3.4.4}
Let \(G\) be the group generated by two elements \(x\) and \(y\)
subject to the relations \(x^2=y^2=1\). Every normal subgroup of
\(G\) containing neither \(x\) nor \(y\) is generated, as a normal
subgroup, by an element of the form \((xy)^n\), where \(n>0\).
\end{lemma}

\begin{proof}
\footnote{Editors' note: The original has been corrected in what
follows.}
Every element of \(G\) can be written as \((xy)^n\), as
\((yx)^n=(xy)^{-n}\), or in one of the forms
\[
\begin{array}{c@{\qquad}c@{\qquad}c}
 (a) & x(yx)^{2n} & \text{or } y(xy)^{2n+1},\\[2mm]
 (b) & y(xy)^{2n} & \text{or } x(yx)^{2n+1},
\end{array}
\]
where \(n\in\mathbb{N}\). The elements of type (a), respectively of
type (b), are conjugate to \(x\), respectively to \(y\).
\end{proof}

\begin{remark}\label{sga3:XXI:3.4.5}
The integer \(n\) is computed immediately. Put
\[
 (\alpha^\ast,\beta)=p,\qquad
 (\beta^\ast,\alpha)=q.
\]
Then\footnote{Editors' note: The original has been corrected in what
follows.}
\[
 (s_\alpha s_\beta)(\alpha)=(pq-1)\alpha-q\beta,
 \qquad
 (s_\alpha s_\beta)(\beta)=p\alpha-\beta.
\]
Thus \(n\) is the order of the matrix
\[
 \begin{pmatrix}
  pq-1 & p\\
  -q   & -1
 \end{pmatrix}.
\]
If \(pq=0\), then \(p=q=0\) by
Corollary~\ref{sga3:XXI:2.2.2}, whence \(n=2\). Otherwise,
Proposition~\ref{sga3:XXI:2.3.1} shows that \(pq\) equals \(1\),
\(2\), or \(3\), and one finds respectively \(n=3\), \(4\), or
\(6\).

\smallskip
\noindent
\emph{N.B.} Expressing the finiteness of the order of the preceding
matrix recovers inequality \SGARefLink{sga3:XXI:eq:13}{(13)} of
Proposition~\ref{sga3:XXI:2.3.1}.
\end{remark}

% END inputs/XXI/tex/components/10_expose_XXI_scholium3311_through_remark345_en.tex

% BEGIN inputs/XXI/tex/components/11_expose_XXI_definition346_through_lemma349_en.tex
\begin{definition}\label{sga3:XXI:3.4.6}
Let \(\Delta\) be a system of simple roots and let
\(\Delta'\subset\Delta\). We write
\[
 R_{\Delta'}
   =R\cap(\mathbb{Q}\cdot\Delta')
   =R\cap(\mathbb{Z}\cdot\Delta').
\]
\end{definition}

\begin{lemma}\label{sga3:XXI:3.4.7}
The set \(R_{\Delta'}\) is closed and symmetric, and \(\Delta'\) is
a system of simple roots of the root datum
\[
 (M,M^\ast,R_{\Delta'},R_{\Delta'}^\ast).
\]
Its Weyl group is the subgroup \(W_{\Delta'}\) of \(W\) generated by
the \(s_\alpha\), for \(\alpha\in\Delta'\). Moreover,
\[
 \Delta\cap R_{\Delta'}=\Delta'.
\]
\end{lemma}

\begin{proof}
Trivial.
\end{proof}

\begin{proposition}\label{sga3:XXI:3.4.8}
Let \(R'\subset R\). The following conditions are equivalent:
\begin{enumerate}
\item[\SGAInlineRefTarget{sga3:XXI:3.4.8:item:i}(i)] There exists a vector subspace \(V'\) of \(V\), or of
  \(\VR(R)\), such that \(R'=R\cap V'\).
\item[\SGAInlineRefTarget{sga3:XXI:3.4.8:item:ii}(ii)] There exist a system of simple roots \(\Delta\) of \(R\)
  and a subset \(\Delta'\) of \(\Delta\) such that
  \(R'=R_{\Delta'}\).
\end{enumerate}
More precisely, under these conditions, every system of simple roots
\(\Delta'\) of \((M,M^\ast,R',R'^\ast)\) is contained in a system of
simple roots \(\Delta\) of \(R\), and one has
\(R'=R_{\Delta'}\).
\end{proposition}

\begin{proof}
The implication \SGARefLink{sga3:XXI:3.4.8:item:ii}{(ii)} \(\Rightarrow\) \SGARefLink{sga3:XXI:3.4.8:item:i}{(i)} is clear. Suppose \SGARefLink{sga3:XXI:3.4.8:item:i}{(i)}.
Then \(R'\) is closed and symmetric, so
\((M,M^\ast,R',R'^\ast)\) is a root datum. Let \(\Delta'\) be a
system of simple roots of this datum, and put
\[
 R'_+=\mathscr{P}(\Delta').
\]
If \(V'=V\), then \(\Delta'\) is a system of simple roots of
\(\cR\), and there is nothing to prove. Otherwise, there exists
\(x\in V^\ast\) such that
\[
 (x,V')=\{0\},
 \qquad
 (x,\alpha)\ne0
 \quad\text{for every }\alpha\in R\setminus R'.
\]
Put
\[
 R_+(x)=\{\alpha\in R\mid(x,\alpha)>0\},
 \qquad
 R_+=R_+(x)\cup R'_+.
\]
For every \(\alpha\in R\), one has
\[
\begin{aligned}
 (x,\alpha)>0 &\Longleftrightarrow \alpha\in R_+(x),\\
 (x,\alpha)<0 &\Longleftrightarrow \alpha\in-R_+(x),\\
 (x,\alpha)=0 &\Longleftrightarrow \alpha\in R'.
\end{aligned}
\]
It follows at once from Corollary~\ref{sga3:XXI:3.3.8} that \(R_+\)
is a system of positive roots of \(R\). Put
\(\Delta=\mathscr{S}(R_+)\). It is enough to prove that
\(\Delta'\subset\Delta\). Otherwise, let
\(\alpha\in\Delta'\setminus\Delta\). By
Corollary~\ref{sga3:XXI:3.2.14}, there exist
\(\beta,\gamma\in R_+\) such that
\[
 \alpha=\beta+\gamma.
\]
If \(\beta,\gamma\in R_+(x)\), then \(\alpha\in R_+(x)\), which is
impossible because \((x,\Delta')=0\). If one of \(\beta,\gamma\),
say \(\beta\), belongs to \(R'_+\), then, since \(R'\) is symmetric
and closed,
\[
 \gamma=\alpha-\beta\in R_+\cap R'=R'_+.
\]
But then \(\alpha\) is not simple in \(R'_+\), again a
contradiction.
\end{proof}

\begin{lemma}\label{sga3:XXI:3.4.9}
Under the preceding conditions, let
\[
 \alpha\in\mathscr{P}(\Delta)\setminus R'.
\]
For every \(w\in W_{\Delta'}\), one has
\[
 w(\alpha)\in\mathscr{P}(\Delta)\setminus R'.
\]
\end{lemma}

\begin{proof}
It suffices to verify the assertion for
\(w=s_\beta\), \(\beta\in\Delta'\). It then follows from
Lemma~\ref{sga3:XXI:3.3.1} and from
\(s_\beta(R')=R'\).
\end{proof}

% END inputs/XXI/tex/components/11_expose_XXI_definition346_through_lemma349_en.tex

% BEGIN inputs/XXI/tex/components/12_expose_XXI_lemma3410_through_proposition354_en.tex
\begin{lemma}\label{sga3:XXI:3.4.10}
Let \(w\in W\). Under the conditions of
Proposition~\ref{sga3:XXI:3.4.8}, the following conditions are
equivalent:
\begin{enumerate}
\item[\SGAInlineRefTarget{sga3:XXI:3.4.10:item:i}(i)] \(w\in W_{\Delta'}\).
\item[\SGAInlineRefTarget{sga3:XXI:3.4.10:item:ii}(ii)] For every \(m\in M\), one has \(w(m)-m\in V'\).
\item[(iii)] For every \(\alpha\in R\), one has
  \(w(\alpha)-\alpha\in V'\).
\end{enumerate}
\end{lemma}

\begin{proof}
Clearly \SGARefLink{sga3:XXI:3.4.10:item:i}{(i)} implies \SGARefLink{sga3:XXI:3.4.10:item:ii}{(ii)}, and \SGARefLink{sga3:XXI:3.4.10:item:ii}{(ii)} implies \SGARefLink{sga3:XXI:2.3.7:corollary:item:iii}{(iii)}. Suppose \SGARefLink{sga3:XXI:2.3.7:corollary:item:iii}{(iii)}.
Write
\[
 w=s_{\alpha_n}\cdots s_{\alpha_1},
 \qquad \alpha_i\in\Delta,
\]
and argue by induction on \(n\) that every \(\alpha_i\) belongs to
\(\Delta'\). We may assume that
\[
 w'=s_{\alpha_{n-1}}\cdots s_{\alpha_1}\in W_{\Delta'}.
\]
Then, for every \(\alpha\in R\),
\[
 w(\alpha)-\alpha
   =w'(\alpha)-\alpha-(\alpha_n^\ast,w'(\alpha))\alpha_n.
\]
Taking \(\alpha=w'^{-1}(\alpha_n)\), one obtains
\(2\alpha_n\in V'\), hence \(\alpha_n\in\Delta'\). Therefore
\[
 w=s_{\alpha_n}w'\in W_{\Delta'}.
\]
\end{proof}

\subsection{Miscellaneous remarks}
\label{sga3:XXI:3.5}

\begin{proposition}\label{sga3:XXI:3.5.1}
Let \(R_+\) be a system of positive roots. Put
\[
 \rho_{R_+}
   =\frac12\sum_{\alpha\in\operatorname{ind}(R_+)}\alpha.
\]
Then
\[
 (\beta^\ast,\rho_{R_+})=1
 \qquad\text{for every }\beta\in\mathscr{S}(R_+).
\]
\end{proposition}

\begin{proof}
One can write
\[
 2\rho_{R_+}
   =\beta+
     \sum_{\substack{\alpha\in\operatorname{ind}(R_+)\\
                     \alpha\ne\beta}}\alpha.
\]
Hence
\[
 s_\beta(2\rho_{R_+})=2\rho_{R_+}-2\beta
\]
by Lemma~\ref{sga3:XXI:3.3.2}.
\end{proof}

\begin{corollary}\label{sga3:XXI:3.5.2}
Put
\[
 \rho_{R_+}^\ast
   =\frac12
     \sum_{\alpha^\ast\in\operatorname{ind}(R_+^\ast)}
     \alpha^\ast.
\]
Then:
\begin{enumerate}
\item[(i)] \((\rho_{R_+}^\ast,\alpha)>0\) for every
  \(\alpha\in R_+\); equivalently,
  \(\rho_{R_+}^\ast\in\mathcal{C}(R_+)\), in the notation of
  Definition~\ref{sga3:XXI:3.6.8}.
\item[(ii)] For every \(\alpha\in R\),
  \[
   \operatorname{ord}_{\mathscr{S}(R_+)}(\alpha)
      =(\rho_{R_+}^\ast,\alpha).
  \]
\end{enumerate}
\footnote{Editors' note: By Lemma~\ref{sga3:XXI:3.1.5}, it is enough
to verify the formula for \(\alpha\in\Delta\). Proposition
\ref{sga3:XXI:3.5.1}, applied to \(R^\ast\), then gives
\((\rho_{R_+}^\ast,\alpha)=1
=\operatorname{ord}_{\mathscr{S}(R_+)}(\alpha)\).}
\end{corollary}

\begin{remark}\label{sga3:XXI:3.5.3}
If \(w\in W\), then
\[
 \rho_{w(R_+)}=w(\rho_{R_+}),
 \qquad
 \rho_{w(R_+)}^\ast=w(\rho_{R_+}^\ast).
\]
\end{remark}

\begin{proposition}\label{sga3:XXI:3.5.4}
Let \(\alpha\) and \(\gamma\) be two nonproportional roots, with
\(\alpha\) indivisible. There exist a system of simple roots
containing \(\alpha\) and a root \(\beta\) such that
\[
 \gamma=a\alpha+b\beta,
 \qquad a,b\in\mathbb{N}.
\]
\end{proposition}

\begin{proof}
Choose a basis of the vector space \(\VR(R)\) containing
\(\alpha_1=\alpha\) and \(\alpha_2=\gamma\), and consider the
lexicographic order with respect to this basis. Let \(R_+\) be the set
of roots \(>0\). It is clear that \(R_+\) is a system of positive
roots and\footnote{Editors' note: The original was modified here to
take account of the case in which \(2\alpha\) is a root.} that the
least element of \(R_+\) not proportional to \(\alpha\) is simple.
This element has the form
\[
 \beta=p\alpha+q\gamma,
 \qquad 0<q\le1.
\]
Consequently
\[
 \gamma=q^{-1}\beta-q^{-1}p\alpha.
\]
Since \(q^{-1}>0\), one has
\[
 q^{-1}\in\mathbb{N}^{\ast},
 \qquad
 -q^{-1}p\in\mathbb{N}.
\]
\end{proof}

% END inputs/XXI/tex/components/12_expose_XXI_lemma3410_through_proposition354_en.tex

% BEGIN inputs/XXI/tex/components/13_expose_XXI_proposition355_through_lemma361_en.tex
We conclude with two remarks on the group \(\Gamma_0(R)\).

\begin{proposition}\label{sga3:XXI:3.5.5}
Let \(G\) be an abelian group, and let \(f:R\to G\) be a map
satisfying:
\begin{enumerate}
\item[(i)] if \(\alpha\in R\), then
  \(f(-\alpha)=-f(\alpha)\);
\item[(ii)] if \(\alpha,\beta,\alpha+\beta\in R\), then
  \[
   f(\alpha+\beta)=f(\alpha)+f(\beta).
  \]
\end{enumerate}
There exists a unique group homomorphism
\[
 \overline f:\Gamma_0(R)\longrightarrow G
\]
such that \(\overline f(\alpha)=f(\alpha)\) for every
\(\alpha\in R\).
\end{proposition}

\begin{proof}
Let \(\Delta\) be a system of simple roots of \(R\), and write
\[
 \beta=\sum_{\alpha\in\Delta}a(\alpha)\alpha.
\]
Lemma~\ref{sga3:XXI:3.2.16} gives at once
\[
 f(\beta)=\sum_{\alpha\in\Delta}a(\alpha)f(\alpha).
\]
Now \(\Delta\) is a basis of \(\Gamma_0(R)\).
\end{proof}

\begin{proposition}\label{sga3:XXI:3.5.6}
Let \(\Delta\) be a system of simple roots. There exists on
\(\Gamma_0(R)\) a totally ordered group structure such that the roots
\(>0\) are the elements of \(\mathscr{P}(\Delta)\), and such that
\[
 \alpha\longmapsto\operatorname{ord}_\Delta(\alpha)
\]
is increasing.
\end{proposition}

\begin{proof}
Let
\[
 \alpha_1,\alpha_2,\ldots,\alpha_n,
 \qquad n=\rgss(\cR),
\]
be the elements of \(\Delta\). Every \(x\in\Gamma_0(R)\) has a
decomposition
\[
 x=\sum_{i=1}^{n}m_i(x)\alpha_i.
\]
Take the lexicographic order with respect to the functions
\[
 \sum_i m_i,\quad m_n,\quad m_{n-1},\quad\ldots,\quad m_2.
\]
\end{proof}

\begin{remark}\label{sga3:XXI:3.5.7}
The first roots, in order, are
\[
 \alpha_1,\alpha_2,\ldots,\alpha_n;
\]
then come, when they are roots,
\[
 2\alpha_1,\quad
 \alpha_1+\alpha_2,\quad
 \alpha_1+\alpha_3,\quad\ldots.
\]
\end{remark}

\subsection{Weyl chambers}
\label{sga3:XXI:3.6}

\begin{lemma}\label{sga3:XXI:3.6.1}
Let \(V\) be a finite-dimensional real vector
space.\footnote{Editors' note: The original \(\mathbb{Q}\) has been
replaced by \(\mathbb{R}\).} Let \(f_i\) be independent linear forms.
Put
\[
 C=\{x\in V\mid f_i(x)>0\}.
\]
Then \(C\) is a maximal convex subset of
\[
 X=V\setminus\bigcup_i f_i^{-1}(0).
\]
\end{lemma}

\begin{proof}
Trivial.
\end{proof}

% END inputs/XXI/tex/components/13_expose_XXI_proposition355_through_lemma361_en.tex

% BEGIN inputs/XXI/tex/components/14_expose_XXI_definition362_through_definition368_en.tex
\begin{definition}\label{sga3:XXI:3.6.2}
A subset \(C\subset V\) described by the construction of
Lemma~\ref{sga3:XXI:3.6.1} is called, here, a \emph{chamber} of \(V\).
\end{definition}

\begin{definition}\label{sga3:XXI:3.6.3}
A hyperplane \(H\) of \(V\) is called a \emph{wall} of \(C\) if
\[
 H\cap(\overline C\setminus C)
\]
contains a nonempty open subset of \(H\).
\end{definition}

\begin{remark}\label{sga3:XXI:3.6.4}
For a convex subset, the closure can be described without reference
to the topology of \(V\): it is the set of endpoints of all open
segments contained in the given subset.
\end{remark}

\begin{lemma}\label{sga3:XXI:3.6.5}
Under the conditions of Lemma~\ref{sga3:XXI:3.6.1},
\[
 \overline C=\{x\in V\mid f_i(x)\ge0\}.
\]
The walls of \(C\) are the hyperplanes \(f_i^{-1}(0)\).
\end{lemma}

\begin{proof}
The first assertion is clear. The second follows from
\[
 \overline C\setminus C\subset\bigcup_i f_i^{-1}(0)
\]
and from the fact that the \(f_i^{-1}(0)\) are visibly walls of
\(C\).
\end{proof}

\begin{proposition}\label{sga3:XXI:3.6.6}
Let \(C\) be a chamber of \(V\). If
\[
 H_i,\qquad i=1,2,\ldots,n,
\]
are the distinct walls of \(C\), then, for every system of linear
forms \(\{u_i\}\) such that \(H_i=u_i^{-1}(0)\), there exist
\(\varepsilon_i\in\{-1,+1\}\) for which
\[
 C=\{x\in V\mid\varepsilon_i u_i(x)>0\}.
\]
For every wall \(H\) of \(C\),
\[
 H\cap C=\varnothing,
 \qquad
 H\cap\overline C=\overline{C_H},
\]
where \(C_H\) is a chamber in \(H\). The walls of \(C_{H_i}\) are
the \(H_j\cap H_i\), for \(j\ne i\).
\end{proposition}

\begin{proof}
This follows immediately from the lemma.
\end{proof}

\begin{definition}\label{sga3:XXI:3.6.7}
The \(C_{H_i}\) are called the \emph{faces} of \(C\).
\end{definition}

Let now
\[
 \cR=(M,M^\ast,R,R^\ast)
\]
be a root datum, and put
\[
 V_{\mathbb{R}}^\ast=M^\ast\otimes\mathbb{R}.
\]

\begin{definition}\label{sga3:XXI:3.6.8}
\footnote{Editors' note: The definition of the hyperplanes
\(\mathcal{H}_\alpha\) was added here.}
For every \(\alpha\in R\), put
\[
 \mathcal{H}_\alpha
   =\{x\in V_{\mathbb{R}}^\ast\mid(x,\alpha)=0\}.
\]
Write
\[
 X=V_{\mathbb{R}}^\ast
   \setminus\bigcup_{\alpha\in R}\mathcal{H}_\alpha.
\]
For every \(x\in X\), put
\[
 R_+(x)=\{\alpha\in R\mid(x,\alpha)>0\}.
\]
For every system of positive roots \(R_+\), put
\[
 \mathcal{C}(R_+)
   =\{x\in V_{\mathbb{R}}^\ast
       \mid (x,\alpha)>0\text{ for every }\alpha\in R_+\}.
\]
\end{definition}

% END inputs/XXI/tex/components/14_expose_XXI_definition362_through_definition368_en.tex

% BEGIN inputs/XXI/tex/components/15_expose_XXI_proposition369_through_remark3614_en.tex
\begin{proposition}\label{sga3:XXI:3.6.9}
\begin{enumerate}
\item[\SGAInlineRefTarget{sga3:XXI:3.6.9:item:i}(i)] For every \(x\in X\), \(R_+(x)\) is a system of positive
roots. For every system of positive roots \(R_+\),
\(\mathcal{C}(R_+)\) is a chamber in \(V_{\mathbb R}^\ast\). The
\(\mathcal{C}(R_+)\) are the maximal convex subsets of \(X\).
\item[\SGAInlineRefTarget{sga3:XXI:3.6.9:item:ii}(ii)] Let \(\Delta\) be a system of simple roots. One has
\[
 \mathcal{C}(\mathscr{P}(\Delta))
 =\{x\in V_{\mathbb R}^\ast
       \mid (x,\alpha)>0\text{ for every }\alpha\in\Delta\}.
\]
The walls of \(\mathcal{C}(\mathscr{P}(\Delta))\) are the
hyperplanes
\[
 \mathcal{H}_\alpha=\alpha^{-1}(0),
 \qquad \alpha\in\Delta;
\]
its faces are
\[
 C_\alpha
 =\{x\in V_{\mathbb R}^\ast
      \mid (x,\alpha)=0,\ 
      (x,\beta)>0\text{ for }\beta\in\Delta,\ \beta\ne\alpha\}.
\]
\item[\SGAInlineRefTarget{sga3:XXI:3.6:item:iii}(iii)] One has the equivalence
\[
 R_+(x)=R_+
 \quad\Longleftrightarrow\quad
 x\in\mathcal{C}(R_+).
\]
\end{enumerate}
\end{proposition}

\begin{proof}
It is first clear that \(R_+(x)\) is a system of positive roots.
Since \(x\in\mathcal{C}(R_+(x))\), the union of the
\(\mathcal{C}(R_+(x))\) is \(X\). Property \SGARefLink{sga3:XXI:3.6:item:iii}{(iii)} is immediate; it
follows that the \(\mathcal{C}(R_+)\) form a partition of \(X\).
The first assertion of \SGARefLink{sga3:XXI:3.6.9:item:ii}{(ii)} is evident. It follows at once that
\(\mathcal{C}(R_+)\) is a chamber in \(V_{\mathbb R}^\ast\), which
proves the remainder of \SGARefLink{sga3:XXI:3.6.9:item:i}{(i)}. Put
\[
 C=\mathcal{C}(\mathscr{P}(\Delta)).
\]
It remains only to note that
\[
 \overline C\cap
 \left(V_{\mathbb R}^\ast\setminus
       \bigcup_{\alpha\in R}\alpha^{-1}(0)\right)
 =
 \overline C\cap
 \left(V_{\mathbb R}^\ast\setminus
       \bigcup_{\alpha\in\Delta}\alpha^{-1}(0)\right)
\]
to complete the proof of \SGARefLink{sga3:XXI:3.6.9:item:ii}{(ii)} by
Lemma~\ref{sga3:XXI:3.6.1}.
\end{proof}

\begin{definition}\label{sga3:XXI:3.6.10}
The \(\mathcal{C}(R_+)\) are called the \emph{Weyl chambers} of the
root datum. Conversely, for every Weyl chamber \(C\), put
\[
 R_+(C)=R_+(x)
\]
for any \(x\in C\).
\end{definition}

\begin{corollary}\label{sga3:XXI:3.6.11}
The maps
\[
 R_+\longmapsto\mathcal{C}(R_+),
 \qquad
 C\longmapsto R_+(C)
\]
give a bijective correspondence between systems of positive roots
and Weyl chambers.
\end{corollary}

This correspondence is invariant under the Weyl group:

\begin{lemma}\label{sga3:XXI:3.6.12}
If \(w\in W(\mathcal{R})\), then
\[
 \mathcal{C}(w(R_+))=w(\mathcal{C}(R_+)).
\]
\end{lemma}

\begin{corollary}\label{sga3:XXI:3.6.13}
The correspondences
\[
 \Delta\longleftrightarrow R_+\longleftrightarrow C
\]
are isomorphisms of homogeneous spaces under \(W(\mathcal{R})\).
\end{corollary}

We shall see later that these homogeneous spaces are principal
(\SGARefLink{sga3:XXI:5.5}{5.5}).

\begin{remark}\label{sga3:XXI:3.6.14}
If \(C\) is a Weyl chamber, then \(-C\) is also one; it is called the
\emph{opposite} of \(C\). Consequently there exists
\(w_0\in W\)---and in fact exactly one, cf.
\SGARefLink{sga3:XXI:5.5}{5.5}---such that
\[
 w_0(C)=-C.
\]
It is called the \emph{symmetry} of the root datum relative to the
Weyl chamber \(C\) (or relative to \(R_+(C)\), or to
\(\mathscr{S}(R_+(C))\), and so on).
\end{remark}

% END inputs/XXI/tex/components/15_expose_XXI_proposition369_through_remark3614_en.tex

% BEGIN inputs/XXI/tex/components/16_expose_XXI_section4_rank_two_tables_en.tex
\section{Reduced root data of semisimple rank \(2\)}
\label{sga3:XXI:4}

\renewcommand{\thetheorem}{\thesection.\arabic{theorem}}
\setcounter{theorem}{0}

\SGARefTarget{sga3:XXI:4.0}
\noindent\textbf{4.0.}\footnote{Editors' note: The numbering \(4.0\)
was added for later references.}
Let \(\mathcal{R}\) be a root datum of semisimple rank \(2\). Let
\(\{\alpha,\beta\}\) be a system of simple roots. Suppose that
\(\ell(\alpha)\le\ell(\beta)\). By
\ref{sga3:XXI:2.3.1} and \ref{sga3:XXI:3.2.1}, there are then four
possibilities:
\[
\begin{array}{c|c|c|c|c}
\text{Type}
 &\ell(\beta)/\ell(\alpha)
 &\ell(\beta^\ast)/\ell(\alpha^\ast)
 &(\beta^\ast,\alpha)
 &(\alpha^\ast,\beta)\\ \hline
A_1\times A_1&-&-&0&0\\
A_2&1&1&-1&-1\\
B_2&2&1/2&-1&-2\\
G_2&3&1/3&-1&-3
\end{array}
\]
It follows from Remark~\ref{sga3:XXI:3.4.5} that the order of
\(s_\alpha s_\beta\) is respectively \(2,3,4,6\).

We study each of these systems separately and give the list of
indivisible roots.

\medskip
\noindent\emph{Type \(A_1\times A_1\).}
The indivisible roots are
\[
 \alpha,\quad\beta,\quad-\alpha,\quad-\beta.
\]
The corresponding coroots are
\[
 \alpha^\ast,\quad\beta^\ast,\quad-\alpha^\ast,\quad-\beta^\ast.
\]

\medskip
\noindent\emph{Type \(A_2\).}
The indivisible positive roots are as follows:
\[
\begin{array}{c|ccc}
\text{root }\gamma
 &\alpha&\beta&\alpha+\beta\\ \hline
\ell(\gamma)/\ell(\alpha)
 &1&1&1\\ \hline
\text{coroot }\gamma^\ast
 &\alpha^\ast&\beta^\ast&\alpha^\ast+\beta^\ast\\ \hline
\ell(\gamma^\ast)/\ell(\beta^\ast)
 &1&1&1
\end{array}
\]
The half-sum of the indivisible positive roots is
\[
 \rho=\alpha+\beta.
\]

\medskip
\noindent\emph{Type \(B_2\).}
The indivisible positive roots are as follows:
\[
\begin{array}{c|cccc}
\text{root }\gamma
 &\alpha&\beta&\alpha+\beta&2\alpha+\beta\\ \hline
\ell(\gamma)/\ell(\alpha)
 &1&2&1&2\\ \hline
\text{coroot }\gamma^\ast
 &\alpha^\ast&\beta^\ast
 &2\alpha^\ast+\beta^\ast&\alpha^\ast+\beta^\ast\\ \hline
\ell(\gamma^\ast)/\ell(\beta^\ast)
 &2&1&2&1
\end{array}
\]
The half-sum of the indivisible positive roots is
\[
 \rho=\frac{4\alpha+3\beta}{2}.
\]

\medskip
\noindent\emph{Type \(G_2\).}
The indivisible positive roots are as follows:
\begin{center}
\resizebox{\textwidth}{!}{%
$\displaystyle
\begin{array}{c|cccccc}
\text{root }\gamma
 &\alpha&\beta&\alpha+\beta&2\alpha+\beta
 &3\alpha+\beta&3\alpha+2\beta\\ \hline
\ell(\gamma)/\ell(\alpha)
 &1&3&1&1&3&3\\ \hline
\text{coroot }\gamma^\ast
 &\alpha^\ast&\beta^\ast&\alpha^\ast+3\beta^\ast
 &2\alpha^\ast+3\beta^\ast&\alpha^\ast+\beta^\ast
 &\alpha^\ast+2\beta^\ast\\ \hline
\ell(\gamma^\ast)/\ell(\beta^\ast)
 &3&1&3&3&1&1
\end{array}
$%
}
\end{center}
The half-sum of the indivisible positive roots is
\[
 \rho=5\alpha+3\beta.
\]

% END inputs/XXI/tex/components/16_expose_XXI_section4_rank_two_tables_en.tex

% BEGIN inputs/XXI/tex/components/17_expose_XXI_proposition41_through_lemma42_en.tex
\begin{proposition}\label{sga3:XXI:4.1}
\footnote{The following numbers \(4.1\) to \(4.4\) are used in the
proof of Theorem~\ref{sga3:XXI:5.1}. There are now simpler proofs of
that theorem; see \emph{Groupes et algèbres de Lie}, Chapter~V,
\S3, Theorem~1. Editors' note: The reference was made explicit and
this note, which occurred at \(4.2\) in the original, was moved here.}
Let \(n\) be the order of \(s_\alpha s_\beta\). Put \(u_0=0\), and,
for \(p\in\mathbb N\), define
\[
\left\{
\begin{aligned}
u_{2p+1}&=u_{2p}+(s_\alpha s_\beta)^p(\alpha),\\
u_{2p+2}&=u_{2p+1}+(s_\alpha s_\beta)^p s_\alpha(\beta).
\end{aligned}
\right.
\]
Then:
\footnote{Editors' note: Let \(\rho\) be the half-sum of the positive
roots, cf. Proposition~\ref{sga3:XXI:3.5.1}. If, as below, one puts
\(w_0=1\), \(w_1=s_\alpha\), \(w_2=s_\beta s_\alpha\),
\(w_3=s_\alpha s_\beta s_\alpha\), and so on, then
\(u_k=\rho-w_k^{-1}(\rho)\). This proves \SGARefLink{sga3:XXI:4.1:item:i}{(i)}, since \(w_{2n}=1\).}
\begin{enumerate}
\item[\SGAInlineRefTarget{sga3:XXI:4.1:item:i}(i)] \(u_{k+2n}=u_k\), for every \(k\in\mathbb N\).
\item[\SGAInlineRefTarget{sga3:XXI:4.1:item:ii}(ii)] \(u_0=0\), \(u_1=\alpha\), \(u_{2n-1}=\beta\),
and \(u_{2n}=0\).
\item[(iii)] If \(1<k<2n-1\), then
\[
 u_k=a_k\alpha+b_k\beta,
 \qquad a_k,b_k\in\mathbb N^\ast.
\]
\end{enumerate}
\end{proposition}

\begin{proof}
Assertion \SGARefLink{sga3:XXI:4.1:item:i}{(i)} follows from \((s_\alpha s_\beta)^n=1\) and
\(u_{2n}=0\).

We prove \SGARefLink{sga3:XXI:4.1:item:ii}{(ii)} and \SGARefLink{sga3:XXI:3.6:item:iii}{(iii)}. Direct calculation gives, in the four cases,
the following sequences of values:
\footnote{Editors' note: The original was corrected here to agree
with the convention \(\ell(\alpha)\le\ell(\beta)\) of \(4.0\).}
\[
\begin{array}{c|l}
A_1\times A_1&0,\alpha,\beta+\alpha,\beta,0,\\[2pt]
A_2&0,\alpha,\beta+2\alpha,2\beta+2\alpha,
       2\beta+\alpha,\beta,0,\\[2pt]
B_2&0,\alpha,\beta+3\alpha,2\beta+4\alpha,
       3\beta+4\alpha,3\beta+3\alpha,2\beta+\alpha,\beta,0,\\[2pt]
G_2&0,\alpha,\beta+4\alpha,2\beta+6\alpha,
       4\beta+9\alpha,5\beta+10\alpha,6\beta+10\alpha,\\
   &\quad 6\beta+9\alpha,5\beta+6\alpha,4\beta+4\alpha,
       2\beta+\alpha,\beta,0.
\end{array}
\]
\end{proof}

\begin{lemma}\label{sga3:XXI:4.2}
Put
\[
 w_{2p}=(s_\beta s_\alpha)^p,
 \qquad
 w_{2p+1}=s_\alpha(s_\beta s_\alpha)^p,
\]
so that \(w_0,\ldots,w_{2n-1}\) are the distinct elements of \(W\).
Let \(u_0,\ldots,u_{2n-1}\) be the elements of \(V\) defined in
Proposition~\ref{sga3:XXI:4.1}. For every \(x\in V^\ast\), one has
\[
 x-w_k(x)=n_k\alpha^\ast+m_k\beta^\ast,
\]
where \(n_k,m_k\in\mathbb Q\) and
\[
 n_k+m_k=(x,u_k).
\]
\footnote{Editors' note: In view of the preceding editor's note, this
follows immediately from Propositions~\ref{sga3:XXI:3.5.1} and
\ref{sga3:XXI:1.1.9}: one has
\[
 n_k+m_k=(x-w_k(x),\rho)
 =(x,\rho-w_k^{-1}(\rho))=(x,u_k).
\]}
\end{lemma}

\begin{proof}
The proof is by induction on \(k\). If \(k=0\), the formula is
trivial. We carry out, for example, the passage from \(w_{2p}\) to
\(w_{2p+1}\). Since \(w_{2p+1}=s_\alpha w_{2p}\),
\[
\begin{aligned}
x-w_{2p+1}(x)
 &=x-w_{2p}(x)+w_{2p}(x)-s_\alpha w_{2p}(x)\\
 &=n_{2p}\alpha^\ast+m_{2p}\beta^\ast
   +(w_{2p}(x),\alpha)\alpha^\ast.
\end{aligned}
\]
Consequently
\[
\begin{aligned}
n_{2p+1}+m_{2p+1}
 &=n_{2p}+m_{2p}
   +\bigl((s_\beta s_\alpha)^p(x),\alpha\bigr)\\
 &=(x,u_{2p})+\bigl(x,(s_\alpha s_\beta)^p(\alpha)\bigr)\\
 &=(x,u_{2p+1}).
\end{aligned}
\]
\end{proof}

% END inputs/XXI/tex/components/17_expose_XXI_proposition41_through_lemma42_en.tex

% BEGIN inputs/XXI/tex/components/18_expose_XXI_corollary43_through_lemma52_statement_en.tex
\begin{corollary}\label{sga3:XXI:4.3}
\footnote{Editors' note: The statement was corrected.}
Let \(x\in V^\ast\). For every \(w\in W\), put
\[
 x-w(x)=a_w\alpha^\ast+b_w\beta^\ast.
\]
If \((x,\alpha)\ge0\) and \((x,\beta)\ge0\), then
\[
 a_w+b_w\ge0.
\]
If moreover \((x,\alpha)>0\) (respectively \((x,\beta)>0\)), then
\[
 a_w+b_w>0
\]
for \(w\ne1,s_\beta\) (respectively for \(w\ne1,s_\alpha\)).
\end{corollary}

\begin{proof}
This follows at once from Proposition~\ref{sga3:XXI:4.1} and
Lemma~\ref{sga3:XXI:4.2}.
\end{proof}

\begin{corollary}\label{sga3:XXI:4.4}
Let \(\mathcal{R}\) be any root datum, and let \(\Delta\) be a system
of simple roots. Let \(\gamma\) be a positive root, let
\(\alpha,\beta\) be two simple roots, and let \(W_{\alpha,\beta}\) be
the subgroup of \(W\) generated by \(s_\alpha\) and \(s_\beta\). If
\[
 \operatorname{ord}_\Delta(s_\alpha(\gamma))
 <\operatorname{ord}_\Delta(\gamma),
 \qquad
 \operatorname{ord}_\Delta(s_\beta(\gamma))
 \le\operatorname{ord}_\Delta(\gamma),
\]
then, for every \(w\in W_{\alpha,\beta}\),
\[
 \operatorname{ord}_\Delta(w(\gamma))
 \le\operatorname{ord}_\Delta(\gamma).
\]
Moreover,
\[
 \operatorname{ord}_\Delta(w(\gamma))
 <\operatorname{ord}_\Delta(\gamma)
\]
if \(w\ne1,s_\beta\).
\end{corollary}

\begin{proof}
Consider the dual root datum \(\mathcal{R}^\ast\), and then the datum
\[
 (M^\ast,M,R'^\ast,R'),
\]
where \(R'\) is the set of roots that are rational linear
combinations of \(\alpha\) and \(\beta\). Applying
Corollary~\ref{sga3:XXI:4.3} to this datum gives the stated result.
\footnote{Editors' note: The hypotheses say that
\((\gamma,\alpha^\ast)>0\) and \((\gamma,\beta^\ast)\ge0\).
Consequently
\(\gamma-w(\gamma)\in\mathbb N\alpha+\mathbb N\beta\) for every
\(w\in W_{\alpha,\beta}\), and it is nonzero if
\(w\notin\{1,s_\beta\}\).}
\end{proof}

\section{The Weyl group: generators and relations}
\label{sga3:XXI:5}
\setcounter{theorem}{0}

Let \(\mathcal{R}\) be a root datum. Since the Weyl group is the same
for this datum and for the corresponding reduced datum, we may assume
\(\mathcal{R}\) reduced when studying the Weyl group.

Let
\[
 \Delta=\{\alpha_1,\ldots,\alpha_n\}
\]
be a system of simple roots, where \(n=\rgss(\mathcal{R})\). Let
\(n_{ij}\) be the order of \(s_{\alpha_i}s_{\alpha_j}\) in \(W\).
In particular \(n_{ii}=1\), and Remarks~\ref{sga3:XXI:3.4.2} and
\ref{sga3:XXI:3.4.3} show that the subgroup \(W_{ij}\) of \(W\)
generated by \(s_{\alpha_i}\) and \(s_{\alpha_j}\) is defined by
\[
 s_{\alpha_i}^{\,2}=s_{\alpha_j}^{\,2}
 =(s_{\alpha_i}s_{\alpha_j})^{n_{ij}}=1.
\]

\begin{theorem}\label{sga3:XXI:5.1}
The group \(W\) is the group generated by the elements
\[
 s_{\alpha_i},\qquad i=1,2,\ldots,n,
\]
subject to the relations
\[
 (s_{\alpha_i}s_{\alpha_j})^{n_{ij}}=1.
\]
\end{theorem}

\begin{proof}
We have already seen that the theorem is true when \(n=2\); we shall
use this fact in the proof. Introduce the group \(\widetilde W\)
generated by elements
\[
 T_i,\qquad i=1,2,\ldots,n,
\]
subject to the relations
\[
 (T_iT_j)^{n_{ij}}=1.
\]
In particular \(n_{ii}=1\), whence \(T_i^2=1\). Let
\[
 p:\widetilde W\longrightarrow W
\]
be the group homomorphism sending \(T_i\) to \(s_{\alpha_i}\). We know
that \(p\) is surjective; we shall prove that it is injective.
\end{proof}

\begin{lemma}\label{sga3:XXI:5.2}
For every \(\alpha\in\mathscr{P}(\Delta)\), one can define uniquely an
element \(T_\alpha\in\widetilde W\) with the following properties:
\begin{enumerate}
\item[(i)] \(p(T_\alpha)=s_\alpha\);
\item[(ii)] \(T_{\alpha_i}=T_i\);
\item[(iii)] if \(\beta\) and \(\alpha\) are positive roots such that
\(s_{\alpha_i}(\alpha)=\beta\), then
\[
 T_iT_\alpha T_i=T_\beta.
\]
\end{enumerate}
\end{lemma}

% END inputs/XXI/tex/components/18_expose_XXI_corollary43_through_lemma52_statement_en.tex

% BEGIN inputs/XXI/tex/components/19_expose_XXI_lemma52_proof_partA_en.tex
\begin{proof}
First observe that it follows from
\ref{sga3:XXI:1.2.10} and \ref{sga3:XXI:3.3.6} that (i) is a
consequence of (ii) and \SGARefLink{sga3:XXI:3.6:item:iii}{(iii)}, and that (ii) and \SGARefLink{sga3:XXI:3.6:item:iii}{(iii)} determine the
\(T_\alpha\) completely. We construct them by induction on
\(\operatorname{ord}_\Delta(\alpha)\). If
\(\operatorname{ord}_\Delta(\alpha)=1\), then \(\alpha\in\Delta\),
and we put \(T_\alpha=T_i\) if \(\alpha=\alpha_i\).

Consider the hypothesis
\[
\begin{split}
(\mathrm{H}_p)\quad&
\text{there exist \(T_\alpha\), for
\(\alpha\in\mathscr{P}(\Delta)\) with
\(\operatorname{ord}_\Delta(\alpha)\le p\), satisfying (ii),}\\
&\text{and satisfying (iii) whenever
\(\operatorname{ord}_\Delta(\alpha)\le p\) and
\(\operatorname{ord}_\Delta(\beta)\le p\).}
\end{split}
\]
It holds for \(p=1\). Indeed, if \(\alpha\) and
\(s_{\alpha_i}(\alpha)=\beta\) are simple, then \(\alpha_i\) and
\(\alpha\) are orthogonal. Thus, if
\(\alpha_j=\alpha=\beta\), one has \(n_{ij}=2\), whence
\[
 T_iT_jT_i=T_j.
\]

Suppose \(p>1\) and \((\mathrm{H}_{p-1})\) holds.

\medskip
\noindent\textbf{A) Construction of the \(T_\alpha\) for
\(\operatorname{ord}_\Delta(\alpha)\le p\).}
It is enough to do this when
\(\operatorname{ord}_\Delta(\alpha)=p\). There exists
\(\alpha_i\in\Delta\) such that
\[
 s_{\alpha_i}(\alpha)\in\mathscr{P}(\Delta),
 \qquad
 \operatorname{ord}_\Delta(s_{\alpha_i}(\alpha))<p
\]
by Lemma~\ref{sga3:XXI:3.3.3}. Put
\begin{equation}
\tag{\(\star\)}
\label{sga3:XXI:5.2:construction}
 T_\alpha=T_iT_{s_{\alpha_i}(\alpha)}T_i.
\end{equation}
We verify that \(T_\alpha\) depends only on \(\alpha\). Let
\(\alpha_j\in\Delta\) also satisfy
\[
 s_{\alpha_j}(\alpha)\in\mathscr{P}(\Delta),
 \qquad
 \operatorname{ord}_\Delta(s_{\alpha_j}(\alpha))<p.
\]
We must prove
\begin{equation}
\tag{\(+\)}
\label{sga3:XXI:5.2:independence}
 T_iT_{s_{\alpha_i}(\alpha)}T_i
 =T_jT_{s_{\alpha_j}(\alpha)}T_j.
\end{equation}
We distinguish two cases.

\begin{enumerate}
\item[(1)] Suppose that \(\alpha\) is a linear combination of
\(\alpha_i\) and \(\alpha_j\). The same is then true of
\(s_{\alpha_i}(\alpha)\) and \(s_{\alpha_j}(\alpha)\); by
\((\mathrm{H}_{p-1})\), the elements
\(T_{s_{\alpha_i}(\alpha)}\) and
\(T_{s_{\alpha_j}(\alpha)}\) are words in \(T_i,T_j\). Since the
projection of \eqref{sga3:XXI:5.2:independence} to \(W\) holds and
the theorem is true for \(n=2\), the restriction of \(p\) to the
subgroup of \(\widetilde W\) generated by \(T_i,T_j\) is injective.
Thus \eqref{sga3:XXI:5.2:independence} holds.

\item[(2)] Suppose that \(\alpha\) is not a linear combination of
\(\alpha_i\) and \(\alpha_j\). Then, for
\(w\in W_{\alpha_i,\alpha_j}\), all the \(w(\alpha)\) are positive
(cf. Lemma~\ref{sga3:XXI:3.4.9}). The relation to be verified can
also be written
\begin{equation}
\tag{\({+}{+}\)}
\label{sga3:XXI:5.2:rank-two-relation}
 (T_iT_j)^{n_{ij}-1}T_{s_{\alpha_i}(\alpha)}
 (T_jT_i)^{n_{ij}-1}
 =T_{s_{\alpha_j}(\alpha)}.
\end{equation}
Corollary~\ref{sga3:XXI:4.4} shows that all the \(w(\alpha)\) have
order \(<p\) for
\(w\in W_{\alpha_i,\alpha_j}\), \(w\ne1\). We may therefore apply
\((\mathrm{H}_{p-1})\) \(2(n_{ij}-1)\) times, which proves the
relation.
\end{enumerate}

\medskip
\noindent\textbf{B) Verification of \((\mathrm{H}_p)\).}
\footnote{Editors' note: The argument in the original has been
expanded in what follows.}
We must verify that, if \(\alpha_j\in\Delta\) and
\(\beta=s_{\alpha_j}(\alpha)\) satisfies
\(\operatorname{ord}_\Delta(\beta)\le p\), then
\[
 T_jT_\alpha T_j=T_\beta.
\]
If \(\operatorname{ord}_\Delta(\beta)<p\), this follows from the
preceding construction, since \(T_j^2=1\). We may therefore suppose
\[
 \operatorname{ord}_\Delta(\beta)
 =p=\operatorname{ord}_\Delta(\alpha).
\]
In this case \(\alpha\) and \(\alpha_j\) are orthogonal, so
\(\beta=\alpha\), and it remains to prove
\begin{equation}
\tag{\(\dagger\)}
\label{sga3:XXI:5.2:orthogonal}
 T_jT_\alpha T_j=T_\alpha.
\end{equation}
By \eqref{sga3:XXI:5.2:construction},
\[
 T_\alpha=T_iT_{s_{\alpha_i}(\alpha)}T_i,
\]
so it remains to verify
\begin{equation}
\tag{\({+}{+}{+}\)}
\label{sga3:XXI:5.2:orthogonal-expanded}
 T_jT_iT_{s_{\alpha_i}(\alpha)}T_iT_j
 =T_\alpha
 =T_iT_{s_{\alpha_i}(\alpha)}T_i.
\end{equation}
\end{proof}

% END inputs/XXI/tex/components/19_expose_XXI_lemma52_proof_partA_en.tex

% BEGIN inputs/XXI/tex/components/20_expose_XXI_lemma52_partB_through_theorem56_statement_en.tex
\begin{proof}[Proof of Lemma~\ref{sga3:XXI:5.2}, continued]
Put \(m=n_{ij}\), \(s_i=s_{\alpha_i}\), and
\(s_j=s_{\alpha_j}\). We have
\[
 T_jT_i=(T_iT_j)^{m-1}.
\]
By Corollary~\ref{sga3:XXI:4.4},
\(\operatorname{ord}_\Delta(w(\alpha))<p\) for every
\(w\in W_{ij}\) distinct from
\[
 1=(s_is_j)^m
 \quad\text{and}\quad
 s_j=s_i(s_is_j)^{m-1}.
\]
The induction hypothesis therefore gives
\[
\begin{split}
 &T_j(T_iT_j)^{m-2}T_{s_i(\alpha)}
       (T_jT_i)^{m-2}T_j\\
 &\qquad
 =T_{s_j(s_is_j)^{m-2}s_i(\alpha)}
 =T_{s_is_j(\alpha)}
 =T_{s_i(\alpha)}.
\end{split}
\]
The last equality holds because \(s_j(\alpha)=\alpha\). Hence
\[
 (T_iT_j)^{m-1}T_{s_i(\alpha)}(T_jT_i)^{m-1}
 =T_iT_{s_i(\alpha)}T_i,
\]
which proves \eqref{sga3:XXI:5.2:orthogonal-expanded}.
\end{proof}

\begin{lemma}\label{sga3:XXI:5.3}
Let \(h\in\widetilde W\). Write
\[
 h=T_{\alpha_1}\cdots T_{\alpha_m},
\]
where the \(\alpha_i\in\Delta\), not necessarily distinct, are chosen
so that \(m\) is minimal. Then
\[
 p(h)(\alpha_m)\in-\mathscr{P}(\Delta).
\]
\end{lemma}

\begin{proof}
Since
\[
 p(T_{\alpha_m})(\alpha_m)
 =s_{\alpha_m}(\alpha_m)=-\alpha_m,
\]
if \(p(h)(\alpha_m)\) were positive, there would be an index
\(k\), \(1\le k\le m-1\), such that
\[
 u=s_{\alpha_{k+1}}\cdots s_{\alpha_m}(\alpha_m)
 =-s_{\alpha_{k+1}}\cdots s_{\alpha_{m-1}}(\alpha_m)
 \in-\mathscr{P}(\Delta)
\]
and \(s_{\alpha_k}(u)\in\mathscr{P}(\Delta)\). Necessarily
\(u=-\alpha_k\), by Lemma~\ref{sga3:XXI:3.3.1}, whence
\[
 s_{\alpha_{k+1}}\cdots s_{\alpha_{m-1}}(\alpha_m)=\alpha_k.
\]
Property \SGARefLink{sga3:XXI:3.6:item:iii}{(iii)} of Lemma~\ref{sga3:XXI:5.2} then gives
\[
 T_{\alpha_k}T_{\alpha_{k+1}}\cdots
 T_{\alpha_{m-1}}T_{\alpha_m}
 =
 T_{\alpha_{k+1}}\cdots T_{\alpha_{m-1}},
\]
contradicting the minimality of \(m\).
\end{proof}

Let now \(h\in\widetilde W\) satisfy
\[
 p(h)\bigl(\mathscr{P}(\Delta)\bigr)
 \subset\mathscr{P}(\Delta).
\]
Lemma~\ref{sga3:XXI:5.3} gives \(p(h)=1\). This proves
Theorem~\ref{sga3:XXI:5.1} and also the following two corollaries.

\begin{corollary}\label{sga3:XXI:5.4}
If \(R_+\) is a system of positive roots and \(w\in W\) satisfies
\(w(R_+)=R_+\), then \(w=1\).
\end{corollary}

\begin{corollary}\label{sga3:XXI:5.5}
The Weyl group acts freely and transitively on the set of systems of
positive roots, respectively on the set of systems of simple roots,
respectively on the set of Weyl chambers.
\end{corollary}

Choose a system of simple roots \(\Delta\), and put
\[
 R^+=\mathscr{P}(\Delta).
\]
\footnote{Editors' note: The superscript notation \(R^+\), rather
than \(R_+\), is used here so that
\(R^+_{\alpha,\beta}=R_{\alpha,\beta}\cap R^+\).}
For every pair of simple roots
\((\alpha,\beta)\in\Delta\times\Delta\), let
\(R_{\alpha,\beta}\) be the set of roots that are linear combinations
of \(\alpha\) and \(\beta\). Put
\[
 R^+_{\alpha,\beta}=R^+\cap R_{\alpha,\beta},
\]
and let \(W_{\alpha,\beta}\) be the Weyl group of
\(R_{\alpha,\beta}\), that is, the subgroup of \(W\) generated by
\(s_\alpha\) and \(s_\beta\).

\begin{theorem}[Tits]\label{sga3:XXI:5.6}
Let \(\alpha\) and \(\beta\) be two simple roots, and let \(w\in W\)
satisfy \(w(\alpha)=\beta\). There exist a sequence of simple roots
\(\alpha_0,\ldots,\alpha_m\) and a sequence
\(w_0,\ldots,w_{m-1}\) of elements of \(W\) satisfying:
\begin{enumerate}
\item[(i)] \(\alpha_0=\alpha\) and \(\alpha_m=\beta\);
\item[(ii)] \(w=w_{m-1}w_{m-2}\cdots w_0\);
\item[(iii)] \(w_i(\alpha_i)=\alpha_{i+1}\) for
  \(0\le i\le m-1\);
\item[(iv)] for every \(i\), \(0\le i\le m-1\), such that
  \(\alpha_i\ne\alpha_{i+1}\), one has
  \(w_i\in W_{\alpha_i,\alpha_{i+1}}\);
\item[(v)] for every \(i\), \(0\le i\le m-1\), such that
  \(\alpha_i=\alpha_{i+1}\), there exists a simple root \(\beta_i\)
  such that \(w_i\in W_{\alpha_i,\beta_i}\).
\end{enumerate}
\end{theorem}

% END inputs/XXI/tex/components/20_expose_XXI_lemma52_partB_through_theorem56_statement_en.tex

% BEGIN inputs/XXI/tex/components/21_expose_XXI_theorem56_proof_en.tex
\begin{proof}
Put\footnote{Editors' note: In what follows, the original argument has
been expanded, and equality \((1)\) has been corrected.}
\[
 M(w)
 =\operatorname{Card}\bigl(R^+\cap w^{-1}(-R^+)\bigr)
 =\operatorname{Card}\{\gamma\in R^+\mid w(\gamma)\in-R^+\}.
\]
If \(M(w)=0\), then \(w(R^+)=R^+\), so
\(w=1\) by Corollary~\ref{sga3:XXI:5.4}; the theorem is then
trivial: \(m=0\), and assertions (iii)--(v) are empty.

We argue by induction on \(M(w)\). Suppose \(M(w)>0\). There exists
\(\beta_0\in\Delta\) such that \(w(\beta_0)\in-R^+\). Put
\(\alpha_0=\alpha\), and consider
\[
 A=w^{-1}(R^+)\cap R_{\alpha_0,\beta_0}.
\]
This is a system of positive roots of
\(R_{\alpha_0,\beta_0}\). Hence there exists
\(w_0\in W_{\alpha_0,\beta_0}\) such that
\[
 w_0^{-1}\bigl(R^+_{\alpha_0,\beta_0}\bigr)=A.
\]
Put \(w'=ww_0^{-1}\). Lemma~\ref{sga3:XXI:3.4.9} gives
\[
 R^+-R^+_{\alpha_0,\beta_0}
 =
 w_0\bigl(R^+-R^+_{\alpha_0,\beta_0}\bigr),
\]
and therefore
\begin{equation}
\tag{1}
\label{sga3:XXI:5.6:outside-rank-two}
\begin{split}
 &(R^+-R^+_{\alpha_0,\beta_0})\cap w^{-1}(-R^+)\\
 &\qquad=
 w_0\Bigl((R^+-R^+_{\alpha_0,\beta_0})
             \cap w'^{-1}(-R^+)\Bigr).
\end{split}
\end{equation}
On the other hand,
\begin{equation}
\tag{2}
\label{sga3:XXI:5.6:beta-negative}
 \beta_0\in
 R^+_{\alpha_0,\beta_0}\cap w^{-1}(-R^+).
\end{equation}
Since
\(w_0(R_{\alpha_0,\beta_0})=R_{\alpha_0,\beta_0}\), one has
\[
 w_0(-A)
 =R_{\alpha_0,\beta_0}\cap w'^{-1}(-R^+),
\]
and hence
\begin{equation}
\tag{\(2'\)}
\label{sga3:XXI:5.6:rank-two-empty}
\begin{split}
 R^+_{\alpha_0,\beta_0}\cap w'^{-1}(-R^+)
 &=
 R^+_{\alpha_0,\beta_0}\cap w_0(-A)\\
 &=
 R^+_{\alpha_0,\beta_0}\cap(-R^+_{\alpha_0,\beta_0})
 =\varnothing.
\end{split}
\end{equation}
It follows from \eqref{sga3:XXI:5.6:outside-rank-two},
\eqref{sga3:XXI:5.6:beta-negative}, and
\eqref{sga3:XXI:5.6:rank-two-empty} that
\[
 M(w')<M(w).
\]

Put \(\alpha_1=w_0(\alpha_0)\). We show that
\(\alpha_1\in\Delta\), or equivalently that
\(\alpha_0\in w_0^{-1}(\Delta)\). Since
\(w(\alpha_0)\in\Delta\), we have
\[
 \alpha_0\in
 w^{-1}(\Delta)\cap R_{\alpha_0,\beta_0}.
\]
Thus \(\alpha_0\) is a simple root of
\[
 A
 =w^{-1}(R^+)\cap R_{\alpha_0,\beta_0}
 =w_0^{-1}(R^+_{\alpha_0,\beta_0}),
\]
and consequently belongs to
\[
 w_0^{-1}\bigl(\Delta\cap R^+_{\alpha_0,\beta_0}\bigr)
 =
 w_0^{-1}\{\alpha_0,\beta_0\}
\]
by Lemma~\ref{sga3:XXI:3.4.8}. Hence
\(\alpha_1=w_0(\alpha_0)\) equals \(\alpha_0\) or \(\beta_0\).
If \(\alpha_1\ne\alpha_0\), then
\(\alpha_1=\beta_0\) and
\(w_0\in W_{\alpha_0,\alpha_1}\).

Finally,
\[
 \beta=w'(\alpha_1)
 \qquad\text{and}\qquad
 M(w')<M(w),
\]
so the induction hypothesis completes the proof.
\end{proof}

% END inputs/XXI/tex/components/21_expose_XXI_theorem56_proof_en.tex

% BEGIN inputs/XXI/tex/components/22_expose_XXI_section61_morphisms_en.tex
\section{Morphisms of root data}\label{sga3:XXI:6}
\renewcommand{\thetheorem}{\thesubsection.\arabic{theorem}}

\subsection{Definition}\label{sga3:XXI:6.1}

Let
\[
 \cR=(M,M^\ast,R,R^\ast),
 \qquad
 \cR'=(M',M'^\ast,R',R'^\ast)
\]
be two root data. Let \(f:M'\to M\) be a linear map and
\({}^tf:M^\ast\to M'^\ast\) its transpose.

\begin{definition}\label{sga3:XXI:6.1.1}
The map \(f\) is called a \emph{morphism} from \(\cR'\) to \(\cR\),
and is denoted
\[
 f:\cR'\longrightarrow\cR,
\]
if \(f\) induces a bijection from \(R'\) onto \(R\), and
\({}^tf\) induces a bijection from \(R^\ast\) onto \(R'^\ast\).
\end{definition}

Then \({}^tf\) is a morphism of the dual root data:
\[
 {}^tf:\cR^\ast\longrightarrow\cR'^\ast.
\]
If \(f:\cR'\to\cR\) is a morphism and
\(\alpha=f(\alpha')\) for \(\alpha'\in R'\), then
\[
 \alpha'^\ast={}^tf(\alpha^\ast).
\]
Indeed, if \(p\) and \(p'\) denote the maps of
\ref{sga3:XXI:1.2.1} for \(\cR\) and \(\cR'\), respectively, then
\[
 p'={}^tf\circ p\circ f,
\]
which gives the assertion at once. The reader may verify the
statements below, almost all of which are immediate.

\begin{proposition}\label{sga3:XXI:6.1.2}
Let \(f:\cR'\to\cR\) be a morphism of root data. If
\(\alpha'\in R'\) and \(\alpha=f(\alpha')\), then
\[
 \alpha'^\ast={}^tf(\alpha^\ast).
\]
Moreover, \(f\) induces isomorphisms
\[
 R'\xrightarrow{\sim}R,\qquad
 \Gamma_0(R')\xrightarrow{\sim}\Gamma_0(R),\qquad
 \VR(R')\xrightarrow{\sim}\VR(R),
\]
and \({}^tf\) induces isomorphisms
\[
 R^\ast\xrightarrow{\sim}R'^\ast,\qquad
 \Gamma_0(R^\ast)\xrightarrow{\sim}\Gamma_0(R'^\ast),\qquad
 \VR(R^\ast)\xrightarrow{\sim}\VR(R'^\ast).
\]
The last is the transpose of the corresponding morphism induced by
\(f\). The map
\[
 s_{\alpha'}\longmapsto s_{f(\alpha')}
\]
extends to an isomorphism
\[
 W(\cR')\xrightarrow{\sim}W(\cR)
\]
compatible with the actions of these two groups on the sets of
\ref{sga3:XXI:1.1.13}.
\end{proposition}

\begin{proposition}\label{sga3:XXI:6.1.3}
The maps
\[
 \Delta'\longmapsto f(\Delta'),\qquad
 R'_+\longmapsto f(R'_+),\qquad
 C'\longmapsto({}^tf\otimes\mathbb R)^{-1}(C')
\]
define bijective correspondences between the systems of simple roots,
the systems of positive roots, and the Weyl chambers for \(\cR'\) and
\(\cR\). These correspondences are compatible with the actions of the
Weyl groups and with the correspondences
\[
 \mathscr{S}(R_+)\longleftrightarrow R_+
 \longleftrightarrow\mathscr{C}(R_+).
\]
\end{proposition}

\begin{lemma}\label{sga3:XXI:6.1.4}
Morphisms compose. A morphism \(f:\cR'\to\cR\) is an isomorphism if
and only if the linear map \(f:M'\to M\) is bijective.
\end{lemma}

% END inputs/XXI/tex/components/22_expose_XXI_section61_morphisms_en.tex

% BEGIN inputs/XXI/tex/components/23_expose_XXI_section62_isogenies_en.tex
\subsection{Isogenies}\label{sga3:XXI:6.2}

\begin{definition}\label{sga3:XXI:6.2.1}
A morphism \(f:\cR'\to\cR\) of root data is called an
\emph{isogeny} if \(f:M'\to M\) is injective with finite cokernel.
\end{definition}

If \(f\) is an isogeny, then \({}^tf\) is an isogeny.

\begin{definition}\label{sga3:XXI:6.2.2}
Let \(f:\cR'\to\cR\) be an isogeny. Put
\[
 K(f)=\operatorname{Coker}\bigl(M'\xrightarrow{f}M\bigr).
\]
\end{definition}

\begin{lemma}\label{sga3:XXI:6.2.3}
There is a natural pairing
\[
 K(f)\times K({}^tf)\longrightarrow\mathbb Q/\mathbb Z
\]
that puts these two finite groups in duality.
\end{lemma}

\begin{proof}
This is classical.
\end{proof}

\begin{lemma}\label{sga3:XXI:6.2.4}
If \(f:\cR'\to\cR\) is a morphism, then
\[
 \rgss(\cR')=\rgss(\cR).
\]
If \(f\) is moreover an isogeny, then
\[
 \rgred(\cR')=\rgred(\cR).
\]
\end{lemma}

\begin{proof}
This is immediate.
\end{proof}

\begin{lemma}\label{sga3:XXI:6.2.5}
Every morphism of semisimple root data is an isogeny.
\end{lemma}

\begin{proof}
The map \(f\) must induce an isomorphism from
\(V'=\VR(R')\) onto \(V=\VR(R)\).
\end{proof}

If \(\cR'\) and \(\cR\) are semisimple, every isogeny
\(f:\cR'\to\cR\) defines a commutative diagram:
\begin{center}
\begin{tikzcd}[column sep=huge,row sep=large]
\Gamma_0(R') \arrow[r,"\sim"] \arrow[d]
  & \Gamma_0(R) \arrow[d] \\
M' \arrow[r,"f"'] & M
\end{tikzcd}
\end{center}
If \(M=\Gamma_0(R)\), then \(f\) is necessarily an isomorphism.

\begin{definition}\label{sga3:XXI:6.2.6}
A root datum is called \emph{adjoint}, respectively
\emph{simply connected}, if
\[
 M=\Gamma_0(R),
 \qquad\text{respectively}\qquad
 M^\ast=\Gamma_0(R^\ast).
\]
\end{definition}

An adjoint or simply connected root datum is therefore semisimple.
Moreover, \(\cR\) is adjoint, respectively simply connected, if and
only if \(\cR^\ast\) is simply connected, respectively adjoint. The
preceding result gives the following proposition.

\begin{proposition}\label{sga3:XXI:6.2.7}
Let \(\cR\) be a semisimple root datum. The following conditions are
equivalent:
\begin{enumerate}
\item[(i)] \(\cR\) is adjoint, respectively simply connected;
\item[(ii)] every isogeny \(\cR'\to\cR\), respectively
  \(\cR\to\cR'\), is an isomorphism.
\end{enumerate}
\end{proposition}

\begin{proposition}\label{sga3:XXI:6.2.8}
Let \(\cR\) be an adjoint, respectively simply connected, root datum.
Every indivisible root, respectively coroot, is an indivisible element
of \(M\), respectively of \(M^\ast\).
\end{proposition}

\begin{proof}
Every indivisible root belongs to a basis of \(\Gamma_0(R)\), by
\ref{sga3:XXI:3.3.5}.
\end{proof}

% END inputs/XXI/tex/components/23_expose_XXI_section62_isogenies_en.tex

% BEGIN inputs/XXI/tex/components/24_expose_XXI_section63_radical_coradical_en.tex
\subsection{Radical and coradical}\label{sga3:XXI:6.3}

Let \(\cR\) be a root datum. Put
\[
 N=\{x\in M\mid(\alpha^\ast,x)=0
       \text{ for every }\alpha^\ast\in R^\ast\},
\qquad
 N^\ast=M^\ast/\bigl(\VR(R^\ast)\cap M^\ast\bigr).
\]

\begin{lemma}\label{sga3:XXI:6.3.1}
Consider the canonical morphisms
\[
 N\longrightarrow M,
 \qquad
 M^\ast\longrightarrow N^\ast.
\]
They are transposes of one another, and \(N^\ast\) is identified with
the dual of \(N\).
\end{lemma}

\begin{proof}
This follows immediately from \ref{sga3:XXI:1.2.5}.
\end{proof}

\begin{definition}\label{sga3:XXI:6.3.2}
The \emph{coradical} of \(\cR\), denoted
\(\operatorname{corad}(\cR)\), is the trivial root datum
\[
 \operatorname{corad}(\cR)=(N,N^\ast,\varnothing,\varnothing).
\]
\end{definition}

If
\[
 \cR^0=(M,M^\ast,\varnothing,\varnothing),
\]
which is a trivial root datum, there is therefore a morphism
\[
 \operatorname{corad}(\cR)\longrightarrow\cR^0.
\]

\begin{definition}\label{sga3:XXI:6.3.3}
The \emph{radical} of \(\cR\), denoted
\(\operatorname{rad}(\cR)\), is the trivial root datum
\footnote{Editors' note: Equivalently, if
\[
 P=\{y\in M^\ast\mid(y,\alpha)=0\text{ for every }\alpha\in R\},
 \qquad
 P^\ast=M/\bigl(\VR(R)\cap M\bigr),
\]
then
\(\operatorname{rad}(\cR)=(P^\ast,P,\varnothing,\varnothing)\).}
\[
 \operatorname{rad}(\cR)=\operatorname{corad}(\cR^\ast)^\ast.
\]
\end{definition}

Thus there is a diagram
\begin{center}
\begin{tikzcd}[column sep=huge,row sep=large]
\operatorname{corad}(\cR) \arrow[rr] \arrow[dr]
  && \operatorname{rad}(\cR) \\
& \cR^0 \arrow[ur] &
\end{tikzcd}
\end{center}
whose transpose is the corresponding diagram for \(\cR^\ast\).

\begin{lemma}\label{sga3:XXI:6.3.4}
The canonical morphism
\[
 u:\operatorname{corad}(\cR)\longrightarrow\operatorname{rad}(\cR)
\]
is an isogeny.
\end{lemma}

\begin{definition}\label{sga3:XXI:6.3.5}
Put
\[
 N(\cR)=K(u)
 =
 M\Big/
 \left(
   \bigl(\VR(R)\cap M\bigr)
   +\bigcap_{\alpha\in R}\ker(\alpha^\ast)
 \right).
\]
There is then a canonical pairing
\[
 N(\cR)\times N(\cR^\ast)\longrightarrow\mathbb Q/\mathbb Z.
\]
\end{definition}

\begin{lemma}\label{sga3:XXI:6.3.6}
One has
\[
 \rgred\bigl(\operatorname{rad}(\cR)\bigr)
 =
 \rgred\bigl(\operatorname{corad}(\cR)\bigr)
 =
 \rgred(\cR)-\rgss(\cR),
\]
and the following conditions are equivalent:
\begin{enumerate}
\item[(i)] \(\cR\) is semisimple;
\item[(ii)] \(\operatorname{rad}(\cR)=0\);
\item[(iii)] \(\operatorname{corad}(\cR)=0\).
\end{enumerate}
\end{lemma}

% END inputs/XXI/tex/components/24_expose_XXI_section63_radical_coradical_en.tex

% BEGIN inputs/XXI/tex/components/25_expose_XXI_section64_products_en.tex
\subsection{Products of root data}\label{sga3:XXI:6.4}

\begin{definition}\label{sga3:XXI:6.4.1}
Let
\[
 \cR=(M,M^\ast,R,R^\ast),
 \qquad
 \cR'=(M',M'^\ast,R',R'^\ast)
\]
be two root data. The \emph{product root datum} of \(\cR\) and
\(\cR'\), denoted
\[
 \cR''=\cR\times\cR',
\]
is the root datum \((M'',M''^\ast,R'',R''^\ast)\), where
\[
\begin{aligned}
 M''&=M\times M',
 &M''^\ast&=M^\ast\times M'^\ast,\\
 R''&=(R\times0)\cup(0\times R'),
 &R''^\ast&=(R^\ast\times0)\cup(0\times R'^\ast),
\end{aligned}
\]
and the map \(\alpha\mapsto\alpha^\ast\) is the evident one.
\end{definition}

\begin{proposition}\label{sga3:XXI:6.4.2}
Under the preceding hypotheses there are canonical isomorphisms
\[
 \Gamma_0(R'')\simeq\Gamma_0(R)\times\Gamma_0(R'),
 \qquad
 \VR(R'')\simeq\VR(R)\times\VR(R'),
\]
\[
 W(\cR'')\simeq W(\cR)\times W(\cR'),
\]
and similarly for the other associated constructions. Moreover,
\[
 \rgred(\cR'')=\rgred(\cR)+\rgred(\cR'),
 \qquad
 \rgss(\cR'')=\rgss(\cR)+\rgss(\cR').
\]
There is also a canonical isomorphism of root data
\[
 (\cR\times\cR')^\ast\simeq\cR^\ast\times\cR'^\ast.
\]
\end{proposition}

The preceding definitions extend immediately to products of any finite
number of factors.

\begin{proposition}\label{sga3:XXI:6.4.3}
Let
\[
 \cR=\cR_1\times\cR_2\times\cdots\times\cR_n
\]
be a product of root data. The following conditions are equivalent:
\begin{enumerate}
\item[(i)] \(\cR\) is semisimple, respectively simply connected,
  respectively adjoint, respectively reduced;
\item[(ii)] every \(\cR_i\) is semisimple, respectively simply
  connected, respectively adjoint, respectively reduced.
\end{enumerate}
\end{proposition}

Consider the following special case. Let \(\cR_0\) be a trivial root
datum and \(\cR_1\) a semisimple root datum. There is then a commutative
diagram
\begin{center}
\begin{tikzcd}[column sep=large,row sep=large]
& \cR_0\times\cR_1 \arrow[dr] & \\
\cR_1 \arrow[ur] \arrow[rr,"\operatorname{id}"]
  && \cR_1
\end{tikzcd}
\end{center}

\begin{lemma}\label{sga3:XXI:6.4.4}
There are canonical isomorphisms
\begin{center}
\begin{tikzcd}[column sep=huge,row sep=large]
\operatorname{corad}(\cR_0\times\cR_1)
  \arrow[rr,"\sim"] \arrow[dr,"\sim"']
  && \operatorname{rad}(\cR_0\times\cR_1) \\
& \cR_0 \arrow[ur,"\sim"'] &
\end{tikzcd}
\end{center}
In particular,
\[
 N(\cR_0\times\cR_1)=0.
\]
\end{lemma}

We shall see later that, conversely, if \(N(\cR)=0\), then \(\cR\) is
the product of a semisimple root datum and a trivial root datum.

% END inputs/XXI/tex/components/25_expose_XXI_section64_products_en.tex

% BEGIN inputs/XXI/tex/components/26_expose_XXI_section65_induced_through_lemma654_en.tex
\subsection{Induced and coinduced root data}\label{sga3:XXI:6.5}

Let \(\cR=(M,M^\ast,R,R^\ast)\) be a root datum, and let
\(N\subset M\) be a subgroup containing the roots, that is,
\[
 \Gamma_0(R)\subset N\subset M.
\]
The canonical linear map \(i_N:N\to M\) gives, by transposition, a
linear map
\[
 {}^ti_N:M^\ast\longrightarrow N^\ast.
\]
Put
\[
 R_N=R,
 \qquad
 R_N^\ast={}^ti_N(R^\ast).
\]

\begin{lemma}\label{sga3:XXI:6.5.1}
The quadruple
\[
 \cR_N=(N,N^\ast,R_N,R_N^\ast)
\]
is a root datum, and \(i_N\) is a morphism.
\end{lemma}

\begin{proof}
First, \({}^ti_N\) induces an isomorphism from \(R^\ast\) onto
\(R_N^\ast\). Indeed, if \(\alpha,\beta\in R\) and
\[
 {}^ti_N(\alpha^\ast)={}^ti_N(\beta^\ast),
\]
then
\[
 (\alpha^\ast,x)=(\beta^\ast,x)
\]
for every \(x\in N\), in particular for \(x\in R\). Hence
\(\alpha=\beta\) by \ref{sga3:XXI:1.1.4}. The rest follows without
difficulty.
\end{proof}

\begin{definition}\label{sga3:XXI:6.5.2}
The root datum \(\cR_N\) is called the root datum
\emph{induced by \(\cR\) on \(N\)}.
\end{definition}

\begin{lemma}\label{sga3:XXI:6.5.3}
Let \(f:\cR'\to\cR\) be a morphism, and put
\[
 N=f(M')\subset M.
\]
Then \(f\) factors uniquely through \(i_N\).
\end{lemma}

In particular, up to isomorphism, the isogenies
\(\cR'\to\cR\) correspond bijectively to the finite-index subgroups of
\[
 M/\Gamma_0(R).
\]
This sharpens \ref{sga3:XXI:6.2.7}.

Now let \(N^\ast\subset M^\ast\) be a subgroup containing \(R^\ast\).
Define the root datum \emph{coinduced by \(\cR\) on \(N^\ast\)} by
\[
 \cR^{N^\ast}=(\cR^\ast_{N^\ast})^\ast.
\]
There is a canonical morphism
\[
 p^{N^\ast}:\cR\longrightarrow\cR^{N^\ast}.
\]

\begin{lemma}\label{sga3:XXI:6.5.4}
Let \(f:\cR'\to\cR\) be a morphism. There exist subgroups
\(N\subset M\) and \(N'^\ast\subset M'^\ast\) such that \(f\) factors
as
\begin{center}
\begin{tikzcd}[column sep=huge,row sep=large]
\cR' \arrow[r,"f"] \arrow[d,"p^{N'^\ast}"']
  & \cR \\
\cR'^{\,N'^\ast} \arrow[r,"f_0","\sim"']
  & \cR_N \arrow[u,"i_N"']
\end{tikzcd}
\end{center}
where \(f_0\) is an isomorphism.
\end{lemma}

\begin{proof}
Take \(N=f(M')\), as in \ref{sga3:XXI:6.5.3}. The resulting morphism
\(M'\to N\) is surjective, so its transpose is injective. Take its
image for \(N'^\ast\).
\end{proof}

% END inputs/XXI/tex/components/26_expose_XXI_section65_induced_through_lemma654_en.tex

% BEGIN inputs/XXI/tex/components/27_expose_XXI_section65_special_cases_en.tex
We now treat some special cases. If \(N=\Gamma_0(R)\), write
\[
 \cR_N=\operatorname{ad}(\cR).
\]
If \(N=\VR(R)\cap M\), write
\[
 \cR_N=\operatorname{ss}(\cR).
\]
Thus there is a diagram
\[
 \operatorname{ad}(\cR)\longrightarrow
 \operatorname{ss}(\cR)\longrightarrow\cR.
\]
Put
\[
 \operatorname{der}(\cR)=\operatorname{ss}(\cR^\ast)^\ast,
 \qquad
 \operatorname{sc}(\cR)=\operatorname{ad}(\cR^\ast)^\ast.
\]
By duality there is a diagram
\[
 \cR\longrightarrow
 \operatorname{der}(\cR)\longrightarrow
 \operatorname{sc}(\cR).
\]

\begin{proposition}\label{sga3:XXI:6.5.5}
\begin{enumerate}
\item[(i)] In the first row of the diagram
\begin{center}
\begin{tikzcd}[column sep=large,row sep=large]
\operatorname{ad}(\cR) \arrow[r]
  & \operatorname{ss}(\cR) \arrow[rr] \arrow[dr]
  && \operatorname{der}(\cR) \arrow[r]
  & \operatorname{sc}(\cR) \\
&& \cR \arrow[ur] &&
\end{tikzcd}
\end{center}
the four root data are semisimple and the three morphisms are
isogenies.
\item[(ii)] The datum \(\operatorname{ad}(\cR)\) is adjoint, and
\(\cR\) is adjoint if and only if
\(\operatorname{ad}(\cR)\to\cR\) is an isomorphism.
\item[(iii)] The datum \(\operatorname{sc}(\cR)\) is simply connected,
and \(\cR\) is simply connected if and only if
\(\cR\to\operatorname{sc}(\cR)\) is an isomorphism.
\item[(iv)] The following conditions are equivalent:
\begin{enumerate}
\item[(a)] \(\cR\) is semisimple;
\item[(b)] \(\operatorname{ss}(\cR)\to\cR\) is an isomorphism;
\item[(c)] \(\cR\to\operatorname{der}(\cR)\) is an isomorphism.
\end{enumerate}
\end{enumerate}
\end{proposition}

Consider the morphism
\[
 \operatorname{ss}(\cR)\longrightarrow\operatorname{der}(\cR).
\]
The definitions immediately give the following result.

\begin{lemma}\label{sga3:XXI:6.5.6}
Let
\[
 h:\operatorname{ss}(\cR)\longrightarrow\operatorname{der}(\cR)
\]
be the canonical isogeny. Then
\[
 K(h)\simeq N(\cR).
\]
\end{lemma}

\SGARefTarget{sga3:XXI:6.5.7}
Consider other special cases of induced root data. Put
\[
 N=
 \{x\in M\mid(\alpha^\ast,x)=0\text{ for every }\alpha\in R\}
 \times\Gamma_0(R).
\]
The sum is direct by \ref{sga3:XXI:1.2.5}. It follows that
\(\cR_N\) is identified with
\[
 \operatorname{ad}(\cR)\times\operatorname{corad}(\cR).
\]
One may replace \(\Gamma_0(R)\) by \(\VR(R)\cap M\), and then dualize
both constructions. This gives the following diagram of root data:
\begin{center}
\resizebox{\textwidth}{!}{%
\begin{tikzpicture}[
  >=stealth,
  every node/.style={font=\small},
  x=1cm,y=1cm
]
\node (at) at (0,2.1) {$\operatorname{ad}(\cR)$};
\node (st) at (4,2.1) {$\operatorname{ss}(\cR)$};
\node (dt) at (10,2.1) {$\operatorname{der}(\cR)$};
\node (ct) at (14,2.1) {$\operatorname{sc}(\cR)$};

\node (am) at (0,0) {$\operatorname{ad}(\cR)\times\operatorname{corad}(\cR)$};
\node (sm) at (4,0) {$\operatorname{ss}(\cR)\times\operatorname{corad}(\cR)$};
\node (rm) at (7,0) {$\cR$};
\node (dm) at (10,0) {$\operatorname{der}(\cR)\times\operatorname{rad}(\cR)$};
\node (cm) at (14,0) {$\operatorname{sc}(\cR)\times\operatorname{rad}(\cR)$};

\node (ab) at (0,-2.1) {$\operatorname{ad}(\cR)$};
\node (sb) at (4,-2.1) {$\operatorname{ss}(\cR)$};
\node (db) at (10,-2.1) {$\operatorname{der}(\cR)$};
\node (cb) at (14,-2.1) {$\operatorname{sc}(\cR)$};

\draw[->] (at) -- (st);
\draw[->] (st) -- (dt);
\draw[->] (dt) -- (ct);
\draw[->] (am) -- (sm);
\draw[->] (sm) -- (rm);
\draw[->] (rm) -- (dm);
\draw[->] (dm) -- (cm);
\draw[->] (ab) -- (sb);
\draw[->] (sb) -- (db);
\draw[->] (db) -- (cb);

\draw[->] (at) -- (am);
\draw[->] (am) -- (ab);
\draw[->] (st) -- (sm);
\draw[->] (sm) -- (sb);
\draw[->] (dt) -- (dm);
\draw[->] (dm) -- (db);
\draw[->] (ct) -- (cm);
\draw[->] (cm) -- (cb);
\end{tikzpicture}
}
\end{center}
It is commutative, as is checked at once, and is self-dual in the
evident sense. The horizontal morphisms are isogenies, and the
composites of the vertical arrows are the identity.

\begin{lemma}\label{sga3:XXI:6.5.8}
Let \(h_1\) and \(h_2\) be the canonical isogenies
\[
 \operatorname{ss}(\cR)\times\operatorname{corad}(\cR)
 \xrightarrow{h_1}\cR
 \xrightarrow{h_2}
 \operatorname{der}(\cR)\times\operatorname{rad}(\cR).
\]
Then
\[
 K(h_1)\simeq K(h_2)\simeq N(\cR).
\]
\end{lemma}

\begin{proof}
This is immediate from the definitions.
\end{proof}

\begin{corollary}\label{sga3:XXI:6.5.9}
For a root datum \(\cR\), the following conditions are equivalent:
\begin{enumerate}
\item[(i)] \(N(\cR)=0\), that is,
  \(\operatorname{corad}(\cR)\to\operatorname{rad}(\cR)\) is an
  isomorphism;
\item[(ii)] the canonical map
  \(\operatorname{ss}(\cR)\to\operatorname{der}(\cR)\) is an
  isomorphism;
\item[(iii)] \(h_1\) is an isomorphism;
\item[(iv)] \(h_2\) is an isomorphism;
\item[(v)] \(\cR\) is the product of a semisimple root datum and a
  trivial root datum.
\end{enumerate}
\end{corollary}

\begin{corollary}\label{sga3:XXI:6.5.10}
For every root datum \(\cR\), there are isogenies
\[
 \operatorname{ad}(\cR)\times\cR_0
 \longrightarrow\cR\longrightarrow
 \operatorname{sc}(\cR)\times\cR_0,
\]
where \(\cR_0\) is the trivial root datum of rank
\(\rgred(\cR)-\rgss(\cR)\).
\end{corollary}

\begin{lemma}\label{sga3:XXI:6.5.11}
Let \(\cR\) be a root datum, let \(\Delta\) be a system of simple
roots, and let \(\Delta'\subset\Delta\). Consider the root datum
\[
 \cR_{\Delta'}=(M,M^\ast,R_{\Delta'},R^\ast_{\Delta'})
\]
of \ref{sga3:XXI:3.4.7}.
\begin{enumerate}
\item[(i)] If \(\cR\) is simply connected, then
  \(\operatorname{der}(\cR_{\Delta'})\) is simply connected.
\item[(ii)] If \(\cR\) is adjoint, then
  \(\operatorname{ss}(\cR_{\Delta'})\) is adjoint.
\end{enumerate}
\end{lemma}

\begin{proof}
The two assertions are equivalent by duality. The second reduces to
\[
 M\cap\VR(R_{\Delta'})=\Gamma_0(R_{\Delta'}).
\]
When \(M=\Gamma_0(R)\), both sides are the subgroup of \(M\) generated
by \(\Delta'\).
\end{proof}

% END inputs/XXI/tex/components/27_expose_XXI_section65_special_cases_en.tex

% BEGIN inputs/XXI/tex/components/28_expose_XXI_section66_weights_partA_en.tex
\subsection{Weights}\label{sga3:XXI:6.6}

\begin{definition}\label{sga3:XXI:6.6.1}
For a root datum \(\cR\), put
\footnote{Editors' note: The notation \(\Gamma(\cR)\) has been replaced
by \(\Lambda(\cR)\) to avoid any possible confusion with
\(\Gamma_0(R)\).}
\[
 \Lambda(\cR)=
 \{x\in\VR(R)\mid(\alpha^\ast,x)\in\mathbb Z
   \text{ for every }\alpha^\ast\in R^\ast\}.
\]
The elements of \(\Lambda(\cR)\) are called the \emph{weights} of
\(\cR\). The weights of \(\cR^\ast\) are called the \emph{coweights}
of \(\cR\).
\end{definition}

One has
\[
 \Gamma_0(R)\subset\Lambda(\cR),
\]
and \(\Lambda(\cR)\) is stable under \(W(\cR)\).

\begin{lemma}\label{sga3:XXI:6.6.2}
The bilinear map
\[
 V^\ast\times V\longrightarrow\mathbb Q
\]
induces a duality
\[
 \Gamma_0(R^\ast)\times\Lambda(\cR)\longrightarrow\mathbb Z.
\]
\end{lemma}

\begin{proof}
Immediate.
\end{proof}

\begin{corollary}\label{sga3:XXI:6.6.3}
Let
\[
 \Delta^\ast=(\alpha_1^\ast,\ldots,\alpha_n^\ast)
\]
be a system of simple coroots, and let \(p_i\in\VR(R)\),
\(i=1,\ldots,n\), be defined by
\[
 (\alpha_i^\ast,p_j)=\delta_{ij}.
\]
Thus
\[
 s_{\alpha_i}(p_i)=p_i-\alpha_i,
 \qquad
 s_{\alpha_i}(p_j)=p_j\quad(i\ne j).
\footnote{Caution: when the datum is not reduced, the \(\alpha_i\)
need not form a system of simple roots.}
\]
Then \(\Lambda(\cR)\) is the free abelian group generated by the
\(p_i\).
\end{corollary}

The \(p_i\) are called the \emph{fundamental weights} corresponding to
the system of simple coroots \(\Delta^\ast\).

\begin{corollary}\label{sga3:XXI:6.6.4}
For every \(\alpha^\ast\in\Delta^\ast\),
\[
 \left(\alpha^\ast,\sum_i p_i\right)=1.
\]
Consequently,
\[
 \sum_i p_i=\rho_{R_+},
 \qquad
 R_+=\mathscr P\bigl(\operatorname{ind}(\Delta)\bigr),
\]
in the notation of \ref{sga3:XXI:3.5.1}.
\end{corollary}

\begin{corollary}\label{sga3:XXI:6.6.5}
For every \(x\in\VR(R)\),
\[
 x=\sum_i(\alpha_i^\ast,x)p_i.
\]
\end{corollary}

Since
\[
 R^\ast\subset\Gamma_0(R^\ast),
 \qquad
 R\subset\Lambda(\cR),
\]
the quadruple
\[
 \bigl(\Lambda(\cR),\Gamma_0(R^\ast),R,R^\ast\bigr)
\]
is a root datum.

\begin{corollary}\label{sga3:XXI:6.6.6}
The canonical morphism
\[
 \Gamma_0(R^\ast)\longrightarrow M^\ast
\]
is the transpose of
\[
 x\longmapsto\sum_i(\alpha_i^\ast,x)p_i.
\]
This map defines a morphism of root data, and there is a commutative
diagram
\begin{center}
\begin{tikzcd}[column sep=huge,row sep=large]
\cR \arrow[rr] \arrow[dr]
  && \operatorname{sc}(\cR) \arrow[dl,"\wr"'] \\
& \bigl(\Lambda(\cR),\Gamma_0(R^\ast),R,R^\ast\bigr) &
\end{tikzcd}
\end{center}
\end{corollary}

% END inputs/XXI/tex/components/28_expose_XXI_section66_weights_partA_en.tex

% BEGIN inputs/XXI/tex/components/29_expose_XXI_section66_partB_through_lemma671_en.tex
Thus we have an explicit description of \(\operatorname{sc}(\cR)\)
in terms of the weights of \(\cR\). Likewise, one obtains a
commutative diagram
\begin{center}
\begin{tikzcd}[column sep=huge,row sep=large]
\operatorname{ad}(\cR) \arrow[rr] \arrow[d,"\wr"]
  && \cR \\
\bigl(\Gamma_0(R),\Lambda(\cR^\ast),R,R^\ast\bigr)
  \arrow[urr] &&
\end{tikzcd}
\end{center}

\begin{corollary}\label{sga3:XXI:6.6.7}
The datum \(\cR\) is simply connected if and only if
\[
 M=\Lambda(\cR).
\]
\end{corollary}

\begin{remark}\label{sga3:XXI:6.6.8}
One has
\[
 \Lambda(\cR)\cap M=\VR(R)\cap M.
\]
Consequently, \(\cR\) is semisimple if and only if
\[
 M\subset\Lambda(\cR).
\]
\end{remark}

The results of \S\ref{sga3:XXI:6.5} also give:

\begin{corollary}\label{sga3:XXI:6.6.9}
The datum \(\cR\) is the product of a simply connected datum and a
trivial datum if and only if
\[
 M\supset\Lambda(\cR).
\]
\end{corollary}

Consider now the canonical isogeny
\[
 f:\operatorname{ad}(\cR)\longrightarrow\operatorname{sc}(\cR),
\]
and put \(Z(\cR)=K(f)\). One has
\[
 Z(\cR)\simeq Z(\operatorname{sc}\cR)
 \simeq Z(\operatorname{ad}\cR).
\]

\begin{corollary}\label{sga3:XXI:6.6.10}
There is a canonical isomorphism
\[
 Z(\cR)=\Lambda(\cR)/\Gamma_0(R).
\]
More precisely, there is an exact sequence of
\(W(\cR)\)-modules
\[
 0\longrightarrow\Gamma_0(R)\longrightarrow\Lambda(\cR)
 \longrightarrow Z(\cR)\longrightarrow0.
\]
\end{corollary}

\begin{corollary}\label{sga3:XXI:6.6.11}
There is a canonical pairing
\[
 Z(\cR^\ast)\times Z(\cR)\longrightarrow\mathbb Q/\mathbb Z
\]
which puts these groups in duality.
\end{corollary}

\begin{remark}\label{sga3:XXI:6.6.12}
One has
\[
 Z(\cR\times\cR')\simeq Z(\cR)\times Z(\cR').
\]
In particular, consider simply connected data
\(\cR_i\), \(i=1,2,\ldots,n\), and a trivial datum \(\cR_0\).
Put
\[
 \cR=\cR_0\times\cR_1\times\cdots\times\cR_n.
\]
If
\[
 \cR=(M,M^\ast,R,R^\ast),
 \qquad
 \cR_0=(M_0,M_0^\ast,\varnothing,\varnothing),
\]
then
\[
 M/\Gamma_0(R)\simeq
 M_0\times Z(\cR_1)\times\cdots\times Z(\cR_n).
\]
\end{remark}

\subsection{Automorphisms}\label{sga3:XXI:6.7}

By \ref{sga3:XXI:6.1.4}, an automorphism of \(\cR\) is an
automorphism \(u\) of \(M\) such that
\[
 u(R)=R,
 \qquad
 {}^tu(R^\ast)=R^\ast.
\]
In particular, every \(w\in W(\cR)\) defines an automorphism of
\(\cR\).

\begin{lemma}\label{sga3:XXI:6.7.1}
The group \(W(\cR)\) is a normal subgroup of
\(\operatorname{Aut}(\cR)\). More precisely, if
\(u\in\operatorname{Aut}(\cR)\) and \(\alpha\in R\), then
\[
 us_\alpha u^{-1}=s_{u(\alpha)}.
\]
\end{lemma}

\begin{proof}
The proof is the same as that of \ref{sga3:XXI:1.2.10}.
\end{proof}

% END inputs/XXI/tex/components/29_expose_XXI_section66_partB_through_lemma671_en.tex

% BEGIN inputs/XXI/tex/components/30_expose_XXI_proposition672_through_lemma677_en.tex
\begin{proposition}\label{sga3:XXI:6.7.2}
Let \(\Delta\) be a system of simple roots, and put
\[
 E_\Delta(\cR)=
 \{u\in\operatorname{Aut}(\cR)\mid u(\Delta)=\Delta\}.
\]
Then \(\operatorname{Aut}(\cR)\) is the semidirect product of
\(W(\cR)\) by \(E_\Delta(\cR)\).
\end{proposition}

\begin{proof}
This follows at once from the fact that \(W(\cR)\) acts simply
transitively on the systems of simple roots and from the fact that,
if \(\Delta\) is a system of simple roots of \(\cR\), then
\(u(\Delta)\) is a system of simple roots for every automorphism
\(u\) of \(\cR\).
\end{proof}

We shall later give a simpler description of \(E_\Delta(\cR)\) in
the case of reduced irreducible root data.\footnote{Editors' note:
When \(\cR\) is simply connected or adjoint, see \SGARefLink{sga3:XXI:7.4.5}{7.4.5}.}

\begin{definition}\label{sga3:XXI:6.7.3}
We denote by \(\operatorname{Aut}^{s}(\cR)\)\footnote{Editors' note:
The superscript \(s\) is intended to suggest ``semisimple.''}
the set of \(u\in\operatorname{Aut}(\cR)\) for which the following
diagram is commutative:
\begin{center}
\begin{tikzcd}[column sep=huge,row sep=large]
\cR \arrow[rr,"u"] \arrow[dr]
  && \cR \arrow[dl] \\
& \operatorname{rad}(\cR) &
\end{tikzcd}
\end{center}
We write
\[
 E_\Delta^{s}(\cR)
 =E_\Delta(\cR)\cap\operatorname{Aut}^{s}(\cR).
\]
\end{definition}

\begin{remark}\label{sga3:XXI:6.7.4}
If \(u\in\operatorname{Aut}(\cR)\), then
\(u\in\operatorname{Aut}^{s}(\cR)\) if and only if
\[
 (u-\operatorname{id})(M)\subset\VR(R).
\]
In particular,
\[
 W(\cR)\subset\operatorname{Aut}^{s}(\cR).
\]
\end{remark}

\begin{proposition}\label{sga3:XXI:6.7.5}
For every system of simple roots \(\Delta\), the group
\(\operatorname{Aut}^{s}(\cR)\) is the semidirect product of
\(W(\cR)\) by \(E_\Delta^{s}(\cR)\).
\end{proposition}

By functoriality, every automorphism of \(\cR\) induces an
automorphism of \(\operatorname{ad}(\cR)\). Thus there is a
canonical homomorphism
\[
 \operatorname{Aut}(\cR)\longrightarrow
 \operatorname{Aut}(\operatorname{ad}(\cR)).
\]

\begin{lemma}\label{sga3:XXI:6.7.6}
The homomorphism
\[
 \operatorname{Aut}^{s}(\cR)\longrightarrow
 \operatorname{Aut}(\operatorname{ad}(\cR))
\]
is injective.
\end{lemma}

\begin{proof}
Let \(u\) be an automorphism of \(M\) such that
\[
 (u-\operatorname{id})(M)\subset\VR(R)
\]
and
\[
 {}^tu(\alpha^\ast)=\alpha^\ast
 \qquad(\alpha^\ast\in R^\ast).
\]
For every \(x\in M\), one has
\[
 (\alpha^\ast,u(x)-x)
 =({}^tu(\alpha^\ast)-\alpha^\ast,x)=0.
\]
Hence \(u(x)-x=0\), by \ref{sga3:XXI:1.2.5}.
\end{proof}

\begin{lemma}\label{sga3:XXI:6.7.7}
The group \(\operatorname{Aut}^{s}(\cR)\) is finite.
\end{lemma}

\begin{proof}
It suffices to prove that \(\operatorname{Aut}(\cR)\) is finite when
\(\cR\) is adjoint. In that case \(M\) is generated by \(R\), so
every automorphism of \(\cR\) is determined by the permutation it
induces on \(R\).
\end{proof}

% END inputs/XXI/tex/components/30_expose_XXI_proposition672_through_lemma677_en.tex

% BEGIN inputs/XXI/tex/components/31_expose_XXI_remark678_through_example686_en.tex
\begin{remark}\label{sga3:XXI:6.7.8}
It is immediate that \(\operatorname{Aut}(\cR)\) (respectively
\(E_\Delta(\cR)\)) is finite if and only if
\[
 \rgred(\cR)-\rgss(\cR)\leq1.
\]
\end{remark}

\subsection{\(p\)-morphisms of reduced root data}
\label{sga3:XXI:6.8}

Throughout this subsection, \(p\) is a fixed positive integer.

\begin{definition}\label{sga3:XXI:6.8.1}
Let
\[
 \cR=(M,M^\ast,R,R^\ast),
 \qquad
 \cR'=(M',M'^\ast,R',R'^\ast)
\]
be two reduced root data. A group homomorphism
\[
 f:M'\longrightarrow M
\]
is called a \emph{\(p\)-morphism from \(\cR'\) to \(\cR\)} if
there exist a bijection
\[
 u:R\xrightarrow{\sim}R'
\]
and a map
\[
 q:R\longrightarrow\{p^n\mid n\in\mathbb N\}
\]
such that:
\begin{enumerate}
\item[(i)] \(f(u(\alpha))=q(\alpha)\alpha\) for every
  \(\alpha\in R\);
\item[(ii)] \({}^tf(\alpha^\ast)=q(\alpha)u(\alpha)^\ast\) for
  every \(\alpha\in R\).\footnote{Editors' note: These two
  conditions imply
  \[
   q(\alpha)\bigl(u(\alpha)^\ast,u(\beta)\bigr)
   =q(\beta)(\alpha^\ast,\beta)
  \]
  for all \(\alpha,\beta\in R\).}
\end{enumerate}
\end{definition}

\begin{corollary}\label{sga3:XXI:6.8.2}
A \(1\)-morphism is precisely a morphism.
\end{corollary}

\begin{corollary}\label{sga3:XXI:6.8.3}
The transpose of a \(p\)-morphism is a \(p\)-morphism.
\end{corollary}

\begin{lemma}\label{sga3:XXI:6.8.4}
If \(w\in W(\cR)\) and \(\alpha\in R\), then
\[
 q(w(\alpha))=q(\alpha).
\]
The map
\[
 s_\alpha\longmapsto s_{u(\alpha)}
\]
extends to an isomorphism
\[
 u:W(\cR)\xrightarrow{\sim}W(\cR')
\]
such that
\[
 u(w(\alpha))=u(w)\bigl(u(\alpha)\bigr).
\]
\end{lemma}

\begin{proof}
It suffices to prove, for \(\alpha,\beta\in R\), that
\[
 u(s_\alpha(\beta))=s_{u(\alpha)}u(\beta),
 \qquad
 q(s_\alpha(\beta))=q(\beta).
\]
Successively, one has
\begin{align*}
 f\bigl(s_{u(\alpha)}u(\beta)\bigr)
  &=f(u(\beta))
    -\bigl(u(\alpha)^\ast,u(\beta)\bigr)f(u(\alpha))\\
  &=q(\beta)\beta
    -q(\beta)q(\alpha)^{-1}
       (\alpha^\ast,\beta)q(\alpha)\alpha\\
  &=q(\beta)\bigl(\beta-(\alpha^\ast,\beta)\alpha\bigr)\\
  &=q(\beta)s_\alpha(\beta).
\end{align*}
If
\[
 \gamma=u^{-1}\bigl(s_{u(\alpha)}u(\beta)\bigr),
\]
then
\[
 q(\gamma)\gamma=f(u(\gamma))=q(\beta)s_\alpha(\beta).
\]
The roots \(\gamma\) and \(s_\alpha(\beta)\) are proportional over
\(\mathbb Q\), hence equal or opposite. Since \(q(\gamma)\) and
\(q(\beta)\) are positive, it follows that
\[
 q(\gamma)=q(\beta),
 \qquad
 \gamma=s_\alpha(\beta).
\]
\end{proof}

\begin{definition}\label{sga3:XXI:6.8.5}
The numbers \(q(\alpha)\) are called the \emph{root exponents} of
\(f\).
\end{definition}

\begin{example}\label{sga3:XXI:6.8.6}
Let \(\cR\) be a reduced root datum and let
\(q=p^n\), \(n\in\mathbb N\). Multiplication by \(q\),
\[
 M\longrightarrow M,\qquad x\longmapsto qx,
\]
is a \(p\)-morphism whose root exponents are all equal to \(q\)
(and for which \(u=\operatorname{id}\)). It is denoted
\[
 q:\cR\longrightarrow\cR.
\]
\end{example}

% END inputs/XXI/tex/components/31_expose_XXI_remark678_through_example686_en.tex

% BEGIN inputs/XXI/tex/components/32_expose_XXI_proposition687_through_proposition711_en.tex
\begin{proposition}\label{sga3:XXI:6.8.7}
In the notation of \ref{sga3:XXI:6.8.1}, \(u\) induces an
isomorphism from the set of systems of simple roots (respectively
positive roots) of \(R\) onto the corresponding set for \(R'\).
\end{proposition}

\begin{proof}
This follows immediately from \ref{sga3:XXI:3.1.5}
(respectively \ref{sga3:XXI:3.2.1}).
\end{proof}

\section{Structure}\label{sga3:XXI:7}

\subsection{Decomposition of a root datum}\label{sga3:XXI:7.1}

\begin{proposition}\label{sga3:XXI:7.1.1}
Let \(\cR\) be a root datum and let \(\Delta\) be a system of
simple roots.
\begin{enumerate}
\item[\SGAInlineRefTarget{sga3:XXI:7.1.1:item:i}(i)] Let \(R'\) and \(R''\) be two closed symmetric sets of
  roots forming a partition of \(R\). If
  \[
   \Delta'=\Delta\cap R',
   \qquad
   \Delta''=\Delta\cap R'',
  \]
  then
  \[
   R'=R_{\Delta'},
   \qquad
   R''=R_{\Delta''},
  \]
  and every root of \(\Delta'\) is orthogonal to every root of
  \(\Delta''\).
\item[\SGAInlineRefTarget{sga3:XXI:7.1.1:item:ii}(ii)] Let \(\Delta'\) and \(\Delta''\) be two orthogonal
  subsets forming a partition of \(\Delta\). Then
  \(R_{\Delta'}\) and \(R_{\Delta''}\) form a partition of \(R\).
\end{enumerate}
\end{proposition}

\begin{proof}
We first prove \SGARefLink{sga3:XXI:7.1.1:item:i}{(i)}.

\begin{lemma}\label{sga3:XXI:7.1.2}
Under the hypotheses of \textup{\SGARefLink{sga3:XXI:7.1.1:item:i}{(i)}}, if
\(\alpha\), \(\beta\), and \(\alpha+\beta\) are roots, then all
three belong to \(R'\), or all three belong to \(R''\).
\end{lemma}

Suppose, for example, that \(\alpha+\beta\in R'\). One cannot have
\(\alpha,\beta\in R''\), since \(R''\) is closed. Suppose therefore
that \(\alpha\in R'\). Then \(-\alpha\in R'\), and
\[
 \beta=(\beta+\alpha)-\alpha\in R'.
\]

We now prove \(R'=R_{\Delta'}\) by induction on the order of a
positive root \(\alpha\in R'\cap\mathscr P(\Delta)\). If
\(\operatorname{ord}_\Delta(\alpha)=1\), then
\[
 \alpha\in R'\cap\Delta=\Delta'.
\]
If \(\operatorname{ord}_\Delta(\alpha)>1\), there is
\(\beta\in\Delta\) such that \(\alpha-\beta\in R\). By the lemma,
\[
 \beta\in\Delta',
 \qquad
 \alpha-\beta\in R'.
\]
Hence \(\alpha-\beta\in R_{\Delta'}\) by induction, and finally
\[
 \alpha=(\alpha-\beta)+\beta\in R_{\Delta'}.
\]

We next prove that \(\Delta'\) and \(\Delta''\) are orthogonal. If
\(\alpha\in\Delta'\) and \(\beta\in\Delta''\), then
\[
 (\beta^\ast,\alpha)\leq0.
\]
If \((\beta^\ast,\alpha)\neq0\), then \(\beta+\alpha\) is a root,
contrary to the lemma.

We now prove \SGARefLink{sga3:XXI:7.1.1:item:ii}{(ii)}. If \(\Delta'\) or \(\Delta''\) is empty, the
assertion is immediate. Otherwise, if \(R_{\Delta'}\) and
\(R_{\Delta''}\) do not form a partition of \(R\), there is a root
\(\alpha\) of the form
\[
 \alpha=\sum_i m_i'\alpha_i'
       +\sum_j m_j''\alpha_j'',
 \qquad
 m_i',m_j''\in\mathbb Z_+,
\]
where the \(\alpha_i'\) (respectively \(\alpha_j''\)) belong to
\(\Delta'\) (respectively \(\Delta''\)). Applying
\ref{sga3:XXI:3.1.2}, and interchanging \(\Delta'\) and
\(\Delta''\) if necessary, yields a relation
\[
 \delta=\gamma+\beta,
 \qquad
 \gamma\in R_{\Delta'},
 \quad
 \beta\in\Delta'',
 \quad
 \delta\in R.
\]
But \((\beta^\ast,\gamma)=0\), so
\(\gamma-\beta\) is also a root by \ref{sga3:XXI:2.2.5}, which is
impossible.
\end{proof}

% END inputs/XXI/tex/components/32_expose_XXI_proposition687_through_proposition711_en.tex

% BEGIN inputs/XXI/tex/components/33_expose_XXI_proposition713_through_lemma723_en.tex
\begin{proposition}\label{sga3:XXI:7.1.3}
Let \(\cR\) be a root datum. The following conditions are
equivalent:
\begin{enumerate}
\item[\SGAInlineRefTarget{sga3:XXI:7.1.3:item:i}(i)] There is no nontrivial partition of \(R\) into two
  closed symmetric subsets.
\item[\SGAInlineRefTarget{sga3:XXI:7.1:item:ii}(ii)] For some (respectively every) system of simple roots
  \(\Delta\) of \(R\), there is no partition of \(\Delta\) into
  two nonempty orthogonal subsets.
\item[\SGAInlineRefTarget{sga3:XXI:7.1.3:item:iii}(iii)] The natural representation of \(W(\cR)\) on
  \(\VR(R)\) is irreducible.
\item[\SGAInlineRefTarget{sga3:XXI:7.1.3:item:iv}(iv)] For every pair \((\alpha,\beta)\) of roots, there is a
  sequence of roots
  \[
   \alpha_0,\alpha_1,\ldots,\alpha_n,
   \qquad
   \alpha_0=\alpha,\quad\alpha_n=\beta,
  \]
  such that \(\alpha_i\) and \(\alpha_{i+1}\) are nonorthogonal
  for \(i=0,\ldots,n-1\).
\end{enumerate}
\end{proposition}

\begin{proof}
The equivalence \SGARefLink{sga3:XXI:7.1.3:item:i}{(i)}\(\Leftrightarrow\)\SGARefLink{sga3:XXI:7.1:item:ii}{(ii)} follows from
\ref{sga3:XXI:7.1.1}, and \SGARefLink{sga3:XXI:7.1.3:item:iv}{(iv)}\(\Rightarrow\)\SGARefLink{sga3:XXI:7.1.1:item:ii}{(ii)} is immediate.
Conversely, if \SGARefLink{sga3:XXI:7.1.1:item:ii}{(ii)} holds for \(\Delta\), then \SGARefLink{sga3:XXI:7.1.3:item:iv}{(iv)} holds when
\(\alpha,\beta\in\Delta\). Every root is nonorthogonal to some
simple root, by \ref{sga3:XXI:3.1.1}, for example. Moreover,
\SGARefLink{sga3:XXI:7.1.3:item:iii}{(iii)}\(\Rightarrow\)\SGARefLink{sga3:XXI:7.1.3:item:i}{(i)}: under the hypotheses of
\ref{sga3:XXI:7.1.1}, the space \(\VR(R')\) is stable under
\(W(\cR)\). It remains to prove \SGARefLink{sga3:XXI:7.1.1:item:i}{(i)}\(\Rightarrow\)\SGARefLink{sga3:XXI:7.1.3:item:iii}{(iii)}.

Let \(H\) be a vector subspace of \(\VR(R)\), stable under
\(W(\cR)\). For every \(\alpha\in R\), the equality
\(s_\alpha(H)=H\) immediately gives
\[
 \alpha\in H
 \qquad\text{or}\qquad
 \alpha^\ast\in H^\perp,
\]
where \(H^\perp\) denotes the orthogonal complement of \(H\) in
\(\VR(R^\ast)\), which is in duality with \(\VR(R)\). Put
\[
 R'=\{\alpha\in R\mid\alpha\in H\},
 \qquad
 R''=\{\alpha\in R\mid\alpha^\ast\in H^\perp\}.
\]
This gives a partition of \(R\) into two closed symmetric subsets.
\end{proof}

\begin{definition}\label{sga3:XXI:7.1.4}
A root datum (respectively root system) satisfying the equivalent
conditions of \ref{sga3:XXI:7.1.3} and having nonzero semisimple
rank is called \emph{irreducible}.
\end{definition}

\begin{corollary}\label{sga3:XXI:7.1.5}
For every root datum \(\cR\), there is a unique partition, up to
order, of \(R\) into closed symmetric irreducible subsets.
\end{corollary}

\begin{corollary}\label{sga3:XXI:7.1.6}
Every adjoint (respectively simply connected) root datum is a
product of irreducible adjoint (respectively simply connected) root
data.
\end{corollary}

\begin{proof}
It suffices to consider the adjoint case. The assertion then follows
from the fact that, under the hypotheses of
\ref{sga3:XXI:7.1.1},
\[
 \Gamma_0(R)=\Gamma_0(R')\times\Gamma_0(R'').
\]
\end{proof}

\begin{corollary}\label{sga3:XXI:7.1.7}
For every root datum (respectively reduced root datum) \(\cR\),
there is an isogeny
\[
 \cR\longrightarrow\cR',
\]
where \(\cR'\) is a product of a trivial root datum and irreducible
simply connected root data (respectively irreducible reduced simply
connected root data).
\end{corollary}

\subsection{Properties of irreducible root data}
\label{sga3:XXI:7.2}

\begin{definition}\label{sga3:XXI:7.2.1}
Let \(\cR\) be an irreducible root datum. For every \(\alpha\in R\),
put
\[
 \operatorname{long}(\alpha)
 =\ell(\alpha)/\ell(\alpha_0),
\]
where \(\alpha_0\in R\) is such that \(\ell(\alpha_0)\) is minimal.
The number \(\operatorname{long}(\alpha)\) is called the
\emph{length} of \(\alpha\).
\end{definition}

\begin{lemma}\label{sga3:XXI:7.2.2}
Let \(\cR\) be an irreducible root datum. The Weyl group acts
transitively on the set of roots of a given length.
\end{lemma}

\begin{proof}
Let \(\alpha,\beta\in R\). Since the representation of \(W(\cR)\)
on \(\VR(R)\) is irreducible, \(\alpha\) cannot be orthogonal to
all \(w(\beta)\), \(w\in W(\cR)\). Hence some \(w(\beta)\) is
nonorthogonal to \(\alpha\). Since
\[
 \ell(w(\beta))=\ell(\beta),
\]
the assertion follows from \ref{sga3:XXI:2.3.2}.
\end{proof}

\begin{lemma}\label{sga3:XXI:7.2.3}
If \(\cR\) is irreducible and reduced, then
\[
 \operatorname{long}(R)
 \in\bigl\{\{1\},\{1,2\},\{1,3\}\bigr\}.
\]
\end{lemma}

\begin{proof}
By the observation used above, for every
\(\alpha,\beta\in R\) there is \(w\in W(\cR)\) such that
\(w(\beta)\) is nonorthogonal to \(\alpha\). Therefore, by
\ref{sga3:XXI:2.3.1},
\[
 \frac{\ell(\alpha)}{\ell(\beta)}
 \in\left\{1,2,3,\frac12,\frac13\right\}.
\]
Thus \(\operatorname{long}(\alpha)\) is \(1\), \(2\), or \(3\).
But lengths \(2\) and \(3\) cannot both occur, since their ratio
\(2/3\) is not in the displayed set.
\end{proof}

% END inputs/XXI/tex/components/33_expose_XXI_proposition713_through_lemma723_en.tex

% BEGIN inputs/XXI/tex/components/34_expose_XXI_remark724_through_remark729_en.tex
\begin{remark}\label{sga3:XXI:7.2.4}
Reasoning similarly gives the following result.  If \(\cR\) is
irreducible and nonreduced, with \(\rgss(\cR)>1\), then
\[
 \operatorname{long}(R)=\{1,2,4\}.
\]
If \(R_i=\operatorname{long}^{-1}(i)\), then
\[
 \operatorname{ind}(R)=R_1\cup R_2,\qquad R_4=2R_1,
\]
and two nonproportional roots of \(R_1\) are orthogonal.
Conversely, suppose that \(R\) is an irreducible reduced root
system such that \(\operatorname{long}(R)=\{1,2\}\).  Put
\(R_i=\operatorname{long}^{-1}(i)\), and suppose that two
nonproportional roots of \(R_1\) are orthogonal.  Then
\(R\cup2R_1\) is irreducible and nonreduced, and
\[
 \operatorname{ind}(R\cup2R_1)=R.
\]
\end{remark}

\begin{lemma}\label{sga3:XXI:7.2.5}
If \(\cR\) is an irreducible root datum, then \(\cR^\ast\) is
irreducible as well, and the product
\[
 \operatorname{long}(\alpha)\operatorname{long}(\alpha^\ast)
\]
is constant as \(\alpha\) runs through \(R\).
\end{lemma}

\begin{proof}
This follows immediately from \ref{sga3:XXI:7.1.3}\,\SGARefLink{sga3:XXI:7.2:item:iv}{(iv)} and
\ref{sga3:XXI:2.2.6}.
\end{proof}

\begin{definition}\label{sga3:XXI:7.2.6}
Let \(\cR\) be any root datum.  The \emph{length} of
\(\alpha\in R\), denoted \(\operatorname{long}(\alpha)\), is the
length of \(\alpha\) in its irreducible component.
\end{definition}

\begin{lemma}\label{sga3:XXI:7.2.7}
There is a unique group homomorphism
\[
 u:\Gamma_0(R)\longrightarrow\Gamma_0(R^\ast)
\]
such that
\[
 u(\alpha)=\operatorname{long}(\alpha)\alpha^\ast
 \qquad(\alpha\in R).
\]
\end{lemma}

\begin{proof}
By \ref{sga3:XXI:3.5.5}, it is enough to verify that, whenever
\(\alpha,\beta,\alpha+\beta\in R\),
\[
 \operatorname{long}(\alpha)\alpha^\ast+
 \operatorname{long}(\beta)\beta^\ast
 =
 \operatorname{long}(\alpha+\beta)(\alpha+\beta)^\ast.
\]
But \(\alpha,\beta,\alpha+\beta\) belong to the same irreducible
component of \(R\), by \ref{sga3:XXI:7.1.2}; the assertion is
therefore reduced to \ref{sga3:XXI:1.2.2}.
\end{proof}

\begin{remark}\label{sga3:XXI:7.2.8}
Let \(u\) be as in \ref{sga3:XXI:7.2.7}.  For
\(\alpha,\beta\in R\),
\[
 (u(\alpha),\beta)=(u(\beta),\alpha).
\]
Indeed, this amounts to
\[
 \operatorname{long}(\alpha)(\alpha^\ast,\beta)
 =
 \operatorname{long}(\beta)(\beta^\ast,\alpha),
\]
which is clear when \(\alpha\) and \(\beta\) are orthogonal.  If
they are not orthogonal, they lie in the same irreducible component
of \(R\), and the assertion reduces to \ref{sga3:XXI:1.2.1},
formula~\SGARefLink{sga3:XXI:eq:9}{(9)}.
\end{remark}

\begin{remark}\label{sga3:XXI:7.2.9}
The symmetric bilinear form
\[
 (x,y)\longmapsto(u(x),y)
\]
is positive and nondegenerate on \(\Gamma_0(R)\).
\end{remark}

\begin{proof}
Let \(R_i\) be the irreducible components of \(R\).  Then
\[
 \Gamma_0(R)=\prod_i\Gamma_0(R_i),
\]
and the bilinear form \((u(x),y)\) is the product, on the groups
\(\Gamma_0(R_i)\), of the forms
\[
 2\ell(\alpha_i)^{-1}(p(x),y),
\]
where \(\ell(\alpha_i)\) is the minimum of \(\ell(\alpha)\) for
\(\alpha\in R_i\).  These symmetric bilinear forms are positive
and nondegenerate by \ref{sga3:XXI:1.2.6}.
\end{proof}

% END inputs/XXI/tex/components/34_expose_XXI_remark724_through_remark729_en.tex

% BEGIN inputs/XXI/tex/components/35_expose_XXI_section73_cartan_matrix_en.tex
\subsection{Cartan matrix}\label{sga3:XXI:7.3}

Let \(\cR\) be a root datum.  If \(\Delta\) is a system of simple
roots, the \emph{Cartan matrix} of \(\cR\) relative to \(\Delta\)
is the square matrix indexed by \(\Delta\) and defined by
\[
 a_{\alpha,\beta}=(\alpha^\ast,\beta),
 \qquad \alpha,\beta\in\Delta.
\]
If \(\Delta'\) is another system of simple roots and
\(w\in W(\cR)\) satisfies \(w(\Delta)=\Delta'\), then
\[
 (w(\alpha)^\ast,w(\beta))=(\alpha^\ast,\beta).
\]
Thus the Cartan matrix relative to \(\Delta'\) is obtained from
that relative to \(\Delta\) by the isomorphism
\(\Delta\to\Delta'\) of index sets defined by \(w\).  Consequently,
up to isomorphism of its index set, the Cartan matrix depends only
on \(\cR\).

\begin{proposition}\label{sga3:XXI:7.3.1}
The Cartan matrix has the following properties.
\begin{enumerate}
\item[\SGAInlineRefTarget{sga3:XXI:7.3.1:item:i}(i)] \(a_{\alpha,\alpha}=2\), and
  \(a_{\alpha,\beta}\leq0\) for \(\alpha\neq\beta\).
\item[\SGAInlineRefTarget{sga3:XXI:7.3.1:item:ii}(ii)] \(a_{\alpha,\beta}=0\) implies
  \(a_{\beta,\alpha}=0\).
\item[\SGAInlineRefTarget{sga3:XXI:7.2:item:iii}(iii)] There are strictly positive integers
  \(m_\alpha=\operatorname{long}(\alpha)\) such that the matrix
  \[
   (m_\alpha a_{\alpha,\beta})
  \]
  is symmetric, positive, and nondegenerate.
\item[\SGAInlineRefTarget{sga3:XXI:7.2:item:iv}(iv)] The principal minors of
  \((a_{\alpha,\beta})_{\alpha,\beta\in\Delta}\), namely
  \[
   \det(a_{\alpha,\beta})_{\alpha,\beta\in\Delta'}
   \qquad(\Delta'\subset\Delta),
  \]
  are strictly positive.
\item[\SGAInlineRefTarget{sga3:XXI:7.2:item:v}(v)] One has
  \[
   s_\alpha(\beta)=\beta-a_{\alpha,\beta}\alpha,
   \qquad
   s_\alpha(\beta^\ast)
   =\beta^\ast-a_{\beta,\alpha}\alpha^\ast.
  \]
\end{enumerate}
\end{proposition}

\begin{proof}
Part \SGARefLink{sga3:XXI:7.2:item:v}{(v)} is a definition; \SGARefLink{sga3:XXI:7.3.1:item:i}{(i)} follows from
\ref{sga3:XXI:3.2.11}, \SGARefLink{sga3:XXI:7.3.1:item:ii}{(ii)} from \ref{sga3:XXI:2.2.2}, and
\SGARefLink{sga3:XXI:7.2:item:iii}{(iii)} from \ref{sga3:XXI:7.2.9}.  Part \SGARefLink{sga3:XXI:7.2:item:iv}{(iv)} follows immediately
from \SGARefLink{sga3:XXI:7.2:item:iii}{(iii)} and
\[
 \det(m_\alpha a_{\alpha,\beta})_{\alpha,\beta\in\Delta'}
 =
 \left(\prod_{\alpha\in\Delta'}m_\alpha\right)
 \det(a_{\alpha,\beta})_{\alpha,\beta\in\Delta'}.
\]
\end{proof}

\begin{proposition}\label{sga3:XXI:7.3.2}
Let \(\cR\) and \(\cR'\) be two reduced simply connected
(respectively adjoint) root data, let \(\Delta\) and \(\Delta'\)
be systems of simple roots of \(\cR\) and \(\cR'\), and let
\(u:\Delta\to\Delta'\) be an isomorphism.  Denote the corresponding
Cartan matrices by \((a_{\alpha,\beta})\) and
\((a'_{\alpha',\beta'})\), and suppose that
\[
 a'_{u(\alpha),u(\beta)}=a_{\alpha,\beta}.
\]
Then there is a unique isomorphism from \(\cR\) onto \(\cR'\)
that induces \(u\) on \(\Delta\).
\end{proposition}

\begin{proof}
It is enough to give the proof in the adjoint case.  Then
\(M=\Gamma_0(R)\) and \(M'=\Gamma_0(R')\) are the free abelian
groups generated by \(\Delta\) and \(\Delta'\).  Hence there is a
unique group isomorphism \(M\to M'\) inducing \(u\) on \(\Delta\);
denote it again by \(u\).

We show that \(u(R)\subset R'\).  Every root \(\alpha\) of \(R\)
can be written
\[
 \alpha=s_{\alpha_1}\cdots s_{\alpha_n}(\alpha_{n+1}),
 \qquad \alpha_i\in\Delta.
\]
By the hypothesis on \(u\) and the relations
\ref{sga3:XXI:7.3.1}\,\SGARefLink{sga3:XXI:7.2:item:v}{(v)},
\[
 u(\alpha)=
 s_{u(\alpha_1)}\cdots s_{u(\alpha_n)}
 \bigl(u(\alpha_{n+1})\bigr).
\]
It remains to prove that \({}^tu(R'^\ast)\subset R^\ast\).  This
follows from the fact that the elements of \(M^\ast\)
(respectively \(M'^\ast\)) are determined by their duality with
\(R\) or \(\Delta\) (respectively \(R'\) or \(\Delta'\)), by
\ref{sga3:XXI:1.2.5}.
\end{proof}

\begin{corollary}\label{sga3:XXI:7.3.3}
A reduced simply connected or adjoint root datum is determined up
to isomorphism by its Cartan matrix.
\end{corollary}

\begin{corollary}\label{sga3:XXI:7.3.4}
Let \(\cR\) be a reduced simply connected (respectively adjoint)
root datum and let \(\Delta\) be a system of simple roots.  The
group \(E_\Delta(\cR)\) is identified with the group of
automorphisms of the set \(\Delta\) that preserve the Cartan
matrix.
\end{corollary}

\begin{remark}\label{sga3:XXI:7.3.5}
The existence of a root datum corresponding to a given Cartan
matrix satisfying, for example, \SGARefLink{sga3:XXI:7.3.1:item:i}{(i)}, \SGARefLink{sga3:XXI:7.3.1:item:ii}{(ii)}, and \SGARefLink{sga3:XXI:7.2:item:iv}{(iv)} is not easily
settled directly without using the classification.
\end{remark}

% END inputs/XXI/tex/components/35_expose_XXI_section73_cartan_matrix_en.tex

% BEGIN inputs/XXI/tex/components/36_expose_XXI_section74_dynkin_diagrams_en.tex
\subsection{Dynkin diagrams}\label{sga3:XXI:7.4}

\begin{definition}\label{sga3:XXI:7.4.1}
A \emph{Dynkin-diagram structure} on a finite set \(\Delta\)
(the word ``scheme'' being banished for obvious reasons) consists
of a set of pairs of distinct elements of \(\Delta\), called
\emph{linked pairs}, together with a map
\[
 \Delta\longrightarrow\{1,2,3\}.
\]
The notion of isomorphism between such structures is evident.
\end{definition}

\begin{definition}\label{sga3:XXI:7.4.2}
Let \(\cR\) be a root datum and let \(\Delta\) be a system of
simple roots.  The \emph{Dynkin diagram} of \(\cR\) relative to
\(\Delta\) is the set \(\Delta\), in which two simple roots are
linked if and only if they are nonorthogonal, and each root is
assigned its length.
\end{definition}

\begin{proposition}\label{sga3:XXI:7.4.3}
The Dynkin diagram and the Cartan matrix determine one another
uniquely.
\end{proposition}

\begin{proof}
The equivalence
\[
 \alpha\text{ and }\beta\text{ are not linked}
 \quad\Longleftrightarrow\quad
 a_{\alpha,\beta}=0
\]
and the relation
\[
 \operatorname{long}(\alpha)a_{\alpha,\beta}
 =
 \operatorname{long}(\beta)a_{\beta,\alpha},
\]
with \(\inf\operatorname{long}(\alpha)=1\) in each connected
component of the diagram, determine the \(a_{\alpha,\beta}\) from
the links and lengths, and conversely.  The details are left to the
reader.
\end{proof}

\begin{corollary}\label{sga3:XXI:7.4.4}
A reduced simply connected or adjoint root datum is determined by
its Dynkin diagram.
\end{corollary}

\begin{corollary}\label{sga3:XXI:7.4.5}
Let \(\cR\) be a reduced simply connected or adjoint root datum and
let \(\Delta\) be a system of simple roots.  The group
\(E_\Delta(\cR)\) is identified with the group of automorphisms of
the Dynkin diagram of \(\cR\) relative to \(\Delta\); that is, with
the group of permutations of \(\Delta\) preserving lengths and
links.
\end{corollary}

\begin{remark}\label{sga3:XXI:7.4.6}
The usual method\footnote{See Bourbaki, \emph{Groupes et algèbres
de Lie}, Chapter~VI, no.~4.2, or the Sophus Lie Seminar.}
\footnote{Editors' note: for a slightly different proof, see also
[De80].}
classifies the connected Dynkin diagrams and shows that each
actually corresponds to an irreducible reduced simply connected
root datum.  One obtains the familiar types:
\begin{center}
% BEGIN inputs/XXI/tex/figures/xxi_p44_dynkin_classification_native.tex
\begin{tikzpicture}[
  x=.86cm,y=.78cm,
  dyn/.style={circle,fill=black,inner sep=1.35pt},
  wt/.style={font=\small,above=2pt},
  typ/.style={font=\normalsize,anchor=west}
]
% A_n
\foreach \x in {0,1,2,4,5,6} {
  \node[dyn] (A\x) at (\x,0) {};
  \node[wt] at (A\x) {$1$};
}
\draw (A0)--(A1)--(A2);
\draw (A4)--(A5)--(A6);
\node at (3,0) {$\cdots$};
\node[typ] at (7.3,0) {$A_n,\qquad n\geq1.$};

% B_n
\foreach \x in {0,1,2,4,5} {
  \node[dyn] (B\x) at (\x,-1.55) {};
  \node[wt] at (B\x) {$2$};
}
\node[dyn] (B6) at (6,-1.55) {};
\node[wt] at (B6) {$1$};
\draw (B0)--(B1)--(B2);
\draw (B4)--(B5)--(B6);
\node at (3,-1.55) {$\cdots$};
\node[typ] at (7.3,-1.55) {$B_n,\qquad n\geq2.$};

% C_n
\foreach \x in {0,1,2,4,5} {
  \node[dyn] (C\x) at (\x,-3.10) {};
  \node[wt] at (C\x) {$1$};
}
\node[dyn] (C6) at (6,-3.10) {};
\node[wt] at (C6) {$2$};
\draw (C0)--(C1)--(C2);
\draw (C4)--(C5)--(C6);
\node at (3,-3.10) {$\cdots$};
\node[typ] at (7.3,-3.10) {$C_n,\qquad n\geq3.$};

% D_n
\foreach \x in {0,1,2,4} {
  \node[dyn] (D\x) at (\x,-4.65) {};
  \node[wt] at (D\x) {$1$};
}
\node[dyn] (Du) at (5,-4.25) {};
\node[wt] at (Du) {$1$};
\node[dyn] (Dl) at (5,-5.05) {};
\node[wt] at (Dl) {$1$};
\draw (D0)--(D1)--(D2);
\draw (D4)--(Du);
\draw (D4)--(Dl);
\node at (3,-4.65) {$\cdots$};
\node[typ] at (7.3,-4.65) {$D_n,\qquad n\geq4.$};

% E_n
\foreach \x in {0,1,3,4,5} {
  \node[dyn] (E\x) at (\x,-6.35) {};
  \node[wt] at (E\x) {$1$};
}
\node[dyn] (Ev) at (3,-7.15) {};
\node[right=2pt,font=\small] at (Ev) {$1$};
\draw (E0)--(E1);
\draw (E3)--(E4)--(E5);
\draw (E3)--(Ev);
\node at (2,-6.35) {$\cdots$};
\node[typ] at (7.3,-6.65) {$E_n,\qquad n=6,7,8.$};

% F_4
\foreach \x/\w in {0/1,1/1,2/2,3/2} {
  \node[dyn] (F\x) at (\x,-8.15) {};
  \node[wt] at (F\x) {$\w$};
}
\draw (F0)--(F1)--(F2)--(F3);
\node[typ] at (7.3,-8.15) {$F_4.$};

% G_2
\node[dyn] (G0) at (2,-9.70) {};
\node[wt] at (G0) {$1$};
\node[dyn] (G1) at (3,-9.70) {};
\node[wt] at (G1) {$3$};
\draw (G0)--(G1);
\node[typ] at (7.3,-9.70) {$G_2.$};
\end{tikzpicture}

% END inputs/XXI/tex/figures/xxi_p44_dynkin_classification_native.tex

\end{center}
By \ref{sga3:XXI:7.4.5}, the corresponding group
\(E_\Delta(\cR)\) is obtained at once:
\[
\begin{array}{ll}
E_\Delta(\cR)=\{e\},
 & A_1,\ B_n,\ C_n,\ E_7,\ E_8,\ F_4,\ G_2,\\[2pt]
E_\Delta(\cR)=\mathbb Z/2\mathbb Z,
 & A_n\ (n\geq2),\ D_n\ (n\geq5),\ E_6,\\[2pt]
E_\Delta(\cR)=\mathfrak S_3,
 & D_4.
\end{array}
\]
\end{remark}

% END inputs/XXI/tex/components/36_expose_XXI_section74_dynkin_diagrams_en.tex

% BEGIN inputs/XXI/tex/components/37_expose_XXI_section75_p_morphisms_and_bibliography_en.tex
\subsection{Complements on \(p\)-morphisms}\label{sga3:XXI:7.5}

Let \(f:\cR\to\cR'\) be a \(p\)-morphism
(cf.\ \ref{sga3:XXI:6.8}).  It is clear from the definitions that
the bijection \(u:R\xrightarrow{\sim}R'\) associated with \(f\)
makes the systems of simple roots, the systems of positive roots,
the irreducible components, and so forth of \(R\) and \(R'\)
correspond.  We may therefore suppose, for simplicity, that
\(R\) and \(R'\) are irreducible.

\begin{lemma}\label{sga3:XXI:7.5.1}
If \(R\) and \(R'\) are irreducible, there is a \(k\in\mathbb Q\)
such that, for every \(\alpha\in R\),
\[
 k\,\operatorname{long}(u(\alpha))
 =
 q(\alpha)^2\operatorname{long}(\alpha).
\]
\end{lemma}

\begin{proof}
One has
\[
 \operatorname{long}(\alpha)(\alpha^\ast,\beta)
 =
 \operatorname{long}(\beta)(\beta^\ast,\alpha)
\]
and, similarly,
\[
 \operatorname{long}(u(\alpha))
   (u(\alpha)^\ast,u(\beta))
 =
 \operatorname{long}(u(\beta))
   (u(\beta)^\ast,u(\alpha)).
\]
It follows immediately that, for nonorthogonal \(\alpha,\beta\),
\[
 \frac{q(\alpha)^2\operatorname{long}(\alpha)}
      {\operatorname{long}(u(\alpha))}
 =
 \frac{q(\beta)^2\operatorname{long}(\beta)}
      {\operatorname{long}(u(\beta))}.
\]
The conclusion then follows from \ref{sga3:XXI:7.1.3}\,\SGARefLink{sga3:XXI:7.1.3:item:iv}{(iv)}.
\end{proof}

\begin{remark}\label{sga3:XXI:7.5.2}
It follows from \ref{sga3:XXI:7.2.2} and
\ref{sga3:XXI:6.8.4} that \(q(\alpha)\) depends only on
\(\operatorname{long}(\alpha)\).  One then sees readily that, if
\(q(\alpha)\) is not constant, the product
\(q(\alpha)\operatorname{long}(\alpha)\) is constant; hence
\(p=2\) or \(3\).  Inspection of the preceding diagrams gives the
four possible cases (the same letter denotes a Dynkin diagram and
the corresponding reduced simply connected root datum):
\[
\begin{array}{llll}
p=2, & B_n\xrightarrow{\,f_1\,}C_n,
     & C_n\xrightarrow{\,f_2\,}B_n,
     & (C_2=B_2),\\
p=2, & F_4\xrightarrow{\,g\,}F_4,\\
p=3, & G_2\xrightarrow{\,h\,}G_2.
\end{array}
\]
The reader will observe that
\(f_1\circ f_2\), \(f_2\circ f_1\), \(g\circ g\), and
\(h\circ h\) are \(p\)-morphisms of the form described in
\ref{sga3:XXI:6.8.6}.
\end{remark}

\medskip
\noindent\textbf{7.5.3.}\label{sga3:XXI:7.5.3}
The preceding description shows immediately that, if \(\cR\) and
\(\cR'\) are two reduced root data of semisimple rank at most \(2\)
and there is a \(p\)-morphism from \(\cR'\) to \(\cR\), then
\(\cR\) and \(\cR'\) have the same type.  More precisely, one has
the following table.

\medskip
\noindent\emph{Notation.}
Let \(f:\cR'\to\cR\) be a \(p\)-morphism.  The symbols \(q\) and
\(q_1\) denote arbitrary positive powers of \(p\).  For rank-two
systems we use the notation of Section~\ref{sga3:XXI:4}, with
\(\alpha,\beta\) the simple roots and
\(\ell(\alpha)\leq\ell(\beta)\).

\begin{center}
\begingroup
\renewcommand{\arraystretch}{1.35}
\resizebox{\linewidth}{!}{%
\begin{tabular}{|c|c|c|c|}
\hline
\textbf{Type} & \(p\) & \textbf{values of \(f\)}
 & \textbf{values of \({}^tf\)}\\
\hline
Trivial & arbitrary & --- & ---\\
\hline
\(A_1\) & arbitrary
 & \(f(\alpha')=q\alpha\)
 & \({}^tf(\alpha^\ast)=q\alpha'^\ast\)\\
\hline
\(A_1\times A_1\) & arbitrary
 & \(\begin{array}{c}f(\alpha')=q\alpha\\
                      f(\beta')=q_1\beta\end{array}\)
 & \(\begin{array}{c}{}^tf(\alpha^\ast)=q\alpha'^\ast\\
                      {}^tf(\beta^\ast)=q_1\beta'^\ast\end{array}\)\\
\hline
\(A_2,B_2,G_2\) & arbitrary
 & \(\begin{array}{c}f(\alpha')=q\alpha\\
                      f(\beta')=q\beta\end{array}\)
 & \(\begin{array}{c}{}^tf(\alpha^\ast)=q\alpha'^\ast\\
                      {}^tf(\beta^\ast)=q\beta'^\ast\end{array}\)\\
\hline
\(B_2\) & \(p=2\)
 & \(\begin{array}{c}f(\alpha')=q\beta\\
                      f(\beta')=2q\alpha\end{array}\)
 & \(\begin{array}{c}{}^tf(\alpha^\ast)=2q\beta'^\ast\\
                      {}^tf(\beta^\ast)=q\alpha'^\ast\end{array}\)\\
\hline
\(G_2\) & \(p=3\)
 & \(\begin{array}{c}f(\alpha')=q\beta\\
                      f(\beta')=3q\alpha\end{array}\)
 & \(\begin{array}{c}{}^tf(\alpha^\ast)=3q\beta'^\ast\\
                      {}^tf(\beta^\ast)=q\alpha'^\ast\end{array}\)\\
\hline
\end{tabular}}
\endgroup
\end{center}

\section*{Bibliography}
\addcontentsline{toc}{section}{Bibliography}
\label{sga3:XXI:bibliography}

\begin{description}
\item[{[BLie]}] N.~Bourbaki,
  \emph{Groupes et algèbres de Lie}, Chapters IV--VI,
  Hermann, 1968.
\item[{[De80]}] M.~Demazure,
  ``A, B, C, D, E, F, etc.,'' pp.~221--227 in
  \emph{Séminaire sur les singularités des surfaces}
  (Palaiseau, 1976--1977), eds.\ M.~Demazure,
  H.~C.~Pinkham, and B.~Teissier,
  Lecture Notes in Mathematics \textbf{777},
  Springer-Verlag, 1980.
\end{description}

% END inputs/XXI/tex/components/37_expose_XXI_section75_p_morphisms_and_bibliography_en.tex


